% Santa Clara University PhD Thesis Template
% Document class: report, 12pt, letter paper, single-sided
%
% Structure: 10 chapters organized by estimator fit order (zeroth/first/
% second), per the fit-order-ladder outline agreed 2026-09-29.
%
% Chapter prose was rebuilt 2026-09-29 by copying from the three papers
% (IDETC 2025, the critical-points Systems paper, Draft_11). The scripts in
% sources/assemble/ regenerate each chapter file; every passage carries a
% "% Ported from" or "% New text" comment. Rerunning a script overwrites
% its chapter file.
\documentclass[12pt, letterpaper, oneside]{report}

% -----------------------------------------------------------------------
% Packages
% -----------------------------------------------------------------------
\usepackage[left=1.5in, right=1in, top=1in, bottom=1in, headheight=15pt]{geometry}
\usepackage{setspace}
\usepackage{amsmath, amsfonts, amssymb}
\usepackage{amsthm}
\usepackage{graphicx}
\graphicspath{{figures/}}
\usepackage{array}
\usepackage{caption}
\usepackage{algorithm}
\usepackage{algorithmic}
\usepackage{fancyhdr}
\usepackage{booktabs}
\usepackage{subcaption}
\usepackage{tikz}
\usepackage{url}
\usepackage{textcomp}
\usepackage{xcolor}
\usepackage[hidelinks]{hyperref}

% -----------------------------------------------------------------------
% Theorem environments -- one running counter, tied to chapter number
% -----------------------------------------------------------------------
\newtheorem{theorem}{Theorem}[chapter]
\newtheorem{lemma}[theorem]{Lemma}
\newtheorem{proposition}[theorem]{Proposition}
\newtheorem{corollary}[theorem]{Corollary}
\newtheorem{remark}[theorem]{Remark}
\newtheorem{assumption}[theorem]{Assumption}

% -----------------------------------------------------------------------
% Thesis metadata -- edit these fields
% -----------------------------------------------------------------------
\newcommand{\thesistitle}{Cooperative Multirobot Estimation and Navigation in Vector Fields: From Hardware Testbed to Second-Order Field Structure}
\newcommand{\thesisauthor}{Christopher J. Waight}
\newcommand{\thesisdepartment}{Mechanical Engineering}
\newcommand{\thesisdegreefield}{Mechanical Engineering}
\newcommand{\thesisadvisor}{Christopher A. Kitts}
\newcommand{\thesisreaderA}{Reader Name A}
\newcommand{\thesisreaderB}{Reader Name B}
\newcommand{\thesisreaderC}{Reader Name C}
\newcommand{\thesisreaderD}{Reader Name D}
\newcommand{\thesisdeptchair}{Department Chair Name}
\newcommand{\thesisyear}{2026}
\newcommand{\thesismonth}{Month}
\newcommand{\thesisday}{Day}

% -----------------------------------------------------------------------
% Page style setup
% -----------------------------------------------------------------------
\pagestyle{fancy}
\fancyhf{}
\fancyhead[R]{\thepage}
\renewcommand{\headrulewidth}{0pt}

\fancypagestyle{plain}{%
  \fancyhf{}
  \fancyhead[R]{\thepage}
  \renewcommand{\headrulewidth}{0pt}
}

% -----------------------------------------------------------------------
% TOC / LOF / LOT formatting
% -----------------------------------------------------------------------
\renewcommand{\contentsname}{Table of Contents}
\renewcommand{\listfigurename}{List of Figures}
\renewcommand{\listtablename}{List of Tables}

% -----------------------------------------------------------------------
% Hyphenation hints
% -----------------------------------------------------------------------
\hyphenation{multi-robot cluster-space Jacobian separ-atrix}
\setlength{\emergencystretch}{3em}  % let long math-heavy lines break

% -----------------------------------------------------------------------
% Pending numbers. \prelim{} prints a number from a reduced-size run in blue.
% It is copied from Draft_14 (Separatrix paper), where the 10^4-trial runs of
% overnight_runs.txt will replace these numbers. After that, rerun
% sources/assemble/all.py against the updated Draft_14 and the blue goes away.
% -----------------------------------------------------------------------
\newcommand{\prelim}[1]{\textcolor{blue}{#1}}

% -----------------------------------------------------------------------
% Placeholder figure macro
% -----------------------------------------------------------------------
\newcommand{\placeholderfig}[1]{\fbox{\rule{0pt}{2in}\rule{0.7\textwidth}{0pt}}\\[-1.6in]\centering\textit{[placeholder: #1]}\\[1.4in]}

% -----------------------------------------------------------------------
% Document
% -----------------------------------------------------------------------
\begin{document}

% Front matter. The approval page carries no number and is not counted; the
% title page is page i and carries no number. The titlepage environment is not
% used because it resets the page counter, which repeated the numerals.
\pagenumbering{gobble}

% ===================================================================
% 1. APPROVAL / SIGNATURE PAGE
% ===================================================================
\thispagestyle{empty}
\begin{center}
  \large
  Santa Clara University \\[6pt]
  Department of \thesisdepartment

  \vspace{0.3in}

  \normalsize
  Date: \thesismonth\ \thesisday, \thesisyear

  \vspace{0.3in}

  I HEREBY RECOMMEND THAT THE THESIS PREPARED UNDER MY SUPERVISION BY

  \vspace{0.2in}
  \MakeUppercase{\thesisauthor}
  \vspace{0.2in}

  ENTITLED

  \vspace{0.15in}
  \textit{\thesistitle}
  \vspace{0.15in}

  BE ACCEPTED IN PARTIAL FULFILLMENT OF THE REQUIREMENTS FOR THE DEGREE

  \vspace{0.15in}
  \textbf{DOCTOR OF PHILOSOPHY IN \MakeUppercase{\thesisdegreefield}}
\end{center}

\vspace{0.25in}

\noindent
\begin{tabular}{p{2.6in} p{0.2in} p{2.6in}}
  \hrulefill & & \hrulefill \\[2pt]
  Thesis Advisor & & Thesis Reader \\[20pt]
  \hrulefill & & \hrulefill \\[2pt]
  Chairman of Department & & Thesis Reader \\[20pt]
  & & \hrulefill \\[2pt]
  & & Thesis Reader \\[20pt]
  & & \hrulefill \\[2pt]
  & & Thesis Reader
\end{tabular}

\clearpage
\pagenumbering{roman}
\setcounter{page}{1}

% ===================================================================
% 2. TITLE PAGE
% ===================================================================
\thispagestyle{empty}
\begin{center}
  \vspace*{1in}

  {\Large \thesistitle}

  \vspace{0.6in}
  by \\[6pt]
  \thesisauthor

  \vspace{1.2in}
  Dissertation \\[6pt]
  Submitted in Partial Fulfillment of the Requirements \\
  for the Degree of Doctor of Philosophy \\
  in \thesisdegreefield \\[6pt]
  in the School of Engineering at \\
  Santa Clara University, \thesisyear

  \vspace{1in}
  Santa Clara, California \\
  \thesisyear
\end{center}
\clearpage

% ===================================================================
% 3. COPYRIGHT
% ===================================================================
\clearpage
\thispagestyle{plain}
\vspace*{3in}
\begin{center}
  \copyright\ Copyright by \thesisauthor\ \thesisyear \\
  All Rights Reserved
\end{center}
\clearpage

% ===================================================================
% 4. ACKNOWLEDGEMENTS
% ===================================================================
\chapter*{Acknowledgements}
\addcontentsline{toc}{chapter}{Acknowledgements}

% Placeholder, for the author to write.

I would like to thank my advisor, Christopher A. Kitts, and the members of
the Santa Clara University Robotic Systems Laboratory for their support
throughout this work.

Portions of this work appear in the papers listed in
Section~\ref{sec:ch1:publications}.

\clearpage

% ===================================================================
% 5. ABSTRACT
% ===================================================================
\chapter*{Abstract}
\addcontentsline{toc}{chapter}{Abstract}

\begin{center}
  \thesistitle \\[6pt]
  \thesisauthor \\[6pt]
  Department of \thesisdepartment \\
  Santa Clara University \\
  Santa Clara, California \\
  \thesisyear
\end{center}

\vspace{0.4in}

% Assembled from the three paper abstracts; linking sentences marked.
\noindent
Adaptive navigation is the process of modifying a robot's or a system of
robots' motion in real time based on measurements taken while moving. This
dissertation applies that process to vector fields and shows that the
order of the instantaneous local fit a cluster makes to its measurements
determines the field information available to it, and so the feature-based
navigation it can perform. An indoor testbed encodes direction in hue and
magnitude in saturation on printed floor maps, and a cross-robot neural
network calibration achieves an encoding-decoding $R^2$ of 0.96 in direction
and 0.91 in magnitude. After maintenance, a recalibration with pooled
training reads direction to within 6.1$^\circ$ and magnitude to within
3.8\% for the trained robots. On the testbed, direction-only primitives
reach a target the flow carries them to, but drift outward when asked to
orbit. A first-order fit from a minimum of three robots estimates critical
point locations without prior field knowledge, using only real-time local
measurements. In 459 hardware runs with three and four robots, the
cluster converged to within 0.05~m of the critical point in 99\% of saddle
and 91\% of vortex trials with three robots and in every vortex trial with
four, with mean position errors of 0.034~m or less, and orbits held the
commanded radius to within a quarter of that radius from 0.20~m upward
with three robots. A
second-order fit from six robots recovers the local quadratic model of a
flow from one synchronized set of measurements. In simulation, two
primitives built on it navigate to and along the trenches of two fields
computed from that fit, the determinant of the velocity gradient and the
smaller strain eigenvalue, which lie closest to a separatrix where strain
dominates rotation, with no integration horizon and no pre-straddling. The
determinant tracker withstands more measurement noise than the strain
tracker from every start tested, while in a rotating-frame trial only the
strain tracker, which is built from objective quantities, holds the same
path. On Santa Barbara Channel radar currents, both primitives acquire the
dominant transport corridor, but under small changes to the start only the
determinant tracker keeps its downstream direction. Across the three rungs,
the noise on a fitted coefficient grows as the inverse formation radius
raised to the order of the coefficient, and the steady offset of a cluster
from its target is the speed that its loop does not cancel divided by the
loop gain.

\clearpage

% ===================================================================
% 6. GLOSSARY OF TERMS
% ===================================================================
\chapter*{Glossary of Terms}
\addcontentsline{toc}{chapter}{Glossary of Terms}

\begin{table}[h]
\centering
\begin{tabular}{l p{4.2in}}
$\mathbf{p} = (x, y)$ & Position in the plane \\
$\mathbf{p}_c$, $\mathbf{p}_i$ & Cluster centroid, position of robot $i$ \\
$\mathbf{p}^*$ & Critical point, $\mathbf{v}(\mathbf{p}^*) = \mathbf{0}$ \\
$\mathbf{v} = (u, v)$ & Velocity field and its components \\
$\mathbf{J} = \nabla\mathbf{v}$ & Velocity gradient (field Jacobian) \\
$\mathbf{J}_c^{-1}$ & Inverse kinematic Jacobian of the cluster \\
$\mathbf{S}$, $\mathbf{W}$ & Rate-of-strain tensor and spin, the symmetric and antisymmetric parts of $\mathbf{J}$ \\
$\omega$ & Vorticity, $\partial v/\partial x - \partial u/\partial y$ \\
$s_1 \leq s_2$ & Eigenvalues of $\mathbf{S}$, with eigenvectors $\mathbf{e}_1$, $\mathbf{e}_2$ \\
$D$ & Determinant of the velocity gradient, $\det\mathbf{J}$ \\
$\mathbf{H}_D$ & Hessian of $D$ \\
$\alpha \pm i\nu$ & Complex eigenvalue pair of $\mathbf{J}$ (Chapter~\ref{ch:first_order}) \\
$\omega_o$ & Orbital angular rate, $k_t/r_d$ (Chapter~\ref{ch:first_order}) \\
$\mathbf{A}$ & Three-robot formation matrix of the first-order fit \\
$\boldsymbol{\Phi}$, $\boldsymbol{\Phi}_k$ & Formation matrix of the second-order fit, and of an order-$k$ fit \\
$n_k$ & Number of coefficients per component in an order-$k$ fit \\
$\kappa(\cdot)$ & Condition number \\
$p, \beta, q$ & SAS side lengths and interior angle of a cluster triangle \\
$\rho$ & Formation radius \\
$r$, $r_d$ & Orbital radius and commanded orbital radius \\
$\sigma_{uv}$, $\sigma_p$ & Standard deviations of measurement noise and position noise \\
$g$ & Share of a command that a robot attains, 0.86 (Section~\ref{sec:ch2:dynamics}) \\
$L$ & Transport delay of a robot command, 0.055~s \\
$\tau$ & Actuator time constant, 0.02~s \\
$v_c$ & Dead zone of a robot, 0.025~m/s \\
$\alpha_{\text{mom}}$ & Momentum coefficient of the sampled robot model, $e^{-\Delta t/\tau}$ \\
$\Delta t$ & Control period, 0.1~s (10~Hz) \\
$\Omega$ & Rotation rate of an observer frame \\
\end{tabular}
\end{table}

\clearpage

% ===================================================================
% 7. TABLE OF CONTENTS
% ===================================================================
\tableofcontents
\clearpage

% ===================================================================
% 8. LIST OF FIGURES
% ===================================================================
\listoffigures
\addcontentsline{toc}{chapter}{List of Figures}
\clearpage

% ===================================================================
% 9. LIST OF TABLES
% ===================================================================
\listoftables
\addcontentsline{toc}{chapter}{List of Tables}
\clearpage

% ===================================================================
% BODY -- switch to arabic page numbering
% ===================================================================
\pagenumbering{arabic}
\setcounter{page}{1}
\doublespacing

% ===================================================================
% CHAPTER 1: INTRODUCTION
% ===================================================================
\chapter{Introduction}
\label{ch:intro}

\section{Adaptive Navigation in Vector Field Environments}
\label{sec:ch1:an}

% Ported from: Draft_14.tex
Multirobot systems (MRS) offer redundancy, increased
throughput, and cooperative capabilities a single robot cannot produce, and have
succeeded across land, sea, air, and space \cite{1,2}. In adaptive
navigation (AN), a vehicle or group of vehicles modifies its trajectory in
real time from measurements taken while moving \cite{3}. A single vehicle
must move to sense a gradient, which costs time and misleads in a changing
field, while an MRS samples many points at once.

% Ported from: Draft_14.tex
Some flows are better described as vector fields, which carry features a
scalar map cannot. Critical points, such as sinks, sources, and vortices,
can model convergence zones \cite{4} and subsurface leaks \cite{5}.

% Ported from: cwaight_systems_paper_with_revisions2.tex
The oceanic and atmospheric environments that motivate this work are governed by a spatially varying velocity or force, and their transport and accumulation behavior is organized by the critical points of the field, the coordinates where the flow vanishes \cite{6,7}. A surface convergence zone acts as a sink, gathering buoyant debris and drifting bodies; this makes such zones high probability search regions after a loss at sea \cite{8,9}. Containing a ruptured wellhead or a harbor discharge requires reaching the source against its own outward flow \cite{5,10}. A mesoscale ocean eddy serves as a center, trapping a body of water, along with the plankton bloom inside it, and carrying both across a basin for weeks \cite{11,12,13}. Monitoring the bloom over its lifetime therefore requires a sensor to orbit the drifting core at a fixed standoff; a ground-fixed station would not stay with the eddy. These applications require the coordinate of the feature itself, whether to drive the multirobot team onto it or to hold it at a fixed standoff radius; a heading toward the feature is not sufficient.

\section{Literature Review}
\label{sec:ch1:lit}

\subsection{Adaptive Navigation in Scalar Fields}

% Ported from: Draft_14.tex
AN in scalar fields is well developed. \"{O}gren, Fiorelli, and Leonard
established cooperative gradient climbing with a formation of mobile
sensors using plane fitting \cite{14}. Bri\~{n}\'{o}n-Arranz, Renzaglia, and Schenato
extended the fit to recover the Hessian from five robots on a ring plus
one at the center \cite{15}. Kitts, McDonald, and
Neumann developed primitives for extrema finding, contour following,
ridge/trench following, and saddle point station keeping \cite{3}, and
verified ridge, trench, saddle, and front tracking on an indoor testbed
\cite{16,17}.

% Ported from: cwaight_systems_paper_with_revisions2.tex
Strategies include gradient following \cite{2}, Hessian estimation \cite{15}, extended information consensus filters \cite{18,19}, constrained nonsmooth distributed optimization \cite{20}, and bio-inspired search \cite{21}, and several such systems have been field-deployed on ground, surface, and aerial vehicles for environmental monitoring, pollution tracking, and search and rescue \cite{2,22,23,24}.

\subsection{Navigation in Vector Fields}

% Ported from: IDETC 2025 (DETC2025-167604), Sec. 1, transcribed from the published PDF
While multirobot adaptive navigation in scalar field environments is
well-studied both in simulation and practice, adaptive navigation through
vector fields remains largely theoretical, compared to extensive
experimental confirmation in scalar fields \cite{16}. Existing research
in vector fields primarily focuses on constructing artificial vector
fields for single-robot navigation, formation control, and obstacle
avoidance \cite{25}. In contrast, this research involves navigation
through naturally existing vector fields. The few studies examining
multirobot systems with distributed sensor measurements in physical
vector fields have identified promising control primitives for features
like flow maxima, sinks, sources, and vortices, but these approaches lack
comprehensive real-world validation: they were validated mainly in
simulation \cite{26}, and the PIM-triple approach of \cite{27} was
extended to robots in a flow tank \cite{28}.

% Ported from: cwaight_systems_paper_with_revisions2.tex
By contrast, only a few works address navigation of physically sensed vector fields \cite{29,30,27}. Much of the broader vector field literature constructs artificial fields for path-following and obstacle avoidance, a separate problem from sensing a real flow \cite{31,32,25,33}, and field campaigns that track features carried by an actual ocean current underscore the distinction \cite{34}.

\subsection{Robotic Sensing of Physical Flows}

% Ported from: cwaight_systems_paper_with_revisions2.tex
Work that senses the physical field falls into two camps. One camp presumes prior knowledge of the environment, which is unavailable in the unmapped settings that motivate this work \cite{35,36}. The other operates online in two behaviors: minimally actuated drifters that ride and station-keep along ambient streamlines, and PIM-triple-inspired teams that straddle a saddle and track a single stable or unstable manifold from local velocity measurements \cite{27,6,37}. Both supply a heading along the flow or track the manifold of a single assumed saddle; neither recovers the coordinate of a critical point together with its identity across the full taxonomy of feature types. Estimating the locations of critical points in vector fields remains an open problem for multirobot systems. To our knowledge, this dissertation gives the first hardware demonstration of a multirobot cluster that locates, reaches, and orbits the critical points of a sensed vector field from instantaneous measurements alone.

\subsection{Experimental Testbeds}

% Ported from: IDETC 2025 (DETC2025-167604), Sec. 1, transcribed from the published PDF
The gap between simulation and reality highlights the need to develop
experimental testbeds that allow validation of vector field navigation
strategies under real-world conditions. A testbed of this kind offers
researchers a cost-effective way to test vector field navigation designs
before deployment in challenging environments like ocean currents,
airflow systems, electromagnetic fields, and industrial magnetic settings
\cite{38,24}. Building upon the framework of Mokhtarian et al.
\cite{39}, such a platform validates theoretical models prior to
real-world implementation \cite{40,41}.

\section{Problem Statement}
\label{sec:ch1:problem}

% New text
This dissertation asks what a small cluster of robots can learn about a
planar vector field from one synchronized set of measurements, and how it
can steer on what it learns. The cluster has no map, no forecast, and no
prior knowledge of the features in the field. Each robot reads the local
velocity and knows its own position. The features of interest are the
ones that organize transport, the critical points where the flow vanishes
and the separatrices that divide it into regions.

\section{Thesis Statement}
\label{sec:ch1:thesis}

% New text
The order of the instantaneous local fit that a cluster makes to its
measurements determines the field information available to it, and so the
feature-based navigation it can perform. The statement concerns what one
set of simultaneous measurements supports; methods that add temporal
measurements or prior knowledge lie outside it. Fit order here means the
order to which the vector field itself is reconstructed. A heading, or the
slope of the sensed magnitude, read directly from the samples supports
direction-following behaviors, which reach a target that the flow carries
the cluster to and drift when asked to orbit. A first-order fit from three robots recovers the velocity gradient, and
with it the location and type of a critical point. A second-order fit
from six robots recovers the velocity gradient as a field over the
formation, which is enough to acquire and follow a separatrix. Each rung
needs more robots than the one below it, and its highest-order
coefficients are more sensitive to measurement noise.

\section{Contributions}
\label{sec:ch1:contributions}

% New text, assembled from the contribution lists of the three papers
This dissertation makes seven contributions. First, it presents an indoor
testbed for vector field adaptive navigation, with direction encoded in
hue and magnitude in saturation on printed floor maps, and a cross-robot
neural network calibration that reads both from the robots' color sensors
(Chapter~\ref{ch:tools}). Second, it tests the vector-sum and
vector-to-scalar primitives on that testbed and shows where
direction-only navigation reaches its target and where it drifts
(Chapter~\ref{ch:zeroth}). Third, it presents a cooperative estimation algorithm that
estimates the location and type of a critical point from instantaneous
measurements alone, and proves that three robots are the minimum for an
unrestricted affine fit from one set of measurements
(Chapter~\ref{ch:first_order}). Fourth, it validates attraction
and orbital control laws built on that estimate in 459 hardware runs
with three and four robots (Chapter~\ref{ch:first_order}). Fifth, it presents a
second-order estimator that recovers the local quadratic model of a flow
from six robots, the minimum, and fails only when the robots lie on a
common conic (Chapter~\ref{ch:second_order}). Sixth, it builds two
primitives on that fit, the $D$ tracker and the $s_1$ tracker, which
acquire and ride the trench of their field from instantaneous measurements
without pre-straddling, and it tests them in simulation on a double gyre,
on three further fields, and on measured ocean currents
(Chapter~\ref{ch:second_order}). The $D$ tracker tolerates more
measurement noise, while the $s_1$ tracker, being built from objective
quantities, holds its path under a rotating observer. Finally, it compares
the three rungs on one flow, on formation geometry, on noise, and on steady
offsets, and it states for each rung when the feature that it steers to is
the feature of the flow (Chapter~\ref{ch:crosscutting}).

\section{Publications}
\label{sec:ch1:publications}

% New text
Portions of this dissertation appear in the following papers.
\begin{itemize}
\item C. Waight and C. A. Kitts, ``A functional indoor testbed for
multirobot adaptive navigation in vector field environments,'' in
\textit{Proc. ASME Int. Design Engineering Technical Conf. and Computers
and Information in Engineering Conf.}, 2025, DETC2025-167604. Published.
Chapters~\ref{ch:tools} and~\ref{ch:zeroth}.
\item C. Waight and C. A. Kitts, ``Adaptive navigation of multirobot
systems to critical points in 2D vector fields in simulation and
experiment,'' submitted to \textit{IEEE Systems Journal}.
Chapter~\ref{ch:first_order}.
%% TODO(status): update when the critical-points paper is accepted.
\item C. Waight and C. A. Kitts, ``Second-order cooperative field
estimation for multirobot tracking of coherent flow structures,'' in
preparation for \textit{IEEE Systems Journal}.
Chapter~\ref{ch:second_order}.
\end{itemize}

\section{Availability of Code and Data}
\label{sec:ch1:availability}

% New text. Plan0 optional reproducibility statement. Each campaign folder
% has a README that lists its files, robots, and maps.
The code for every simulation, the scripts behind each figure, the
calibration data, and the logs of the three hardware campaigns are
available from the author on request. Each hardware campaign has a README
that lists its files, its robots, and its maps.

\section{Dissertation Organization}
\label{sec:ch1:organization}

% New text
Chapter~\ref{ch:tools} describes the Decabot testbed, the HSV encoding of
vector fields, the sensor calibration, and the robot dynamics.
Chapter~\ref{ch:preliminaries} defines the vector field quantities the
later chapters estimate, from critical points to objective Eulerian
coherent structures, and the benchmark fields.
Chapter~\ref{ch:estimation} describes the control architecture shared by
every primitive and the local polynomial fit, and it lays out the
fit-order ladder that organizes the rest of the dissertation.
Chapters~\ref{ch:zeroth}, \ref{ch:first_order}, and~\ref{ch:second_order}
climb that ladder one rung at a time, from direction-only primitives on
hardware, to critical point estimation and control on hardware, to
separatrix tracking in simulation. Chapter~\ref{ch:crosscutting} compares
the rungs directly, on the scalar fields each makes available, on
formation geometry, on noise, and on the number of robots.
Chapter~\ref{ch:limitations} states the limitations and future work, and
Chapter~\ref{ch:conclusion} concludes. The appendices give the field
definitions, the cluster kinematics for three and six robots, the frame
equivariance of the $s_1$ tracker, the stability arguments, and the
statistical methods.

% ===================================================================
% CHAPTER 2: VERIFICATION TOOLS
% ===================================================================
\chapter{Verification Tools}
\label{ch:tools}

% New text
The hardware results of this dissertation come from the Decabot testbed
of the Santa Clara University Robotic Systems Laboratory. The robots, the
motion capture system, and the cluster space software predate this work
\cite{42}. This dissertation adds the HSV encoding of vector fields
on printed floor maps, the neural network calibration of the robots'
color sensors, a remeasurement of the robots' dynamics, and the
simulation environments built around them.

\section{The Decabots}
\label{sec:ch2:decabots}

% Ported from: IDETC 2025 (DETC2025-167604), Sec. 2.2, transcribed from the published PDF
The robots used for this testbed are Decabots, a fleet of
omnidirectional rovers developed at Santa Clara University
\cite{42}. They are equipped with TCS34725 RGB sensors on their
undercarriage, enabling them to sense the color patterns of the colored
floor mat underneath and estimate the vector field being represented at
that point. Each is powered by an Intel Edison on an Arduino breakout board,
with a Teensy 3.2 that reads the encoders and the light sensor, and
communicates over TCP/IP to
a host computer that interprets the robot's sensed value, uses the
information to calculate and return robot velocity commands to each
robot. Each robot has a maximum linear velocity of 0.3~m/s. No
modifications to these robots from their original design were made for
this study.

A perspective view of the Decabot can be seen in
Fig.~\ref{fig:ch2:decabot}. It has omnidirectional wheels, as well as
small reflective spheres on its roof, which enable tracking with the
motion capture cameras.

\begin{figure}[htbp]
    \centering
    \includegraphics[width=0.5\textwidth]{ch2_decabot.jpg}
    \caption{A Decabot rover with OptiTrack motion tracking balls attached.}
    \label{fig:ch2:decabot}
\end{figure}

Fig.~\ref{fig:ch2:undercarriage} shows the undercarriage of the Decabot
rover. The TCS34725 has a white light LED onboard to illuminate the floor
map being sensed. The undercarriage also has a small shroud to minimize
the effects of external lighting. A simple test of turning the overhead
lab lights on and off confirmed that external lighting did not affect
readings from the test palette.

\begin{figure}[htbp]
    \centering
    \includegraphics[width=0.5\textwidth]{ch2_undercarriage.png}
    \caption{Undercarriage of a Decabot rover with the LED turned on and the
    external light shroud.}
    \label{fig:ch2:undercarriage}
\end{figure}

\section{Testbed and Motion Capture}
\label{sec:ch2:testbed}

% Ported from: IDETC 2025 (DETC2025-167604), Sec. 3.1, transcribed from the published PDF
This research used an existing testbed as a starting point
\cite{42}. Unlike general-purpose multi-agent platforms such as
Robotic Park \cite{43} that emphasize heterogeneous robot
integration, our testbed specifically focuses on vector field
representation and interpretation through a homogeneous robot fleet with
consistent sensing capabilities. It includes three or four omnidirectional
mobile rovers (Decabots), 1.6~m by 1.6~m multicolored floor maps representing
different vector fields, and a motion capture system (OptiTrack) to
verify the position of the robots relative to the floor mat with
centimeter level accuracy. The Decabots and OptiTrack communicate to a
central computer via WiFi. The central computer acts as the main
controller, hosting the control architecture implemented in
MATLAB/Simulink. The central computer processes the information received
then communicates robot velocity commands back to each individual robot.
Fig.~\ref{fig:ch2:testbed_overview} shows the layout of all the
equipment.

\begin{figure}[htbp]
    \centering
    \includegraphics[width=0.75\textwidth]{ch2_testbed_overview.png}
    \caption{Overview of the testbed developed for testing vector field
    adaptive navigation.}
    \label{fig:ch2:testbed_overview}
\end{figure}

% Ported from: cwaight_systems_paper_with_revisions2.tex
The testbed of Chapters~\ref{ch:zeroth} and~\ref{ch:first_order}, with three robots in the first and three or four in the second, is SCU's Robotic Systems Laboratory Decabot testbed \cite{42} (Fig.~\ref{fig:ch2:fullbed}), which has been used to validate multirobot adaptive navigation algorithms for cluster control and swarm particle optimization \cite{16,44}.

% Ported from: cwaight_systems_paper_with_revisions2.tex
\begin{figure}[htbp]
    \centering
    \includegraphics[width=0.5\textwidth]{ch2_testbed_with_four.png}
    \caption{Experimental testbed: 1.6 m × 1.6 m printed vector field map, Decabot rovers, and OptiTrack system. Four Decabots are shown; the experiments here used three or four of this fleet.}
    \label{fig:ch2:fullbed}
    \end{figure}

% Ported from: cwaight_systems_paper.tex
The multilayer control architecture of Chapter~\ref{ch:estimation} runs in MATLAB/Simulink on a central computer that communicates with the robots over Wi-Fi at 10~Hz: at each control cycle it receives OptiTrack positions and RGB sensor readings, computes the control laws, and returns velocity commands, which each robot's onboard computer (Intel Edison) converts into drive motor commands.

\section{Vector Field Encoding in HSV Color Space}
\label{sec:ch2:hsv}

\subsection{Limitations of RGB Encoding}

% Ported from: IDETC 2025 (DETC2025-167604), Sec. 1, transcribed from the published PDF
Previous attempts at creating vector field testbeds for multirobot
adaptive navigation had limited success due to technical barriers
\cite{45}. Approaches using red and green colors to represent
orthogonal vectors suffered from sensor channel interference on robotic
RGB sensors. This interference made encoding vectors to color maps
problematic, preventing accurate decoding of direction and magnitude
information. This issue was particularly severe when dealing with vectors
of small magnitude, as minor value changes could result in large
direction shifts. Additionally, the limited discrete sampling speed of
0.15~Hz made continuous motion of the robots impossible, limiting the
study of motion dynamics.

\subsection{Hue as Direction, Saturation as Magnitude}

% Ported from: IDETC 2025 (DETC2025-167604), Sec. 3.2, transcribed from the published PDF
Previous attempts to use RGB colors to represent orthogonal vector sets
showed limited success, as only a small portion of the spectrum remained
usable after eliminating high-interference regions. To allow for greater
color variance, an HSV color space was implemented to model vector space.
Hue, with its inherently circular nature, provided a natural
representation for vector direction. For magnitude representation,
saturation (ranging from 0 to 1) was chosen after early tests revealed
that the value channel was highly sensitive to printer ink variations,
ambient lighting conditions, and shadow effects.
Fig.~\ref{fig:ch2:hsv_space} shows a saturation-versus-hue plot
representing all possible vector representations.

\begin{figure}[htbp]
    \centering
    \includegraphics[width=0.55\textwidth]{ch2_hsv_space.png}
    \caption{A plot of saturation versus hue representing vector space.}
    \label{fig:ch2:hsv_space}
\end{figure}

% Ported from: cwaight_systems_paper_with_revisions2.tex
Vector fields are physically realized as printed floor maps using HSV color encoding, where hue (0--2$\pi$) represents flow direction and saturation (0--1) encodes magnitude. Printed magnitudes are scaled to the 0.3--1.0 saturation range to avoid undersaturated areas and ensure reliable sensor readings, then remapped to 0--1 in software, with saturation 1.0 equal to a field speed of 0.3~m/s.

\subsection{Printed Field Construction}

% New text
Each printed map is generated from an analytical field. The fixed and
sinking vortex maps of Chapter~\ref{ch:zeroth} and the vortex and saddle
maps of Chapter~\ref{ch:first_order} were built this way, and their
equations are given in Appendix~\ref{app:fields}.

% Ported from: IDETC 2025 (DETC2025-167604), Sec. 3.2, transcribed from the published PDF
Once vectors were created at all plot points, they were scaled to the
range of 0.3 to 1. Values less than 0.2 were not used in training, and
values less than 0.3 produced velocities insufficient to overcome static
friction.

\section{Neural Network Sensor Calibration}
\label{sec:ch2:calibration}

% Ported from: IDETC 2025 (DETC2025-167604), Sec. 3.3, transcribed from the published PDF
The RGB sensor data suffered from channel interference as noted in
previous studies \cite{45}. The TCS34725 RGB sensor specification
sheet confirms that blue, green, and red channels cannot be fully
isolated. This effect persisted even after conversion to an HSV color
space. Traditional calibration methods neither significantly improved
performance nor supported effective scaling across multiple robots,
prompting exploration of a supervised learning approach.

\subsection{Calibration Palette and Data Collection}

A calibration palette with varying hues ($0$ to $2\pi$) and saturations
was created, with values spaced at 24 equal intervals in each dimension.
Three robots were positioned above each subplot, recording RGB and K
measurements alongside the intended hue and saturation levels (see
Fig.~\ref{fig:ch2:training_data}).

\begin{figure}[htbp]
    \centering
    \includegraphics[width=0.6\textwidth]{ch2_training_data.png}
    \caption{A Decabot collecting training data for the vector estimation
    models.}
    \label{fig:ch2:training_data}
\end{figure}

The dataset from three robots was combined and augmented by introducing
$\pm1\%$ noise to RGB readings to ensure model robustness. The collected
RGBK values were normalized between 0 and 1, converted to HSV using
standard MATLAB libraries, and appended to the original data, resulting
in a 7-dimensional input space (R,G,B,K,H,S,V). Data points with
saturation values below 0.2 were removed to boost model performance. All
inputs and targets were normalized to the $[-1,1]$ range before training.

\subsection{Network Architecture and Training}

Two separate neural networks were created using MATLAB's Deep Learning
Toolbox. The architecture for both networks consisted of 7 input neurons
(R,G,B,K,H,S,V), followed by two fully-connected hidden layers with tanh
activation functions. The saturation network had a single output neuron
for direct magnitude prediction, while the hue network had two output
neurons representing sine and cosine components to handle the circular
nature of directional data.

A random search approach determined the optimal hidden layer sizes,
exploring 2 to 12 neurons in the first layer and 4 to 10 neurons in the
second. Model patience and augmentation parameters were also randomized,
while the Levenberg-Marquardt optimizer and default learning rate
remained constant. Data was split into training (70\%), validation
(15\%), and test (15\%) sets.

The models were evaluated against 8 different validation sets including
calibration and verification plots from all four robots. The selection
criterion focused on worst-case performance. For each candidate model,
the lowest $R^2$ score across all validation sets was identified, and the
model with the highest minimum score was selected for each network. This
min-max selection strategy ensured consistent performance across all test
conditions.

\subsection{Cross-Robot Generalization}

% Ported from: IDETC 2025 (DETC2025-167604), Sec. 4.1, transcribed from the published PDF
The supervised learning approach for sensor calibration demonstrated
excellent vector field interpretation capabilities across all validation
datasets. Using a cross-robot training strategy, we achieved high
accuracy in both direction and magnitude sensing, with robust
generalization to new robots.

As shown in Table~\ref{tab:ch2:calibration}, the cross-trained model
maintained circular $R^2$ values consistently above 0.96 for direction
(hue) and above 0.90 for magnitude (saturation) across all robots and
datasets. The model's generalization capability is most evident in its
performance on Robot 4 (R4), which wasn't included in the training
dataset. Despite having no individual calibration data for R4, the
cross-trained model achieved a circular $R^2$ of 0.993 (RMSE 0.024) for
direction and 0.908 (RMSE 0.073) for magnitude on the calibration plot.

\begin{table}[htbp]
\centering
\caption{Comparative performance of cross-trained and individually tuned
models. R4 was not in the training set and has no individually tuned
model.}
\label{tab:ch2:calibration}
\begin{tabular}{llcccc}
\hline
 & & \multicolumn{2}{c}{Hue} & \multicolumn{2}{c}{Saturation} \\
Robot-dataset & Model & $R^2$ & RMSE & $R^2$ & RMSE \\
\hline
R1-Cal & Cross & 0.966 & 0.017 & 0.950 & 0.054 \\
       & Indiv & 0.999 & 0.009 & 0.972 & 0.040 \\
R1-Ver & Cross & 0.969 & 0.050 & 0.950 & 0.064 \\
       & Indiv & 0.995 & 0.020 & 0.963 & 0.055 \\
R2-Cal & Cross & 0.998 & 0.013 & 0.935 & 0.061 \\
       & Indiv & 0.996 & 0.018 & 0.968 & 0.043 \\
R2-Ver & Cross & 0.967 & 0.052 & 0.943 & 0.069 \\
       & Indiv & 0.991 & 0.027 & 0.952 & 0.063 \\
R3-Cal & Cross & 0.998 & 0.014 & 0.948 & 0.055 \\
       & Indiv & 0.999 & 0.007 & 0.980 & 0.034 \\
R3-Ver & Cross & 0.991 & 0.027 & 0.953 & 0.062 \\
       & Indiv & 0.997 & 0.015 & 0.976 & 0.044 \\
R4-Cal & Cross & 0.993 & 0.024 & 0.908 & 0.073 \\
R4-Ver & Cross & 0.964 & 0.054 & 0.925 & 0.079 \\
\hline
\end{tabular}
\end{table}

While individually-tuned models achieved marginally better performance
(typically $R^2 > 0.99$ for hue, $R^2 > 0.95$ for saturation), the
cross-trained approach offered substantial time savings by eliminating
several hours of calibration per robot. This cross-trained approach
significantly improves fleet scalability for vector field navigation. New
robots can simply be loaded with the pre-trained model and verified using
a smaller validation plot, eliminating the need for extensive individual
calibration procedures while maintaining high accuracy.

\subsection{Recalibration After Maintenance}

% New text
The robots were recalibrated after preventive maintenance, before the
critical point experiments of Chapter~\ref{ch:first_order}. The
recalibrated models are the ones the later hardware results rely on.
Their errors are reported as RMS errors in degrees of hue and percent of
saturation, where Table~\ref{tab:ch2:calibration} gives the circular $R^2$
and the RMSE of hue in turns and of saturation on a unit scale. An RMS error
of 6.1$^\circ$ in hue is an RMSE of 0.017 on the scale of that table, and
3.8\% in saturation is 0.038.

% Ported from: cwaight_systems_paper_with_revisions2.tex
The recalibration repeats the palette and the two networks of the previous sections with two changes. The input is the four normalized raw channels $(R,G,B,K)$ in place of seven inputs, and training pools data across robots so that a single model can be applied to an uncalibrated unit without per-robot tuning. On the printed maps, this calibration gave the three trained robots RMS errors of 6.1$^\circ$ in hue and 3.8\% in saturation. The fourth robot, never used in training, read hue as accurately (6.3$^\circ$) but reported field magnitude 12 to 14\% low.

% Ported from: cwaight_systems_paper.tex
\begin{figure}[htbp]
    \centering
    \includegraphics[width=0.8\textwidth]{ch2_color_sensor_hsv_predictions.png}
    \caption{Neural network predictions for HSV field encoding on unit-normalized axes.}
    \label{fig:ch2:HuePredictions}
\end{figure}

\section{Saturation-to-Velocity Characterization}
\label{sec:ch2:saturation}

% Ported from: IDETC 2025 (DETC2025-167604), Sec. 3.4, transcribed from the published PDF
To validate the relationship between detected vector magnitude and robot
velocity, lanes of a single hue at different saturation levels were
printed (Fig.~\ref{fig:ch2:lanes}). Each robot was placed on each of the
different saturation lanes and commanded to follow the direction based on
hue, and to use the magnitude indicated by saturation to scale its
velocity.

\begin{figure}[htbp]
    \centering
    \includegraphics[width=0.6\textwidth]{ch2_saturation_lanes.png}
    \caption{Lanes of constant hue and varying saturation for testing
    Decabot velocities.}
    \label{fig:ch2:lanes}
\end{figure}

Each robot's position was tracked over a 3 second period as it followed
each lane. This continuous motion capability stands in contrast to sparse
sampling approaches necessary in marine environments \cite{46}, where
energy constraints and environmental factors limit sampling frequency.
The average velocity for each saturation run was computed. Across all
robots, the velocities scaled linearly with saturation at values above
0.3. As noted earlier, saturation values below 0.3 did not produce
velocity commands strong enough to overcome static friction of the
robots. The robots achieved a max velocity of 0.3~m/s. This consistency
between saturation and robot velocity across robots is key for
interpreting vector primitive performance.

\section{Robot Dynamics Identification}
\label{sec:ch2:dynamics}

% Ported from: cwaight_systems_paper_with_revisions2.tex
In simulation, each robot is modeled as a single omnidirectional point mass that follows its command after a transport delay $L$, through a first order response with actuator time constant $\tau$, a dead zone $v_c$, and a speed limit $v_{\max}$. The robot attains a fraction $g$ of each command, so its velocity $\mathbf{v}_{r,i} = \dot{\mathbf{p}}_i$ follows
\begin{equation}
    \tau\,\dot{\mathbf{v}}_{r,i}(t) = -\mathbf{v}_{r,i}(t) + g\,\mathbf{u}_i(t - L),
    \label{eq:ch2:continuous}
\end{equation}
where $\mathbf{u}_i$ is zero for commands below $v_c$. Through system characterization of the hardware, it was found that $g = 0.86$, $L + \tau = 0.075$~s (split as $\tau = 0.02$~s and $L = 0.055$~s), $v_{\max} = 0.3$~m/s, and $v_c = 0.025$~m/s.

For digital implementation, each command is held over the control period $\Delta t$ of cycle $j$, with a fractional delay $f = L/\Delta t \in [0, 1)$, yielding
\begin{equation}\label{eq:ch2:momentum}
\begin{split}
&\bar{\mathbf{u}}_i[j] = (1-f)\,\mathbf{u}_i[j] + f\,\mathbf{u}_i[j-1],\\
&\tilde{\mathbf{v}} = \alpha_{\text{mom}}\,\mathbf{v}_{r,i}[j] + (1-\alpha_{\text{mom}})\,g\,\bar{\mathbf{u}}_i[j],\\
&\mathbf{v}_{r,i}[j+1] = \tilde{\mathbf{v}}\,\min\{1,\ v_{\max}/\|\tilde{\mathbf{v}}\|\},\\
&\mathbf{p}_i[j+1] = \mathbf{p}_i[j] + \Delta t\,\mathbf{v}_{r,i}[j+1],
\end{split}
\end{equation}
where $\alpha_{\text{mom}} = e^{-\Delta t/\tau}$ and the dead zone acts on $\bar{\mathbf{u}}_i$. At a 10~Hz control rate, $f = 0.55$ and $\alpha_{\text{mom}} \approx 0.007$, so the robot reaches $g$ times its command within one cycle.

% New text. The fitted dead zone and the 0.05 m/s of the 2025 tests.
The dead zone of 0.025~m/s is a fitted value. The 2025 tests place the
floor near 0.05~m/s: lanes below a saturation of 0.3, a command of about
0.09~m/s, did not overcome static friction (Section~\ref{sec:ch2:saturation}).
The two numbers measure different things. A lane test combines the decoding
of a low saturation with the friction of the robot, while the fit uses the
commanded and the measured velocity directly. The runs of both campaigns
come to rest where the attraction command falls to about 0.03 to
0.05~m/s, and no run commanded less than about 0.06~m/s, so neither number
is a direct measurement of the threshold.

\section{Field Reconstruction as a Noise Model}
\label{sec:ch2:reconstruction}

% Ported from: cwaight_systems_paper_with_revisions2.tex
This method of representing vector fields in printed floor maps has imperfections in both printing and sensing. Fig.~\ref{fig:ch2:measurement_error} shows the analytical vortex and saddle fields, reconstructions of each field from sensed robot measurements, and the reconstruction error (right column). These reconstructed fields were used to fit multilayer perceptron (MLP) interpolation models, used in Chapter~\ref{ch:first_order} as the realistic noise model for simulation.

% Ported from: cwaight_systems_paper_with_revisions2.tex
\begin{figure}[htbp]
    \centering
    \includegraphics[width=0.75\textwidth]{ch2_measurement_error_comparison.png}
    \caption{Analytical fields, sensor-based reconstructions, and reconstruction error for vortex (top) and saddle (bottom) fields.}
    \label{fig:ch2:measurement_error}
\end{figure}

\section{Simulation Environment}
\label{sec:ch2:simulation}

% New text
Both simulation studies use a Python implementation of the three-layer
architecture of Section~\ref{sec:ch4:architecture}, with the robot model
of Section~\ref{sec:ch2:dynamics}. The three-robot simulator of
Chapter~\ref{ch:first_order} reads either the analytical fields of
Appendix~\ref{app:fields} or the reconstructed fields above, which carry
the testbed's own sensing error. The six-robot simulator of
Chapter~\ref{ch:second_order} reads analytical or gridded fields and adds
optional measurement noise on each reading and position noise at
measurement time, leaving the true robot state untouched.

\section{Testbed Limitations}
\label{sec:ch2:limitations}

% Ported from: IDETC 2025 (DETC2025-167604), Sec. 5, transcribed from the published PDF
Despite significant improvements over previous testbeds, important
limitations persist. Indoor small-scale testbeds often inadequately
represent large-scale environments \cite{41}, and our system is
currently restricted to static vector fields, preventing reactive control
testing in dynamic environments \cite{38} with no clear path to enable
this using colormaps. Constrained to two-dimensional vector fields, 3D
extensions would require additional sensing capabilities.

The robots sample only a single point in the field at their current
location without look-ahead capability, hindering testing of predictive
path planning methods. The testbed is limited to homogeneous robots and
lacks obstacle avoidance sensors, restricting studies of heterogeneous
multi-robot behaviors in complex environments. Physical constraints
include a maximum 3~m by 6~m work area, approximately one hour battery
life per robot, and a limited velocity range of 0.05 to 0.3~m/s in the 2025
tests (Section~\ref{sec:ch2:dynamics}).
Additionally, the Decabots' wheels degrade both paper and printed colors
on colormaps, necessitating replacement after a few dozen runs and adding
significant operational costs.

\section{Verification by Substitution}
\label{sec:ch2:substitution}

% New text. The studies of this dissertation earn their claims one way: swap
% one part of the system, hold the rest fixed, and see what changes.
The claims of this dissertation rest on one habit: swap one part of the
system, hold the rest fixed, and see what changes.
Table~\ref{tab:ch2:substitution} lists the substitutions, and each isolates
a single question.

\begin{table}[htbp]
\centering
\caption{Substitutions used to isolate one question at a time.}
\label{tab:ch2:substitution}
\small
\begin{tabular}{>{\raggedright\arraybackslash}p{5.2cm} >{\raggedright\arraybackslash}p{7.0cm} c}
\hline
Substitution & What it isolates & Chapter \\
\hline
Robot R4 held out of the calibration training & whether the sensor calibration generalizes to a new robot & \ref{ch:tools} \\
Analytical field against reconstructed field & the fixed distortion of the printed map against measurement noise that changes in time & \ref{ch:first_order} \\
A simulation matched to each hardware run & what the model reproduces and what it omits & \ref{ch:first_order} \\
The vector-to-scalar primitive alone, with an oracle target, and with a speed estimate & the missing critical point coordinate against the steering direction & \ref{ch:first_order} \\
Measurement noise and position noise swept separately & which noise limits each tracker & \ref{ch:second_order} \\
Scoring against the separatrix and against the tracker's own trench, with analytic trenches & the field against the tracker & \ref{ch:second_order} \\
A rotating observer & objectivity & \ref{ch:second_order} \\
A replay of the ocean start with small perturbations & whether the headline run was selected & \ref{ch:second_order} \\
\hline
\end{tabular}
\end{table}

The studies also form a sequence. The testbed paper found an outward drift
that simulation had not predicted, and it promised digital twins of the
Decabots that capture their dynamics for simulation, together with
controller substitution on a fixed lower architecture. The first-order
study delivered both. It fitted the robot model of
Section~\ref{sec:ch2:dynamics}, used it to explain the drift as a limit
cycle, and substituted the vector-to-scalar primitive for the Jacobian
method with the lower layers held fixed. The second-order study runs its
six-robot simulations on the same model.

% ===================================================================
% CHAPTER 3: VECTOR FIELD PRELIMINARIES
% ===================================================================
\chapter{Vector Field Preliminaries}
\label{ch:preliminaries}

% New text
The features a cluster can navigate to are properties of the velocity
field and its derivatives. This chapter defines them on the true field.
Chapters~\ref{ch:first_order} and~\ref{ch:second_order} read the same
quantities from local fits.

\section{Planar Vector Fields}
\label{sec:ch3:fields}

% Ported from: cwaight_systems_paper_with_revisions2.tex
In its most general form, a vector field assigns a vector $\mathbf{v}(\mathbf{p}) \in \mathbb{R}^m$ to each point $\mathbf{p} \in \mathbb{R}^n$ within a given domain, where $m \geq 2$, making vector fields well suited to represent directional physical quantities such as fluid velocity, electromagnetic force, or gravitational force. For the planar fields analyzed in this work, this takes the form of a 2D vector $\mathbf{v}(\mathbf{p}) = (u(\mathbf{p}), v(\mathbf{p}))$ assigned to each coordinate $\mathbf{p} = (x, y)$.

\section{Critical Points and Their Classification}
\label{sec:ch3:critical_points}

% Ported from: cwaight_systems_paper_with_revisions2.tex
Unlike scalar fields where $m=1$, vector fields can exhibit features such as critical points $\mathbf{p}^*$ where the field itself vanishes, $\mathbf{v}(\mathbf{p}^*) = \mathbf{0}$. Critical points model phenomena such as vortices arising from tidal flows over submarine caves, or plumes when searching for pollution sources \cite{47}.

% Ported from: cwaight_systems_paper_with_revisions2.tex
Near a critical point, the field is well approximated by its linearization, and, except at a center, the local topology is determined by the Jacobian matrix $\mathbf{J} = \nabla\mathbf{v}|_{\mathbf{p}}$. The eigenvalues of $\mathbf{J}$ evaluated at $\mathbf{p}^*$ determine the topological type. For $2 \times 2$ systems with nonzero determinant, six standard critical point types arise from the eigenvalue structure, as shown in Fig.~\ref{fig:ch3:allfield} and classified in Table~\ref{tab:ch3:field_summary}.

% Ported from: cwaight_systems_paper_with_revisions2.tex
Eigenvalue analysis can also identify degenerate scenarios such as cases where the determinant is zero. In uniform flow fields, for instance, no critical point exists. While these scenarios are detectable through the estimation framework of Chapter~\ref{ch:first_order}, they fall outside the scope of the present work. A nonlinear field can also hold several critical points, and those of opposite determinant sign are separated by curves on which the determinant is zero.

% Ported from: cwaight_systems_paper_with_revisions2.tex
\begin{figure}[htbp]
    \centering
    \includegraphics[width=0.6\textwidth]{ch3_six_vector_fields.png}
    \caption{Six linear vector field types with critical points.}
    \label{fig:ch3:allfield}
\end{figure}

% Ported from: cwaight_systems_paper_with_revisions2.tex
\begin{table}[htbp]
    \centering
    \caption{Classification of critical points by Jacobian eigenvalue structure ($\lambda$ = eigenvalue; a complex pair $\alpha \pm i\nu$ has real part $\alpha$, imaginary part $\nu$).}
    \label{tab:ch3:field_summary}
    \footnotesize
    \setlength{\tabcolsep}{2pt}
    \begin{tabular*}{\columnwidth}{@{\extracolsep{\fill}}lll}
    \hline
    \textbf{Field Type} & \textbf{Critical Point} & \textbf{Eigenvalues} \\
    \hline
    Vortex & Center & $\pm i\nu$ \\
    Sinking Vortex & Stable Spiral & $\alpha \pm i\nu$, $\alpha < 0$ \\
    Sink & Stable Node & $\lambda_1, \lambda_2 < 0$ \\
    Source & Unstable Node & $\lambda_1, \lambda_2 > 0$ \\
    Spewing Vortex & Unstable Spiral & $\alpha \pm i\nu$, $\alpha > 0$ \\
    Saddle & Saddle Point & $\lambda_1 < 0 < \lambda_2$ \\
    \hline
    \end{tabular*}
\end{table}

\section{Rotation, Strain, and the Okubo-Weiss Partition}
\label{sec:ch3:ow}

% New text
A critical point is one place where the field vanishes. Most of a flow
has no critical point, and there the velocity gradient still says how
nearby fluid moves relative to the point. The gradient
$\mathbf{J} = \nabla\mathbf{v}$ splits into a symmetric part
$\mathbf{S} = \tfrac{1}{2}(\mathbf{J} + \mathbf{J}^{\top})$, the
rate-of-strain tensor, and an antisymmetric part
$\mathbf{W} = \mathbf{J} - \mathbf{S}$, the spin, which carries the
vorticity $\omega = \partial v/\partial x - \partial u/\partial y$. The
eigenvalues $s_1 \leq s_2$ of $\mathbf{S}$ are the principal strain
rates, and their eigenvectors $\mathbf{e}_1$ and $\mathbf{e}_2$ are the
compression and stretching directions. The determinant
$D = \det\mathbf{J}$ combines both parts.

% Ported from: Draft_14.tex
For incompressible flow the determinant is a negative multiple of the
Okubo-Weiss parameter \cite{4,48}, so the two carry the same partition of
the flow. Where $D < 0$ the flow is strain-dominated, and where $D > 0$
it is rotation-dominated.

\section{Objective Eulerian Coherent Structures}
\label{sec:ch3:oecs}

% New text
Coherent structures are the curves that organize transport in a flow.
Lagrangian coherent structures are found by integrating particle
trajectories over a time window. This dissertation uses the
instantaneous, Eulerian kind, read from the strain tensor at one time.

% Ported from: Draft_14.tex
Serra and Haller define objective Eulerian coherent structures (OECS)
from these eigenvalues and eigenvectors \cite{49}. An attracting OECS is a
trench of $s_1$ tangent to $\mathbf{e}_2$, and floating material has been
seen at sea to collect on such curves \cite{50}. A repelling OECS is a
ridge of $s_2$ tangent to $\mathbf{e}_1$, where neighboring material
separates. In incompressible flow $s_2 = -s_1$, so repelling structures are
trenches of $s_1$ too.

\section{Relation Between the Fields}
\label{sec:ch3:relation}

% Ported from: Draft_14.tex
Both fields come from the same velocity gradient. With the vorticity
$\omega = \partial v/\partial x - \partial u/\partial y$, the Jacobian
splits into its mean, strain, and rotation parts, and the two fields become
\begin{equation}
    D = \mu^2 - r^2 + \tfrac{1}{4}\omega^2, \qquad s_1 = \mu - r.
    \label{eq:ch3:D_parts}
\end{equation}
In divergent flow, $D < 0$ therefore means that strain outweighs rotation
and divergence together. A trench of $D$ is a curve where strain most
outweighs rotation, and a trench of $s_1$ is a curve of strongest
instantaneous compression or stretching. Neither mentions streamlines, so
the trenches mark a separatrix only where the rotation on it is small
against the strain. The ratio $|\omega|/2r$ measures how small. A small
ratio is necessary rather than sufficient.

Where the vorticity vanishes, $\mathbf{J} = \mathbf{S}$ and
$D = s_1 s_2$, which is $-s_1^2$ in incompressible flow. The vorticity
also separates the two fields across observers. A frame rotating clockwise
at rate
$\Omega$, with rotation $\mathbf{Q}$, adds a uniform spin that raises
$\omega$ to $\omega + 2\Omega$. The strain tensor only co-rotates, so
$s_1$ is objective \cite{49}, while the determinant keeps the spin,
\begin{equation}
    D' = D + \Omega\,\omega + \Omega^2 .
    \label{eq:ch3:D_not_objective}
\end{equation}
Relatively rotating observers therefore disagree about where the
trenches of $D$ lie, and about $|\omega|/2r$, since the frame adds $2\Omega$
to $\omega$ but leaves $r$ unchanged.

% New text. Instantaneous Eulerian fields are surrogates for material
% transport structures (Draft 14 Sec. VI; Plan0 Sec. 3E).
Both fields are instantaneous and Eulerian, so they are surrogates for the
Lagrangian structures that carry material, and not the structures
themselves \cite{51}. A trench of $D$ or $s_1$ marks a separatrix of
the flow only where the rotation on it is small against the strain, and
even there it need not match the Lagrangian structure. Chapter~\ref{ch:second_order}
states each result for the trench, and it names a separatrix only where the
benchmark supports it.

\section{Benchmark Fields}
\label{sec:ch3:benchmarks}

\subsection{Canonical Linear Fields}
\label{sec:ch3:canonical}

% New text
Each critical point type of Table~\ref{tab:ch3:field_summary} has a
linear representative, shown in Fig.~\ref{fig:ch3:allfield}. Their
equations are given in Appendix~\ref{app:fields}. A linear field has a
constant Jacobian, so a first-order fit recovers it exactly on any
nondegenerate formation.

\subsection{The Steady Double Gyre}
\label{sec:ch3:dg_example}

% New text. The nonlinear gyre cell of Chapter 6 is one cell of this flow.
The gyre cell that tests the first-order estimate in
Chapter~\ref{ch:first_order} is the left cell of this flow, shifted so that
its center sits at the origin.

% Ported from: Draft_14.tex
The double gyre is
a benchmark flow on which both fields mark the material transport structure
exactly, and every
quantity of interest has a closed form. The flow is two counter-rotating gyres in a
closed basin, separated by a separatrix that runs between two saddle points on
the basin walls. We use the steady member of the Shadden
family~\cite{52}, shifted from the canonical domain
$[0,2]\times[0,1]$ to $[-1,1]\times[-0.5,0.5]$ by $x_f = x + 1$,
$y_f = y + 0.5$ to center it on the origin. All double-gyre quantities are
non-dimensional. The field is
\begin{equation}
u = -\pi A \sin(\pi x_f)\cos(\pi y_f), \quad
v = \pi A \cos(\pi x_f)\sin(\pi y_f),
\label{eq:ch3:dg_u}
\end{equation}
with $A = 0.1$. The separatrix is $x = 0$, the saddles are
$\mathbf{p}_1^* = (0, -0.5)$ and $\mathbf{p}_2^* = (0, 0.5)$, and the gyre
centers $(\pm 0.5, 0)$.

The analytical Jacobian entries are
$u_x = -\pi^2 A \cos\pi x_f \cos\pi y_f = -v_y$ and
$u_y = \pi^2 A \sin\pi x_f \sin\pi y_f = -v_x$, with subscripts denoting
partial derivatives of the true field. The flow is therefore
incompressible, and its determinant and Hessian are
\begin{equation}
\begin{aligned}
D(\mathbf{p}) &= -\tfrac{1}{2}\pi^4 A^2 \left(\cos 2\pi x_f + \cos 2\pi y_f\right),\\
\mathbf{H}_D &= 2\pi^6 A^2
\begin{bmatrix} \cos 2\pi x_f & 0 \\ 0 & \cos 2\pi y_f \end{bmatrix}.
\end{aligned}
\label{eq:ch3:det_analytic_main}
\end{equation}
The landscape of $D$ forms a mountain pass over the separatrix
(Fig.~\ref{fig:ch3:dg_fields}(b)). Its peaks are the two gyre centers, where
rotation dominates and $D > 0$. Its lowest points are the flow saddles
$\mathbf{p}_1^*$ and $\mathbf{p}_2^*$, joined by a trench along the
separatrix that rises to a crest at the origin, where $D = 0$ and the
separatrix splits into an upper and a lower segment. Across the trench
$D$ has constant positive curvature. Further trenches run along the four
walls and meet at the flow saddles and the four domain corners.

The $s_1$ landscape has the same trenches (Fig.~\ref{fig:ch3:dg_fields}(c)).
The shear strain vanishes identically, since $v_x = -u_y$, so
$\mathbf{S} = \mathrm{diag}(u_x, -u_x)$ everywhere and
\begin{equation}
s_1 = -\pi^2 A \left|\cos\pi x_f \cos\pi y_f\right|,
\label{eq:ch3:s1_analytic_main}
\end{equation}
with minima $s_1 = -\pi^2 A$ at the six flow saddles, where the trenches
cross. The two fields share these trenches because the vorticity
$\omega = -2\pi^2 A \sin\pi x_f \sin\pi y_f$ vanishes on them, so
$\mathbf{J} = \mathbf{S}$ and $D = -s_1^2$ there. Between the trenches
$s_1$ rises to zero along the lines $x = \pm 0.5$ and $y = 0$, where
$\mathbf{S}$ vanishes, so the crest of the separatrix's $s_1$ trench is a
kink at the origin.

% Ported from: Draft_14.tex
\begin{figure}[htbp]
    \centering
    \includegraphics[width=0.55\textwidth]{ch3_double_gyre_fields.png}
    \caption{Steady double gyre. (a) Streamlines, with the separatrix
    $x = 0$ (dashed), the saddle points (crosses), and the gyre centers
    (dots). (b) $D$ as an elevation surface: the gyre centers are peaks,
    and the separatrix (bold) is a trench, deepest at the saddles and
    rising to $D = 0$ at the origin. (c) $s_1$, with the same trench and a
    kink at its crest at the origin (dot).}
    \label{fig:ch3:dg_fields}
\end{figure}

% ===================================================================
% CHAPTER 4: COOPERATIVE ESTIMATION FRAMEWORK
% ===================================================================
\chapter{Cooperative Estimation Framework}
\label{ch:estimation}

% New text
Every primitive in this dissertation runs on the same control
architecture and follows the same estimation pattern. The robots sample
the field at one instant, a local polynomial is fit to those samples, and
a control law steers the cluster on what the fit reveals. This chapter
describes the architecture and the estimation pattern once.

\section{Multilayer Control Architecture}
\label{sec:ch4:architecture}

% Ported from: Draft_14.tex
The cluster runs the three-layer cluster space architecture of
\cite{2} (Fig.~\ref{fig:ch4:control_architecture}). The adaptive
navigation layer turns the robots' field samples into a cluster command,
the cluster controller layer holds the formation, and the robot controller
layer drives each robot.

% Ported from: IDETC 2025 (DETC2025-167604), Sec. 2.1, transcribed from the published PDF
This architecture was selected because its modular layered approach
allows for independent development and testing of different components
of the overall architecture. This approach also maximizes system
reusability, enabling systematic comparison of navigation strategies
using standardized components in downstream layers. Additionally, the
architecture enables the same control laws to be applied across different
robot platforms and formations.

% Ported from: cwaight_systems_paper.tex
At each control cycle (10~Hz), each robot reports its position and vector field measurement. The cluster space kinematics are computed, and the cluster shape velocity commands needed for formation control are determined. The adaptive navigation layer computes the critical point estimate and generates a cluster velocity command. The cluster space controller maps all velocity commands to individual robot velocities via the inverse kinematic Jacobian.

% Ported from: cwaight_systems_paper_with_revisions2.tex
\begin{figure}[htbp]
    \centering
    \includegraphics[width=0.55\textwidth]{ch4_control_architecture.png}
    \caption{Hierarchical control architecture: adaptive navigation, cluster space, and robot layers, as in \cite{3,16}. The contents of the feature estimation and navigation controller blocks are the contributions of this dissertation.}
    \label{fig:ch4:control_architecture}
\end{figure}

\subsection{Robot Controller Layer}
\label{sec:ch4:robot_layer}

% Ported from: cwaight_systems_paper_with_revisions2.tex
Each robot's controller receives velocity commands and translates them into motion. Each robot also samples the local vector field $\mathbf{v}(\mathbf{p}_i) = (u_i, v_i)$ and reports its global position $\mathbf{p}_i$, providing the inputs to the estimation framework of Section~\ref{sec:ch4:estimation}.

% New text
Section~\ref{sec:ch2:dynamics} gives the first-order lag, speed limit, and
stiction floor identified for the Decabots, which the simulations of
later chapters reproduce.

\subsection{Cluster Space Controller Layer}
\label{sec:ch4:cluster_space}

% Ported from: IDETC 2025, Sec. 2.3, transcribed from the published PDF
To maintain a fixed formation while commanding all robots with individual
velocity commands, a cluster space controller was used. Cluster space
control is a technique developed by Kitts and colleagues at Santa Clara
University \cite{53}. In this technique, all robot positions are
described geometrically relative to a cluster frame via a set of
kinematic equations.

% New text
It is an operational space control approach in which the multirobot
formation is represented as a virtualized full degree-of-freedom
articulating mechanism.

% Ported from: cwaight_systems_paper_with_revisions2.tex
We refer to the group of robots as a cluster: a formation treated as a single virtual rigid body whose centroid, orientation, and internal geometry can be independently specified and controlled \cite{53,54}.

% Ported from: IDETC 2025, Sec. 2.3, transcribed from the published PDF
Cluster space control also provides a closed form method for converting
between cluster-level commands and robot-level commands. This allows for
the user or adaptive navigation layer to specify the behaviour commands
for the entire cluster, while the individual robots are then commanded by
the robot control layer to react accordingly. These conversions happen on
a centralized computer, after the individual robots communicate their
instantaneous position and sensed information.

\subsubsection{Three-Robot SAS Formation}
\label{sec:ch4:sas}

% Ported from: cwaight_systems_paper.tex
The cluster space controller maintains the cluster's formation geometry while executing centroid velocity commands from the adaptive navigation layer. It translates cluster-level velocity commands into individual robot velocity commands using an inverse kinematic Jacobian, and transforms robot position data back into cluster variables for feedback control.

% Ported from: cwaight_systems_paper.tex
Chapters~\ref{ch:zeroth} and~\ref{ch:first_order} consider a cluster of three robots forming a triangular configuration. The shape can be parameterized using Side-Angle-Side (SAS) representation with parameters $(p, \beta, q)$, where $p$ is the distance between robots 1 and 2, $\beta$ is the interior angle at robot 2, and $q$ is the distance between robots 2 and 3, as visualized in Fig.~\ref{fig:ch4:threerobot}.

% Ported from: cwaight_systems_paper.tex
\begin{figure}[htbp]
    \centering
    \includegraphics[width=0.45\textwidth]{ch4_SAS_Robots.png}
    \caption{Three-robot triangular formation parameterized by SAS variables: side lengths $p$, $q$ and interior angle $\beta$.}
    \label{fig:ch4:threerobot}
\end{figure}

% Ported from: cwaight_systems_paper.tex
The forward kinematic equations map the robot positions $(x_1, y_1), (x_2, y_2), (x_3, y_3)$ to the cluster shape parameters, and the inverse kinematic equations map cluster variables back to robot positions. Differentiating them gives analytic forms for the cluster space kinematic Jacobian and its inverse, which translate velocities between robot space and cluster space. The full equations are given in Appendix~\ref{app:kin3}.

\subsubsection{Six-Robot Pentagon-Plus-Center Formation}
\label{sec:ch4:pentagon}

% Ported from: Draft_14.tex
For six planar robots the cluster state has twelve variables. The centroid $(x_c, y_c)$
and heading $\theta_c$ place the formation, and nine shape variables set
its geometry. The robots form three pairs whose midpoints define an SAS
triangle $(p, \beta, q)$, and each pair $i = 2, 3, 4$ adds its length
$L_i$ and orientation $\theta_i$ (Fig.~\ref{fig:ch4:pentagon}). The layer maps
cluster velocities to robot velocities through the analytic inverse
Jacobian of the cluster kinematics (Appendix~\ref{app:kin6}).

% Ported from: Draft_14.tex
\begin{figure}[htbp]
    \centering
    \includegraphics[width=0.45\textwidth]{ch4_pentagon_formation.png}
    \caption{Pentagon-plus-center formation. Robot 1 sits at the centroid
    $\mathbf{p}_c$ and robots 2 to 6 on a ring (dotted). The cluster
    variables shown are the pair separations $L_2$, $L_3$, $L_4$, a pair
    orientation $\theta_3$, the SAS parameters $(p, \beta, q)$ of the
    triangle of pair midpoints (squares), and the heading $\theta_c$.}
    \label{fig:ch4:pentagon}
\end{figure}

\subsection{Adaptive Navigation Layer}
\label{sec:ch4:an_layer}

% Ported from: IDETC 2025, Sec. 2.4, transcribed from the published PDF
The adaptive navigation layer integrates the cluster shape policy,
feature estimator, and control law, determining the desired actions of
the entire cluster in cluster space variables.

% New text
The feature estimator is the part that changes from rung to rung.
Section~\ref{sec:ch4:estimation} describes it in general, and each of
Chapters~\ref{ch:zeroth} to~\ref{ch:second_order} gives the control law
that consumes it.

\section{Local Polynomial Field Estimation}
\label{sec:ch4:estimation}

% New text. Generalizes the first-order fit of the critical-points paper
% and the second-order fit of Draft 11 to order k; checked 2026-09-29.
Each estimator fits a polynomial to the velocity samples the robots take
at one instant. The order of that polynomial decides what the cluster can
learn about the field.

\subsection{Monomial Basis and Formation Matrix}
\label{sec:ch4:basis}

Let $\boldsymbol{\phi}_k(\mathbf{p})$ collect the monomials $x^a y^b$
with $a + b \leq k$, each scaled by $1/(a!\,b!)$, with coordinates
measured from the cluster centroid. The scaling makes each fitted
coefficient an estimate of the corresponding Taylor derivative of the field
at the centroid. The basis
has
\begin{equation}
    n_k = \tfrac{1}{2}(k+1)(k+2)
    \label{eq:ch4:nk}
\end{equation}
entries. Robot $i$ at $\mathbf{p}_i$ contributes the row
$\boldsymbol{\phi}_k(\mathbf{p}_i)^{\top}$ to the formation matrix
$\boldsymbol{\Phi}_k$, and each velocity component is fit separately,
\begin{equation}
    \boldsymbol{\Phi}_k\,\hat{\boldsymbol{\theta}}_u = \mathbf{m}_u, \qquad
    \boldsymbol{\Phi}_k\,\hat{\boldsymbol{\theta}}_v = \mathbf{m}_v ,
    \label{eq:ch4:fit}
\end{equation}
where $\mathbf{m}_u$ and $\mathbf{m}_v$ stack the robots' readings. The
first-order fit of Chapter~\ref{ch:first_order} is the case $k = 1$, and
the second-order fit of Chapter~\ref{ch:second_order} is $k = 2$. The two
chapters order the columns differently, which changes nothing in the fit.

Every robot reads the field with its own sensor, so the measurement noise
on one reading is independent of the noise on another. Suppose each
reading carries zero-mean noise of variance $\sigma^2$ and
$\boldsymbol{\Phi}_k$ has full column rank. The least-squares coefficients
of one velocity component are then unbiased, with covariance
\begin{equation}
    \operatorname{Cov}(\hat{\boldsymbol{\theta}}) =
    \sigma^{2}\,(\boldsymbol{\Phi}_k^{\top}\boldsymbol{\Phi}_k)^{-1},
    \label{eq:ch4:cov}
\end{equation}
conditional on the robot positions being known.

\subsection{Minimality and Degeneracy}
\label{sec:ch4:minimality}

Each robot adds one equation per component, so a unique fit needs at
least $n_k$ robots. The first order needs three and the second order
needs six. With exactly $n_k$ robots $\boldsymbol{\Phi}_k$ is square. It
is singular exactly when a nonzero coefficient vector $\mathbf{c}$
satisfies $\boldsymbol{\Phi}_k\mathbf{c} = \mathbf{0}$, which says that
the polynomial $\mathbf{c}^{\top}\boldsymbol{\phi}_k(\mathbf{p})$
vanishes at every robot. The fit therefore fails exactly when the robots
lie on a common algebraic curve of degree at most $k$. For $k = 1$ that
curve is a line, and for $k = 2$ it is a conic, which may be a pair of
lines. With more than $n_k$ robots the
system is overdetermined and is solved by least squares.

\subsection{Formation Size}
\label{sec:ch4:size}

Formation size enters through the columns of $\boldsymbol{\Phi}_k$. If a
formation keeps its shape and is scaled to radius $\rho$, a monomial of
degree $q$ scales as $\rho^{\,q}$, so the noise reaching a coefficient of
order $q$ scales as $\rho^{-q}$ (Section~\ref{sec:ch8:noise_order}).
Truncation error runs the other way. The neglected terms are of degree
$k + 1$, so the bias they leave grows with $\rho$. Under smoothness and a
fixed, full-rank normalized formation, the truncation error in an
order-$q$ derivative estimated by an order-$k$ fit generally behaves as
$O(\rho^{\,k+1-q})$, and leading terms can cancel by symmetry or by the
structure of the field. A larger formation therefore trades truncation
bias for noise suppression.

\section{The Fit-Order Ladder}
\label{sec:ch4:ladder}

Table~\ref{tab:ch4:ladder} organizes this dissertation by fit order, the
order to which the vector field itself is reconstructed from one set of
simultaneous measurements (Section~\ref{sec:ch1:thesis}). The
vector-to-scalar primitive fits a plane to the sensed magnitudes, which is
a first-order fit of a scalar field, but it reconstructs the vector field
only to zeroth order, so it sits on the lowest rung. Each rung fixes a
failure of the one below it. The zeroth-order primitives of
Chapter~\ref{ch:zeroth} use the mean sensed vector, or the gradient of
the sensed magnitude, and never assemble the velocity gradient. They
supply a heading, which drifts on an orbit. The first-order fit of
Chapter~\ref{ch:first_order} locates and classifies a critical point, but
it gives the velocity gradient at one point only, so one instantaneous
first-order fit does not by itself identify a curve to follow. The second-order fit of Chapter~\ref{ch:second_order} gives the
velocity gradient as a field over the formation, which is enough to
acquire and ride a separatrix. A third-order fit from ten robots would
recover the third derivatives the second-order trackers lack, and
Chapter~\ref{ch:limitations} leaves it as future work.

\begin{table}[htbp]
\centering
\caption{The fit-order ladder. The minimum robot count is $n_k$ of
(\ref{eq:ch4:nk}), and the degeneracy is the curve on which that many
robots cannot support the fit.}
\label{tab:ch4:ladder}
\begin{tabular}{>{\raggedright\arraybackslash}p{1.6cm} c >{\raggedright\arraybackslash}p{2.4cm} >{\raggedright\arraybackslash}p{5.2cm} c}
\hline
Fit order & Robots & Degenerate on & Enabled navigation capabilities & Chapter \\
\hline
Zeroth & 1 & none & heading only, which drifts on orbits & \ref{ch:zeroth} \\
First & 3 & common line & critical points, with location, type, and orbit & \ref{ch:first_order} \\
Second & 6 & common conic & separatrices, OECS, flow saddles as minima & \ref{ch:second_order} \\
Third & 10 & common cubic & third derivatives, the true $\mathbf{H}_D$ and transverse curvature & \ref{ch:limitations} \\
\hline
\end{tabular}
\end{table}

% ===================================================================
% CHAPTER 5: ZEROTH-ORDER BASELINES (IDETC 2025 testbed paper)
% ===================================================================
\chapter{Zeroth-Order Baselines}
\label{ch:zeroth}

\section{Direction-Only Navigation}
\label{sec:ch5:intro}

% New text
The primitives in this chapter steer on the robots' readings without
assembling the velocity gradient. The vector-sum primitive averages the
sensed vectors, and the vector-to-scalar primitive follows the slope of
the sensed magnitude. Both supply a heading. They were run on the testbed
of Chapter~\ref{ch:tools}, and they set the baseline for the higher rungs
of Section~\ref{sec:ch4:ladder}.

% Ported from: IDETC 2025 (DETC2025-167604), Sec. 2.5, transcribed from the published PDF
We implemented two adaptive navigation
primitives as control laws from \cite{26}. These were both shown as
center finding primitives that can find environmental extrema or follow a
closed orbit around the extrema. Both primitives were previously
validated in simulation, with known limitations documented. These
primitives rely solely on instantaneous information, enabling a purely
reactive approach to adaptive navigation control.

\section{Vector-Sum Primitive}
\label{sec:ch5:vector_sum}

% Ported from: IDETC 2025 (DETC2025-167604), Sec. 2.5, transcribed from the published PDF
The vector-sum technique combines the vectors sensed by each robot into a
resultant vector. The normalized direction of this resultant is then used
to specify $\dot{x}_c$ and $\dot{y}_c$ for the cluster.

\section{Vector-to-Scalar Primitive}
\label{sec:ch5:vector_to_scalar}

% Ported from: IDETC 2025 (DETC2025-167604), Sec. 2.5, transcribed from the published PDF
The vector-to-scalar technique uses only the magnitude information from
vector readings, effectively transforming vector field navigation into a
scalar field problem. This allows the application of scalar field
gradient ascent primitives. In the max-seeking implementation,
differences in vector magnitude readings from individual robots determine
the direction of steepest change.

% New text. The hardware implementation used this affine fit. The
% published IDETC paper printed an older cross-product construction
% (its Eqs. 2-7), which the author describes as confusing; it is not
% reproduced here.
The implementation fits a plane to the magnitudes $z_i =
\lVert\mathbf{v}(\mathbf{p}_i)\rVert$ read by the three robots,
\begin{equation}
    \mathbf{A}\,\hat{\boldsymbol{\theta}}_z = \mathbf{z}, \qquad
    \hat{\boldsymbol{\theta}}_z = [\hat{a}, \hat{b}, \hat{c}]^{\top},
    \label{eq:ch5:mag_fit}
\end{equation}
with the same formation matrix $\mathbf{A}$ as the first-order fit of
Section~\ref{sec:ch6:estimation}. The fitted slope
$(\hat{a}, \hat{b})$ estimates the gradient of the magnitude, and the
max-seeking behaviour moves the cluster along it. A minimum-seeking
behaviour moves against it. This is a first-order fit of a scalar. It
never assembles the Jacobian of the vector field, so it recovers a
direction of steepest change but no critical point.

\section{Hardware Trials on Fixed and Sinking Vortices}
\label{sec:ch5:trials}

\subsection{Setup}

% Ported from: IDETC 2025 (DETC2025-167604), Sec. 2.4, transcribed from the published PDF
For this research, we used a 3-robot cluster with a constant shape policy
as shown in Fig.~\ref{fig:ch4:threerobot}. The cluster parameters were set
to $p = 0.35$~m, $q = 0.35$~m, and $\beta = 1.05$~radians. These
parameters were chosen to form an equilateral triangle which maintains a
tight grouping and equidistant spacing. Cluster control was chosen to
maintain explicit formation control with full controllability and
observability which would not be easily achievable with alternative
approaches such as swarm architecture or leader-follower systems
\cite{44}. Neither the cluster orientation nor individual robot
orientations were specified, ensuring all velocity commands related
directly to the control primitives and formation control rather than
orientation maintenance.

The feature estimator converts data collected by the Decabots into usable
information for the control law, converting RGB sensor data into
navigation parameters. Section~\ref{sec:ch2:calibration} details the
feature estimator's operation.

%% The formation here is 0.35 m on a side; the critical point experiments
%% of Chapter 6 used 0.33 m. These are two experiments, not a conflict.

% Ported from: IDETC 2025 (DETC2025-167604), Sec. 3.2, transcribed from the published PDF
Two vector fields were constructed for these trials. The first studied
field was a fixed vortex, modeled as
\begin{equation}
    r = \sqrt{(x - c_x)^2 + (y - c_y)^2}, \qquad
    \theta = \operatorname{atan2}(y - c_y,\, x - c_x),
    \label{eq:ch5:rtheta}
\end{equation}
\begin{equation}
    u = r\sin\theta, \qquad v = -r\cos\theta ,
    \label{eq:ch5:fixed_vortex}
\end{equation}
where $(c_x, c_y)$ is the center of the map. Fig.~\ref{fig:ch5:fixed_map}
shows the vector field of (\ref{eq:ch5:rtheta})--(\ref{eq:ch5:fixed_vortex})
as a quiver plot overlaid on an HSV plot.

\begin{figure}[htbp]
    \centering
    \includegraphics[width=0.65\textwidth]{ch5_fixed_vortex_map.png}
    \caption{Quiver plot and HSV plot of the fixed vortex centered at the
    origin.}
    \label{fig:ch5:fixed_map}
\end{figure}

The other vector field studied was that of a sinking vortex. Equations
(\ref{eq:ch5:rtheta})--(\ref{eq:ch5:fixed_vortex}) are used to create the
fixed vortex, and then a sinking component is added to the field,
\begin{equation}
    u \leftarrow u - \frac{c_s\,(x - c_x)}{r^2 + 10^{-6}}, \qquad
    v \leftarrow v - \frac{c_s\,(y - c_y)}{r^2 + 10^{-6}},
    \label{eq:ch5:sinking}
\end{equation}
with $c_s = 0.2$.
%% The published Eqs. 12-13 print this term with a plus sign and 0.1. The
%% plotting package (trunk/robots_3/Results_for TestBed_Paper/
%% plotting_package/plotter_avg_run_Sinking_Vortex.m) uses the inward sign
%% and 0.2, and the trials spiral inward. Author decision 2026-09-29: use
%% the code values.
A small value of $10^{-6}$ was added to the denominator to avoid errors
when attempting to divide by zero. Once again, all vectors in the plot are
rescaled to the range of 0.3 to 1. The sinking vortex vector field
resulting from these equations is shown in Fig.~\ref{fig:ch5:sinking_map}
as a quiver plot overlaid on an HSV plot.

\begin{figure}[htbp]
    \centering
    \includegraphics[width=0.65\textwidth]{ch5_sinking_vortex_map.png}
    \caption{Quiver plot and HSV plot of the sinking vortex centered at the
    origin.}
    \label{fig:ch5:sinking_map}
\end{figure}

\subsection{Test Methodology}

% Ported from: IDETC 2025 (DETC2025-167604), Sec. 3.5, transcribed from the published PDF
A 3-robot cluster with fixed formation parameters ($p = 0.35$~m,
$q = 0.35$~m, and $\beta = 1.05$~radians) was used to test both control
primitives (vector-sum and vector-to-scalar) across the two vector field
environments. Throughout all tests, the OptiTrack motion capture system
recorded robot positions at a frequency of 10~Hz.

For each environment-primitive combination, 10 trial runs were conducted
to gather statistically meaningful data. The fixed vortex environment was
tested with both the vector-sum and vector-to-scalar minimum seeking
primitives, while the sinking vortex was tested exclusively with the
vector-sum primitive.

Both consistent and randomized starting positions were used to evaluate
the robustness of the control primitives. Performance evaluation focused
on three metrics, the robot trajectory relative to the expected path, the
steady-state error between the final cluster position and the target
position (where applicable), and the cluster's ability to maintain its
shape.

\subsection{Results}

% Ported from: IDETC 2025 (DETC2025-167604), Sec. 4.2, transcribed from the published PDF
The high accuracy achieved in hue and saturation interpretation ($R^2$
values of 0.96 and 0.91 respectively) enabled the testing of control
primitives in different environments that were previously only validated
in simulation.

The first set of test runs using a 3-robot cluster to navigate around the
center of a fixed vortex using the vector-sum primitive demonstrated an
unexpected but consistent physical phenomenon.
Fig.~\ref{fig:ch5:fixed_vs} shows the average path travelled over 10 runs
with the shaded region indicating 1 standard deviation, revealing very
high consistency between runs.

\begin{figure}[htbp]
    \centering
    \includegraphics[width=0.6\textwidth]{ch5_fixed_vortex_vector_sum.png}
    \caption{Average path of the robots in the fixed vortex using the
    vector-sum primitive, with one standard deviation shaded.}
    \label{fig:ch5:fixed_vs}
\end{figure}

The outward drift rate appears stable at all angles, a behavior not
predicted in the simulation studies of \cite{26}. This drift occurs
because the robots, following only their perceived direction using the
vector-sum technique, are always moving tangent to the circle, but each
robot answers a new heading late, after a delay and a lag that those
simulations did not model, so every step leaves the cluster slightly
outside the circle it was on. This real-world
validation highlights the importance of physical testing, as the physical
dynamics introduced nuances absent in simulated environments. The ability
to orbit a center with a fixed radius using adaptive navigation
techniques remained an open problem at this rung.

The second set of test runs were performed on the sinking vortex
environment using the vector-sum technique. Fig.~\ref{fig:ch5:sinking_vs}
shows the robots following a smooth path as they spiral inwards towards
the center as predicted in simulation. The solid lines show the average
run, and the blurred area shows 1 standard deviation. The smoothness can
be attributed to the very high $R^2$ value in interpreting hue as well as
to the stability of the center of this vortex field.

\begin{figure}[htbp]
    \centering
    \includegraphics[width=0.6\textwidth]{ch5_sinking_vortex_vector_sum.png}
    \caption{Average path of the robots in the sinking vortex with the
    vector-sum primitive.}
    \label{fig:ch5:sinking_vs}
\end{figure}

Fig.~\ref{fig:ch5:sinking_dist} shows the distance of the cluster center
relative to the origin with time. In these tests, the cluster settles
around 0.1~m away from the center. Looking at
Fig.~\ref{fig:ch5:sinking_vs}, the desired location falls within the
boundaries of the triangle that connects the 3 robots in the cluster.
This suggests that smaller cluster space variable values would lead to a
more accurate prediction.

\begin{figure}[htbp]
    \centering
    \includegraphics[width=0.6\textwidth]{ch5_sinking_vortex_distance.png}
    \caption{Average error of the cluster position relative to the center of
    the sinking vortex.}
    \label{fig:ch5:sinking_dist}
\end{figure}

For the third set of test runs, the vector-to-scalar minimum seeking
primitive was used to find the center of the fixed vortex floor map. The
robots are expected to follow a straight line to the center as seen in
Fig.~\ref{fig:ch5:fixed_v2s}. The trajectory of the robots is very
consistent across runs. The trajectory isn't a perfectly straight line,
as there may be noise in the printing, small fluctuations in robot
behavior, and some variance in the way that saturations affect estimates
of magnitudes.

\begin{figure}[htbp]
    \centering
    \includegraphics[width=0.6\textwidth]{ch5_fixed_vortex_vector_to_scalar.png}
    \caption{Average path of the robots in the fixed vortex with the
    vector-to-scalar primitive.}
    \label{fig:ch5:fixed_v2s}
\end{figure}

The cluster settled within 0.2~m of center (Fig.~\ref{fig:ch5:v2s_dist}),
suggesting the center fell within the circumscribed circle of the robot
triangle. The graph shows steady approach and stable positioning once the
desired location was reached.

\begin{figure}[htbp]
    \centering
    \includegraphics[width=0.6\textwidth]{ch5_vector_to_scalar_distance.png}
    \caption{Average error of the cluster position relative to the center of
    the fixed vortex with the vector-to-scalar primitive.}
    \label{fig:ch5:v2s_dist}
\end{figure}
%% The published caption of this figure (IDETC Fig. 18) says "sinking
%% vortex"; the text beside it and the data describe the fixed vortex.

For a fourth set of tests, the vector-to-scalar minimum seeking trials in
the fixed vortex were repeated with randomized starting positions on the
colormap. They showed a steady state error close to that of the trials
that all had the same starting location. This could be happening because
the velocity command is calculated from the gradient of saturation, and
not from the saturation itself. A small gradient could result in velocity
commands less than 0.05~m/s that are not sufficient for the robots to
overcome static friction.

\section{Formation Maintenance}
\label{sec:ch5:formation}

% Ported from: IDETC 2025 (DETC2025-167604), Sec. 4.3, transcribed from the published PDF
Throughout all experiments, the cluster space controller maintained
formation control, as evidenced by
Table~\ref{tab:ch5:formation}. The mean side lengths stayed within 1~cm of
the commanded values ($p = 0.35$~m and $q = 0.35$~m) and the mean angle
within 0.03~rad of $\beta = 1.05$~rad across all test environments, with
standard deviations of 0.9 to 3.2~cm in the sides and 0.045 to 0.12~rad in
the angle. This demonstrates the
effectiveness of the cluster space controller in maintaining desired
geometric relationships between robots while executing adaptive
navigation primitives.

\begin{table}[htbp]
\centering
\caption{Formation control performance over the primitive testing trial
runs.}
\label{tab:ch5:formation}
\begin{tabular}{lcccccc}
\hline
 & \multicolumn{4}{c}{Fixed vortex} & \multicolumn{2}{c}{Sinking vortex} \\
 & \multicolumn{2}{c}{Vector-sum} & \multicolumn{2}{c}{Vector-to-scalar} & \multicolumn{2}{c}{Vector-sum} \\
Variable & Average & $\sigma$ & Average & $\sigma$ & Average & $\sigma$ \\
\hline
$p$ & 0.36~m & 0.028 & 0.35~m & 0.015 & 0.35~m & 0.019 \\
$q$ & 0.36~m & 0.032 & 0.36~m & 0.009 & 0.35~m & 0.015 \\
$\beta$ & 1.03~rad & 0.122 & 1.02~rad & 0.045 & 1.02~rad & 0.056 \\
\hline
\end{tabular}
\end{table}

\section{Direction Without Location}
\label{sec:ch5:summary}

% New text
Both primitives reached their target when the field carried the cluster
there. On the sinking vortex the flow itself points inward, and the
vector-sum primitive settled about 0.1~m from the center. On the fixed
vortex the vector-to-scalar primitive descended the magnitude to within
0.2~m. Orbiting was different. A heading keeps the cluster tangent to the
flow, and nothing in it measures the distance to the center, so the late
turns carried the cluster steadily outward. Those turns are late because of
a delay and a lag that Section~\ref{sec:ch2:dynamics} fits to the hardware. Correcting that drift needs the
location of the center itself, which the three robots can estimate from
the same readings once they assemble the Jacobian.

% ===================================================================
% CHAPTER 6: FIRST-ORDER PRIMITIVES (critical-points paper)
% ===================================================================
\chapter{First-Order Primitives: Critical Points}
\label{ch:first_order}

\section{Introduction}
\label{sec:ch6:intro}

% New text
Chapter~\ref{ch:zeroth} showed that direction-only primitives reach a
feature only when the field itself carries the cluster there. Orbiting a
fixed vortex with the vector-sum primitive drifted steadily outward,
because nothing in a heading measures the distance to the center. This
chapter moves one rung up the ladder of Section~\ref{sec:ch4:ladder}.

% Ported from: cwaight_systems_paper_with_revisions2.tex
Recovering those coordinates is harder than following a heading, for reasons of both information and implementation. The information problem is that the location and classification of a critical point are absent from any single velocity sample; they are carried by the spatial derivatives of the field \cite{54}. A single platform must move, sample again, and difference its readings, which conflates the spatial gradient with the field's change in time and can return features that are not physically present. A distributed team removes this temporal aliasing by sampling at one instant \cite{28,29,44}, but recovering the derivatives is still not recovering the feature. Magnitude-based methods \cite{26} reduce the flow to a scalar and yield a descent direction, but the magnitude vanishes at every critical point regardless of type, and treating the two velocity components as independent scalar fields still does not locate the point. Both the coordinate and the identity are properties of the assembled linear fit. Inverting its Jacobian recovers the location where the field vanishes, and its eigenvalues determine the type \cite{9}.

% Ported from: cwaight_systems_paper_with_revisions2.tex
The implementation problem is that the local field must be reconstructed on physical hardware from noisy distributed measurements. The conditioning of this reconstruction degrades as the formation approaches collinearity, so the sampling shape must be held by a formation controller under real actuator dynamics and communication latency. Both the estimator and the mechatronics of the cluster limit the accuracy of the recovered features.

% Ported from: cwaight_systems_paper_with_revisions2.tex
This chapter develops a cooperative estimation strategy that recovers the local linear structure of the field. Using a minimum of three robots, the cluster estimates the local Jacobian in a single time step and treats critical point navigation as a local-to-global inference. Inverting the Jacobian predicts the location of the nearby critical point that organizes the surrounding flow, and its eigenvalues indicate the type, with no global map required. The same estimate supports two control laws, a proportional law that drives the cluster centroid onto the critical point, and an orbital law that holds the centroid at a commanded standoff radius. The sensing is spatially distributed across the robots, and in the present implementation a central computer gathers their readings, forms the estimate, and computes the commands.

\section{Critical Point Estimation from Three Robots}
\label{sec:ch6:estimation}

% Ported from: cwaight_systems_paper_with_revisions2.tex
Given a cluster of three robots sampling a planar vector field with robot positions $\mathbf{p}_i = (x_i, y_i)$ and corresponding vector field measurements $(u_i, v_i)$ for $i \in \{1, 2, 3\}$, we can approximate the vector field components across the cluster's spatial extent using linear planar models:
\begin{equation}\label{eq:ch6:2}
\hat{U}(x,y) = ax + by + c = u(x,y) + \mathcal{O}(\ell^2)
\end{equation}
\begin{equation}\label{eq:ch6:3}
\hat{V}(x,y) = dx + ey + f = v(x,y) + \mathcal{O}(\ell^2)
\end{equation}
where hats denote quantities estimated from the measurements and $\ell$ is the size of the region spanned by the robots and the point $(x,y)$. The $\mathcal{O}(\ell^2)$ term captures the field's second-order curvature, which a plane cannot represent.

The coefficients are determined by solving the linear systems:
\begin{equation}\label{eq:ch6:4}
\mathbf{A}\hat{\boldsymbol{\theta}}_u = \mathbf{z}_u, \quad \mathbf{A}\hat{\boldsymbol{\theta}}_v = \mathbf{z}_v
\end{equation}
where:
\begin{equation}\label{eq:ch6:5}
\mathbf{A} = \begin{bmatrix} x_1 & y_1 & 1 \\ x_2 & y_2 & 1 \\ x_3 & y_3 & 1 \end{bmatrix}, \quad \hat{\boldsymbol{\theta}}_u = \begin{bmatrix} a \\ b \\ c \end{bmatrix}, \quad \hat{\boldsymbol{\theta}}_v = \begin{bmatrix} d \\ e \\ f \end{bmatrix}
\end{equation}
and $\mathbf{z}_u = [u_1, u_2, u_3]^T$, $\mathbf{z}_v = [v_1, v_2, v_3]^T$ are the collected field measurements.

The estimated critical point $\hat{\mathbf{p}}^* = (\hat{x}^*, \hat{y}^*)$ satisfies $\hat{U}(\hat{x}^*, \hat{y}^*) = 0$ and $\hat{V}(\hat{x}^*, \hat{y}^*) = 0$, which in matrix form reads $\hat{\mathbf{J}}\hat{\mathbf{p}}^* = -\hat{\mathbf{h}}$, where the estimated Jacobian matrix and offset vector are:
\begin{equation}\label{eq:ch6:8}
\hat{\mathbf{J}} = \begin{bmatrix} a & b \\ d & e \end{bmatrix}, \quad \hat{\mathbf{h}} = \begin{bmatrix} c \\ f \end{bmatrix}
\end{equation}

Thus, provided $\det(\hat{\mathbf{J}}) \neq 0$:
\begin{equation}\label{eq:ch6:9}
\hat{\mathbf{p}}^* = -\hat{\mathbf{J}}^{-1}\hat{\mathbf{h}}
\end{equation}

% Ported from: cwaight_systems_paper_with_revisions2.tex
The eigenvalues of the Jacobian estimate $\hat{\mathbf{J}}$ can be used to identify the type of nearby critical point (Table~\ref{tab:ch3:field_summary}).

% Ported from: cwaight_systems_paper_with_revisions2.tex
Equation (\ref{eq:ch6:9}) returns the zero of the fitted affine field $\hat{\mathbf{v}}(\mathbf{p}) = \hat{\mathbf{J}}\mathbf{p} + \hat{\mathbf{h}} = \hat{\mathbf{J}}(\mathbf{p} - \hat{\mathbf{p}}^*)$, which need not be a zero of the true field. Evaluating it at any critical point $\mathbf{p}^*$ of the true field gives
\begin{equation}\label{eq:ch6:scope}
\hat{\mathbf{p}}^* - \mathbf{p}^* = -\hat{\mathbf{J}}^{-1}\hat{\mathbf{v}}(\mathbf{p}^*).
\end{equation}
The true field vanishes at $\mathbf{p}^*$, so $\hat{\mathbf{v}}(\mathbf{p}^*)$ is the error of the affine model there. If the true field is affine over a region containing both the cluster and $\mathbf{p}^*$, the models (\ref{eq:ch6:2})--(\ref{eq:ch6:3}) reproduce it without error and the recovery is exact at any formation size. On a smooth nonlinear field, the fit carries a linearization error that scales with the local curvature of the field; this error decreases as the formation contracts toward the critical point. The inference from $\hat{\mathbf{p}}^*$ to a true critical point is therefore valid only where the affine model remains accurate from the cluster out to $\mathbf{p}^*$.

The errors in $\hat{\mathbf{J}}$ and $\hat{\mathbf{h}}$ are correlated, since both come from the same measurements, but (\ref{eq:ch6:scope}) carries their combined effect through a single quantity, the model error $\hat{\mathbf{v}}(\mathbf{p}^*)$ at the true critical point. Taking norms turns the exact relation into a bound on how far the estimate can sit from the true point, $\|\hat{\mathbf{p}}^* - \mathbf{p}^*\| \le \|\hat{\mathbf{v}}(\mathbf{p}^*)\|/\sigma_{\min}(\hat{\mathbf{J}})$. The singular value $\sigma_{\min}(\hat{\mathbf{J}})$ is the slowest rate at which the fitted field changes in any direction, so where the field is nearly flat along some direction it is small and amplifies the model error.

The size of this model error can be bounded from the formation geometry. Let $\mathbf{p}_c = \frac{1}{3}\sum_i\mathbf{p}_i$ be the cluster centroid, $\mathbf{e} = \mathbf{p}_c - \mathbf{p}^*$ its displacement from the critical point, and $R$ the distance from $\mathbf{p}_c$ to each robot of the equilateral formation. Let $M$ bound the Hessian norms of $u$ and $v$ over the region holding the robots and $\mathbf{p}^*$, and let $\bar{\varepsilon}$ bound each robot's measurement error.

Two sources make a reading differ from the linearization of the field about $\mathbf{p}_c$. Field curvature contributes at most $MR^2/\sqrt{2}$, and the sensor contributes at most $\bar{\varepsilon}$. The fitted plane evaluated at $\mathbf{p}^*$ is a weighted sum of the three readings. The weights are the barycentric weights $w_i = \frac{1}{3} - \frac{2}{3}(\mathbf{p}_i - \mathbf{p}_c)^T\mathbf{e}/R^2$, and their magnitudes sum to at most $1 + \sqrt{2}\|\mathbf{e}\|/R$. The linearization about $\mathbf{p}_c$ also has its own error of at most $M\|\mathbf{e}\|^2/\sqrt{2}$ at $\mathbf{p}^*$. Adding the three contributions and merging the two curvature terms gives
\begin{equation}\label{eq:ch6:bound}
\|\hat{\mathbf{p}}^* - \mathbf{p}^*\| \le \frac{\overbrace{\tfrac{M}{\sqrt{2}}(R + \|\mathbf{e}\|)^2}^{\text{curvature}} + \overbrace{\bigl(1 + \tfrac{\sqrt{2}\|\mathbf{e}\|}{R}\bigr)\bar{\varepsilon}}^{\text{measurement error}}}{\sigma_{\min}(\hat{\mathbf{J}})}.
\end{equation}
Equation (\ref{eq:ch6:bound}) makes a design trade-off in formation sizing explicit, since the two terms move in opposite directions with $R$. A large formation raises the curvature term, which grows with the square of the largest distance from a robot to $\mathbf{p}^*$. A small formation raises the measurement term, which grows with the extrapolation distance $\|\mathbf{e}\|/R$ measured in formation sizes. The bound presumes that a true critical point lies nearby rather than establishing it.

Each robot carries its own sensor, so measurement noise is independent across robots. For zero-mean noise of covariance $\sigma_v^2\mathbf{I}_2$ at each robot on a locally affine field, linearizing (\ref{eq:ch6:scope}) in the noise gives
\begin{equation}\label{eq:ch6:cov}
\operatorname{Cov}(\hat{\mathbf{p}}^*) \approx \sigma_v^2\, c_w\, \mathbf{J}^{-1}\mathbf{J}^{-T},
\end{equation}
where $c_w = \sum_i w_i^2 = \frac{1}{3} + \frac{2}{3}\|\mathbf{e}\|^2/R^2$ for an equilateral formation. As the cluster reaches the critical point, $c_w \to 1/3$, which matches the variance of the mean of three independent readings.

\subsection{Minimum Robot Requirements}
\label{sec:ch6:minimum}

% Ported from: cwaight_systems_paper_with_revisions2.tex
Equations (\ref{eq:ch6:2})--(\ref{eq:ch6:3}) each contain three coefficients. From (\ref{eq:ch6:4}), each robot contributes one equation to each system: robot $i$ provides $u_i = ax_i + by_i + c$ for the U-component and $v_i = dx_i + ey_i + f$ for the V-component. To uniquely determine three unknowns requires at least three equations. Therefore: three robots are the minimum required for critical point estimation.


Three non-collinear robots are also sufficient: $\det(\mathbf{A})$ equals twice the signed area of the formation triangle, so $\mathbf{A}$ is invertible exactly when the robots are not collinear. With more than three robots ($n > 3$), the system becomes overdetermined and is solved via least-squares minimization.

\begin{proposition}[Minimality]\label{prop:min}
For an affine field $\mathbf{v}(\mathbf{p}) = \mathbf{J}(\mathbf{p}-\mathbf{p}^*)$ with $\det(\mathbf{J}) \neq 0$, locating $\mathbf{p}^*$ of every such field from instantaneous measurements requires at least three non-collinear robots, and three suffice.
\end{proposition}

\begin{proof}
Any $n \le 2$ robots, or any collinear robots, lie on a line $L = \{\mathbf{p}_0 + s\mathbf{d}\}$ and measure $\mathbf{m}_0 + s_i\mathbf{m}_1$, with $\mathbf{m}_0 = \mathbf{J}(\mathbf{p}_0 - \mathbf{p}^*)$ and $\mathbf{m}_1 = \mathbf{J}\mathbf{d}$. Take a field with $\mathbf{p}^* \notin L$. For any $\mathbf{q} \notin L$, the field $\mathbf{J}_q(\mathbf{p} - \mathbf{q})$ with $\mathbf{J}_q = [\,\mathbf{m}_0\ \ \mathbf{m}_1\,][\,\mathbf{p}_0 - \mathbf{q}\ \ \mathbf{d}\,]^{-1}$ is nonsingular and reproduces every measurement on $L$, so the measurements cannot determine $\mathbf{p}^*$. Its determinant $\det(\mathbf{J}_q) = \det(\mathbf{J})\det[\,\mathbf{p}_0 - \mathbf{p}^*\ \ \mathbf{d}\,]/\det[\,\mathbf{p}_0 - \mathbf{q}\ \ \mathbf{d}\,]$ changes sign as $\mathbf{q}$ crosses $L$, so the type is also undetermined. Conversely, for three non-collinear robots $\mathbf{A}$ is invertible, so $\hat{\mathbf{J}} = \mathbf{J}$ and $\hat{\mathbf{p}}^* = \mathbf{p}^*$, and the same holds for least squares with $n > 3$ robots that are not all collinear.
\end{proof}

\section{Control Laws}
\label{sec:ch6:control_laws}

% Ported from: cwaight_systems_paper_with_revisions2.tex
We develop two control primitives that drive the cluster centroid $\mathbf{p}_c$ toward or around the estimated critical point $\hat{\mathbf{p}}^*$.

\subsection{Critical Point Attraction}
\label{sec:ch6:attraction_def}

% Ported from: cwaight_systems_paper_with_revisions2.tex
To navigate the cluster center $\mathbf{p}_c$ to $\hat{\mathbf{p}}^*$, we apply a proportional control law:
\begin{equation}\label{eq:ch6:10}
\dot{\mathbf{p}}_c = k(\hat{\mathbf{p}}^* - \mathbf{p}_c)
\end{equation}
where $k > 0$ is the proportional gain. With $\hat{\mathbf{e}} = \mathbf{p}_c - \hat{\mathbf{p}}^*$ and a stationary estimate, the closed-loop error dynamics are $\dot{\hat{\mathbf{e}}} = -k\hat{\mathbf{e}}$, so $\hat{\mathbf{e}}(t) = \hat{\mathbf{e}}(0)e^{-kt}$. This guarantees global exponential convergence to the estimated critical point with time constant $1/k$, provided the estimate is stationary.

\subsection{Orbital Navigation}
\label{sec:ch6:orbital_def}

% Ported from: cwaight_systems_paper_with_revisions2.tex
To maintain a circular orbit of radius $r_d$ around $\hat{\mathbf{p}}^*$, we define the radial vector $\mathbf{r} = \hat{\mathbf{p}}^* - \mathbf{p}_c$ with magnitude $r = \|\mathbf{r}\|$. The unit radial vector is $\mathbf{n} = \mathbf{r}/r$, and the unit tangent vector $\mathbf{t}$ is obtained via a $-90^\circ$ rotation of $\mathbf{n}$. The orbital velocity command is:

\begin{equation}\label{eq:ch6:11}
\dot{\mathbf{p}}_c = k_t \mathbf{t} + k_r(r - r_d)\mathbf{n}
\end{equation}

Here, $k_t$ determines the orbital speed (where $\text{sgn}(k_t)$ dictates direction) and $k_r > 0$ drives radial convergence.
This control law decouples radial and angular dynamics. The radial error $e_r = r - r_d$ follows $\dot{e}_r = -k_r e_r$, ensuring exponential convergence to $r_d$. As $r \to r_d$, the angular velocity stabilizes to $\omega_o = k_t/r_d$, producing steady circular motion.

\section{Closed-Loop Stability and Tracking Bounds}
\label{sec:ch6:stability}

% Ported from: cwaight_systems_paper_with_revisions2.tex
Section~\ref{sec:ch6:control_laws} analyzed the control laws assuming a stationary estimate and ideal centroid kinematics. In practice, the estimate updates every cycle while the centroid responds through actuator delay, lag, and friction. Here we establish closed-loop stability and bound the resulting tracking errors.

\subsection{Sampled Loop}
\label{sec:ch6:sampled_loop}

% Ported from: cwaight_systems_paper_with_revisions2.tex
Let $\hat{\mathbf{p}}^*[j] = \mathbf{p}^* + \boldsymbol{\delta}[j]$ be the estimate at cycle $j$ with error $\boldsymbol{\delta}[j]$, so the attraction law commands $-k(\mathbf{e} - \boldsymbol{\delta})$, where $\mathbf{e} = \mathbf{p}_c - \mathbf{p}^*$. If the formation is held and the speed limit (inactive within $v_{\max}/k$ of $\hat{\mathbf{p}}^*$), the dead zone, and the lag ($\alpha \approx 0$) are neglected, the centroid follows (\ref{eq:ch2:momentum}) and its axes decouple. For each axis $a \in \{x, y\}$, the state $\boldsymbol{\eta}_a[j] = (e_a[j], e_a[j-1])^\top$ evolves as
\begin{equation}\label{eq:ch6:cl}
\boldsymbol{\eta}_a[j{+}1] = \mathbf{M}\,\boldsymbol{\eta}_a[j] + \mathbf{b}\,\bar{\delta}_a[j],
\end{equation}
where
\[
\mathbf{M} = \begin{bmatrix} 1-gk\Delta t(1-f) & -gk\Delta t f \\ 1 & 0 \end{bmatrix}, \quad \mathbf{b} = \begin{bmatrix} gk\Delta t \\ 0 \end{bmatrix},
\]
and $\bar{\delta}_a[j] = (1-f)\,\delta_a[j] + f\,\delta_a[j-1]$ is the estimate error blended over the period as in (\ref{eq:ch2:momentum}). Applying the Jury test to the characteristic polynomial $\det(z\mathbf{I} - \mathbf{M}) = z^2 - (1 - gk\Delta t(1-f))z + gk\Delta t f$ yields exponential stability if and only if $gk\Delta t f < 1$ and $gk\Delta t(1-2f) < 2$ for $k > 0$. For $f \ge 1/4$, the first condition binds; since $f\Delta t = L$, stability requires
\begin{equation}\label{eq:ch6:kmax}
0 < k < 1/(gL).
\end{equation}
For $g = 0.86$ and $L = 0.055$~s, the upper bound is 21.1~s$^{-1}$, and at $k = 1$~s$^{-1}$ the closed-loop poles have magnitudes 0.91 and 0.05. Discretizing (\ref{eq:ch2:continuous}) exactly, lag included, lowers the bound to 16.1~s$^{-1}$, and no split of the fitted $L+\tau$ lowers it below $1/(g(L+\tau)) = 15.5$~s$^{-1}$. A stable linear loop is input-to-state stable, so a bounded $\bar{\delta}_a$ and the dead zone, a command disturbance of at most $v_c$, leave the centroid error bounded. On a curved field the bound (\ref{eq:ch6:bound}) on $\|\boldsymbol{\delta}\|$ grows with $\|\mathbf{e}\|^2$, so the result there is local.

\subsection{Friction and Estimate Error}
\label{sec:ch6:friction}

% Ported from: cwaight_systems_paper_with_revisions2.tex
Because the attraction command $k\|\mathbf{p}_c - \hat{\mathbf{p}}^*\|$ falls below the static friction dead zone $v_c$ within $v_c/k$ of $\hat{\mathbf{p}}^*$, the cluster stops inside this friction ring (0.025~m at $k = 1$~s$^{-1}$). By (\ref{eq:ch6:bound}), the bound on $\|\boldsymbol{\delta}\|$ shrinks as $\mathbf{p}_c \to \mathbf{p}^*$. Under ideal kinematics, provided $\hat{\mathbf{p}}^*$ shifts slower than $\mathbf{p}_c$, the cluster settles without hunting to a final error $\|\mathbf{e}_{\text{final}}\| \le v_c/k + \|\boldsymbol{\delta}_{\text{final}}\|$.

\subsection{Near-Singular Estimates}
\label{sec:ch6:near_singular}

% Ported from: cwaight_systems_paper_with_revisions2.tex
Because the formation keeps the robots far from collinear, $\mathbf{A}$ stays well conditioned, and the estimate (\ref{eq:ch6:9}) is amplified most when $\det\hat{\mathbf{J}} \to 0$, as in near-uniform flow or along curves where $\det\mathbf{J} = 0$. These singular cases are detectable through the eigenvalues of the estimated Jacobian; a fallback could hold the last valid estimate or resize the cluster. The speed limit already bounds the cluster's response to a transient jump in the estimate. We evaluate (\ref{eq:ch6:9}) directly, without a conditioning threshold, to keep the baseline estimator's behavior transparent.

\subsection{Orbits}
\label{sec:ch6:orbits}

% Ported from: cwaight_systems_paper_with_revisions2.tex
The orbital law (\ref{eq:ch6:11}) is undefined at $r = 0$; if noise places $r < 10^{-6}$~m, the controller steps in a random direction, which clears the singularity on the next cycle. Furthermore, actuator delay and lag induce a steady-state radial offset even with an exact estimate. With orbital speed $V = gk_t$ and radial gain $gk_r$, the cluster follows the turning command a time $T \approx L + \tau + \Delta t/2$ late (the delay, the lag, and half a period of command hold), so in a steady orbit of radius $r^*$ it drifts outward at $TV^2/r^*$. Balancing this drift against the inward radial command $gk_r(r^* - r_d)$ gives the steady-state radius
\begin{equation}\label{eq:ch6:orbit_cf}
r^* = \tfrac{1}{2}\bigl(r_d + \sqrt{r_d^2 + 4C}\bigr), \qquad C = TV^2/(gk_r).
\end{equation}

\section{Control Performance in Simulation}
\label{sec:ch6:sim}

\subsection{Simulation Setup}
\label{sec:ch6:sim_setup}

% Ported from: cwaight_systems_paper_with_revisions2.tex
A Python-based simulation implementing the three-layer control architecture described in Chapter~\ref{ch:estimation} was developed to verify the critical point estimation and control primitives of this chapter. Unless stated otherwise, all simulation trials used a cluster of three robots in an equilateral triangular formation with a 0.30~m side length, matching the hardware configuration (Section~\ref{sec:ch6:hw}), with control gains of $k = 1$~s$^{-1}$ in (\ref{eq:ch6:10}), and $k_t = 0.15$~m/s, $k_r = 1$~s$^{-1}$ in (\ref{eq:ch6:11}). The robots follow the model of Section~\ref{sec:ch2:dynamics}, and the commanded centroid velocity was limited to 0.3~m/s on each axis.

Six noise-free analytical fields, one for each of the canonical critical point types listed in Table~\ref{tab:ch3:field_summary}, were used for trials, each with its critical point at the origin: sink $(-x,-y)$, source $(x,y)$, saddle $(-y,-x)$, vortex $(-y,x)$, sinking vortex $0.4(-y,x)-0.15(x,y)$, and spewing vortex $0.4(y,-x)+0.15(x,y)$. After performance in noise-free environments was verified, simulations were performed on two fields reconstructed from robot-sensed measurements collected on the hardware testbed (vortex and saddle). These reconstructed fields are a fixed spatial distortion of the analytical fields, fit to the robots' readings on the printed maps, and they add no time-varying noise.

Three nonlinear fields test how the estimate extrapolates. One cell of the double gyre, $\mathbf{v} = 0.1\pi(\cos\pi x \sin\pi y,\allowbreak\ -\sin\pi x \cos\pi y)$, has a center at the origin, saddles at $(\pm 0.5, \pm 0.5)$~m, and $\det\mathbf{J} = 0$ on the diamond $|x| + |y| = 0.5$~m. A curved saddle, $(x, -y)$ at the origin, bends the inflow axis to $v = -y + \chi y^3$, in an irrotational form, $u - iv = z + \chi z^3$ with $z = x + iy$, and a rotational form with stream function $S(x)S(y)$, where $S(s) = \sin(\gamma s)/\gamma$ and $\gamma^2 = 6\chi$. At the default $\chi = 1$, $\det\mathbf{J}$ vanishes only at $(0, \pm 0.58)$~m in the irrotational form and on the diamond $|x| + |y| = 0.64$~m in the rotational form.

\subsection{Metrics}
\label{sec:ch6:metrics}

% Ported from: cwaight_systems_paper_with_revisions2.tex
\emph{Attraction.} A trial of duration $T$ succeeds if the centroid $\mathbf{p}_c$ stays on the map and $\mathbf{p}^*$ lies inside the formation at time $T$; the remaining metrics use only successful trials. Accuracy is measured by the final error $e_f = \|\mathbf{p}_c(T) - \mathbf{p}^*\|$, reported by its mean $\bar{e}_f$, its standard deviation $s_e$, and the share of trials ending within 0.03, 0.04, and 0.05~m. Bias and precision separate where the trials end from how widely they scatter: the bias is the distance from $\mathbf{p}^*$ to the mean final position $\bar{\mathbf{p}}_c(T)$, the precision $\sigma$ is the RMS distance of the final positions from $\bar{\mathbf{p}}_c(T)$, and together they recover the mean squared error, since bias$^2 + \sigma^2$ equals the mean of $e_f^2$.

Three metrics describe the approach: the convergence time $T_{\mathrm{conv}}$ (the median time after which the centroid stays within 0.05~m of $\mathbf{p}^*$); the localization error (the mean of $\|\hat{\mathbf{p}}^* - \mathbf{p}^*\|$ over the last 20\% of a trial); and the control effort $E_u = \sum_j \|\mathbf{p}_c[j+1] - \mathbf{p}_c[j]\|^2/\Delta t$.

\emph{Orbit.} Each orbit run was scored on its radial error $\|\mathbf{p}_c - \mathbf{p}^*\| - r_d$ after the first 3~s. The mean of this error is the radial bias, its RMS value is the orbit RMSE, and its standard deviation is the within-run spread. Across repeated runs, the standard deviation of the run means gives the between-trial spread. A radius was considered practically achievable if the magnitude of its radial bias was at most $0.25\,r_d$. A run is lost if the centroid leaves the map or turns less than a quarter revolution about $\mathbf{p}^*$ after the first 3~s, and the statistics use only the runs that were not lost.

\emph{Classification.} Classification accuracy was measured as the fraction of cycles on which the type of $\hat{\mathbf{J}}$ matched the type of the critical point the trial reached or orbited.

\subsection{Critical Point Attraction}
\label{sec:ch6:sim_attraction}

% Ported from: cwaight_systems_paper_with_revisions2.tex
For each of the eleven fields, the cluster was initialized at a random position drawn uniformly from a 1~m $\times$ 1~m box centered on the target critical point and the simulation was run for 100 time steps at 10~Hz. This was performed 1000 times per field, from the same starting positions and headings in every field, and the final centroid position of the cluster in each trial was used to determine the metrics. In the gyre and the rotational saddle, trials started with all three robots inside the $\det\mathbf{J} = 0$ diamond around the target are scored separately.

 The cluster converged to the critical point in 100\% of 1000 trials for every analytical field at every threshold (Table~\ref{tab:ch6:simulation_results}). There the estimate is exact, and every trial settles on the friction ring of Section~\ref{sec:ch6:stability}. The reconstructed fields' fixed distortion moves the rest point off that ring and lowers success at 0.03~m, while every trial still ends within 0.05~m.

With the delay, lag, and friction removed, white noise up to 0.03~m/s changes the reconstructed fields' final error by at most 3.0~mm, so nearly all of their error comes from the fixed distortion: printing and sensing errors bend the field the robots read, the same way at the same place on every pass.

% Ported from: cwaight_systems_paper_with_revisions2.tex
\begin{table}[htbp]
\scriptsize
\centering
\caption{Simulation results on eleven fields. Brackets mark trials started with all three robots inside the $\det\mathbf{J} = 0$ diamond, and accuracy rows use successful trials. A dash marks a radius where most orbits were lost.}
\label{tab:ch6:simulation_results}
\setlength{\tabcolsep}{1.2pt}
\begin{tabular}{l|c|cc|ccc}
\hline
 & \textbf{Anal.} & \multicolumn{2}{c|}{\textbf{Recon.}} & \multicolumn{3}{c}{\textbf{Nonlinear}} \\
 & \textbf{(6)} & \textbf{Vort.} & \textbf{Sadd.} & \textbf{Gyre} & \textbf{Irrot.} & \textbf{Rot.} \\
\hline
\multicolumn{7}{l}{\textit{(a) Attraction}} \\
\hline
Succ.\ (\%) & 100 & 100 & 100 & 50 [100] & 100 & 72 [100] \\
Bias (m) & 0.0012 & 0.0031 & 0.0091 & 0.0017 & 0.0011 & 0.0012 \\
$\sigma$ (m) & 0.0226 & 0.0212 & 0.0242 & 0.0240 & 0.0224 & 0.0235 \\
$\bar{e}_f \pm s_e$ (mm) & 22.6$\pm$0.9 & 20.1$\pm$7.3 & 24.7$\pm$7.5 & 23.8$\pm$3.4 & 22.1$\pm$3.5 & 23.4$\pm$2.2 \\
Loc.\ error (m) & $<$0.0001 & 0.0101 & 0.0134 & 0.0043 & 0.0049 & 0.0026 \\
$T_{\mathrm{conv}}$ (s) & 2.3 & 2.4 & 2.5 & 1.9 & 2.4 & 2.1 \\
$E_u$ (m$^2$/s) & 0.068 & 0.065 & 0.066 & 0.042 & 0.070 & 0.055 \\
$\le$\,3/4/5~cm (\%) & 100 & 91/100/100 & 68/100/100 & 99/100/100 & 100 & 100 \\
\hline
\multicolumn{7}{l}{\textit{(b) Orbit radial bias (m), by $r_d$ (m)}} \\
\hline
\quad 0.01 & +0.042 & +0.039 & +0.034 & +0.044 & +0.042 & +0.043 \\
\quad 0.10 & +0.018 & +0.006 & 0.000 & +0.018 & +0.018 & +0.018 \\
\quad 0.20 & +0.010 & +0.038 & +0.045 & $-$0.013 & +0.011 & $-$0.005 \\
\quad 0.30 & +0.007 & +0.048 & +0.050 & $-$0.064 & +0.021 & $-$0.044 \\
\quad 0.40 & +0.005 & +0.046 & +0.043 & -- & -- & $-$0.097 \\
\quad 0.50 & +0.004 & +0.027 & +0.019 & -- & -- & -- \\
\hline
\multicolumn{7}{l}{\textit{(c) Classification (\%)}} \\
\hline
No noise & 100 & 73 & 100 & 100 & 100 & 100 \\
Noise 0.02~m/s & 86--100 & 52 & 100 & 83 & 100 & 100 \\
\hline
\end{tabular}
\end{table}

% Ported from: cwaight_systems_paper_with_revisions2.tex
In the nonlinear fields, curvature affected both the localization error and which critical point the cluster reached. Trials that reached the target settled as tightly as in the analytical fields, despite the estimate staying 2.6 to 4.9~mm from the true critical point. Evaluating the fitted plane at the centroid, $\hat{\mathbf{v}}(\mathbf{p}_c) = \hat{\mathbf{J}}(\mathbf{p}_c - \hat{\mathbf{p}}^*)$, and substituting into (\ref{eq:ch6:10}) rewrites the attraction law as $\dot{\mathbf{p}}_c = -k\hat{\mathbf{J}}^{-1}\hat{\mathbf{v}}(\mathbf{p}_c)$. This is a continuous-time Newton step toward a zero of the local fit. Because a continuous Newton flow cannot cross a $\det\hat{\mathbf{J}} = 0$ contour, this curve acts as a basin boundary. The initial sign of $\det\hat{\mathbf{J}}$ determined the critical point type reached in all 2000 gyre and rotational-saddle trials (Fig.~\ref{fig:ch6:nonlinear}(a)). Every trial initialized with all three robots inside the diamond converged to the target.

% Ported from: cwaight_systems_paper_with_revisions2.tex
\begin{figure}[htbp]
    \centering
    \includegraphics[width=0.8\textwidth]{ch6_fig_nonlinear_triangles_2026.pdf}
    \caption{The gyre cell in simulation, with the $\det\mathbf{J} = 0$ diamond dashed and the formation to scale. (a) Attraction starting positions, blue where the cluster reached the center and orange where it reached a saddle. The orange trial starts with a robot outside the diamond. (b) An orbit held at $r_d = 0.3$~m and one lost at 0.4~m.}
    \label{fig:ch6:nonlinear}
\end{figure}

\subsection{Orbital Navigation}
\label{sec:ch6:sim_orbit}

% Ported from: cwaight_systems_paper_with_revisions2.tex
Orbital control was evaluated on the same eleven fields at different commanded radii. For each radius, the cluster was positioned at the desired radius to initialize the run, then the orbital controller (\ref{eq:ch6:11}) was used for 300 time steps (600 on the nonlinear fields). The mean radial error and standard deviation were computed after discarding the first 30 time steps (3~s) as a settling transient. Each run was repeated 100 times.


    Table~\ref{tab:ch6:simulation_results}(b) presents the orbital tracking results. The six noise-free fields produced identical performance at each radius, with standard deviations below 0.001~m. This is expected: on an exactly linear noise-free field the three-robot fit recovers the critical point exactly at every step, so the controller input, and hence the trajectory, is independent of the field type. Mean radial errors had a positive bias, and decreased with increasing radius. The reconstructed fields share the bias at $r_d = 0.01$~m and nearly cancel it at 0.10~m. From 0.20~m their bias lies between 0.019 and 0.050~m, with within-run standard deviations that increase with the commanded radius. By the $0.25\,r_d$ criterion the practical range is $r_d \ge 0.10$~m on the analytical and reconstructed fields.

In the nonlinear fields, curvature shifted the steady-state orbit radius, and orbits were lost once the formation reached a set where $\det\mathbf{J} = 0$ (Fig.~\ref{fig:ch6:nonlinear}(b)). The gyre's speed grows more slowly than in proportion to the distance from its center, so the slope the fit measures at the formation is shallower than the average slope back to the center. Extrapolating that slope places $\hat{\mathbf{p}}^*$ beyond $\mathbf{p}^*$, and the cluster orbited inside $r_d$, by 0.064~m at $r_d = 0.3$~m. Varying the curvature $\chi$ from $-1$ to 1 shifted the rotational saddle's mean radius by up to 0.17~m, while the irrotational saddle's opposite shifts along its two axes cancelled to within 0.01~m.

\subsection{Classification}
\label{sec:ch6:sim_class}

% Ported from: cwaight_systems_paper_with_revisions2.tex
Table~\ref{tab:ch6:simulation_results}(c) gives the classifier's accuracy over the last 2~s of each attraction trial. On the noise-free analytical and nonlinear fields, it labeled the reached critical point correctly on every cycle. With white measurement noise of 0.02~m/s on each reading, accuracy held at 94 to 95\% on the nodes and spirals and 100\% on the saddles. Accuracy on the centers fell furthest: to 86\% in the vortex and 83\% in the gyre. On the reconstructed vortex, where the fixed distortion biases $\hat{\mathbf{J}}$, it fell to 73\% without noise and 52\% with it. Because a center ($\operatorname{tr}\mathbf{J} = 0$) lies on the boundary between stable and unstable spirals, and the canonical sink and source ($\mathbf{J} = \mp\mathbf{I}$) lie on the boundary between nodes and spirals, small errors in $\hat{\mathbf{J}}$ can push their eigenvalues across those boundaries, whereas a saddle's eigenvalues of opposite sign lie far from every boundary. The classifier counts an eigenvalue pair within about $6^\circ$ of the imaginary or real axis as a center or a node, but noise still carries some estimates across the edge of that band. Every mislabel was a swap with a neighboring type in Table~\ref{tab:ch3:field_summary}.

\section{Hardware Validation}
\label{sec:ch6:hw}

\subsection{Hardware Testbed and Field Representation}
\label{sec:ch6:hw_testbed}

% New text
The testbed, the printed color maps and their sensor calibration, and the
field reconstructions behind the reconstructed simulation fields are
described in Sections~\ref{sec:ch2:testbed}, \ref{sec:ch2:hsv},
\ref{sec:ch2:calibration}, and~\ref{sec:ch2:reconstruction}.

% Ported from: cwaight_systems_paper_with_revisions2.tex
The hardware controller also limited the commanded centroid velocity to 0.3~m/s per axis. In the orbit runs, the robot model of Section~\ref{sec:ch2:dynamics} reproduced the measured velocities with an RMS error of 0.029~m/s.

\subsection{Validation Test Plan}
\label{sec:ch6:hw_plan}

% Ported from: cwaight_systems_paper_with_revisions2.tex
Hardware validation focused on two representative field types that encompass the key challenges of vector field navigation: rotational dynamics (vortex) and hyperbolic instability (saddle). The vortex field tests pure rotational dynamics and provides continuity with previous studies in our group~\cite{28,26}. The saddle is the only canonical field type that contains separatrices, which act as Lagrangian coherent structures in steady flows. Testing on the saddle field validates that the framework operates reliably in the presence of these structures.

Two formations were used. The three-robot equilateral triangle held a mean side of about 0.3~m. The four-robot square estimates $\hat{\mathbf{J}}$ by an overdetermined least squares fit to its four readings. Its diagonals were commanded to 0.30~m and measured 0.33~m in attraction and 0.31~m in orbit.

For each field, eight starting locations were selected around the perimeter of the workspace and the attraction controller (\ref{eq:ch6:10}) was run 10 times from each location for 10~s in randomized order, giving 80 trials per field with three robots and 79 vortex trials with four. To test the orbiting controller (\ref{eq:ch6:11}), each field was run at six commanded radii from 0.01 to 0.50~m with three robots, matching simulation, and at five radii from 0.01 to 0.40~m with four, ten times per radius. Each orbit run lasted 200 to 300 cycles at 10~Hz.

Each hardware run was also simulated from its own starting position, with the formation it held and the same commanded radius and duration, on the analytical and the reconstructed field. The classifier of Section~\ref{sec:ch6:sim_class} labels $\hat{\mathbf{J}}$ on every cycle without feeding the controller. Every trial is scored with the metrics of Section~\ref{sec:ch6:metrics}.

\subsection{Critical Point Attraction}
\label{sec:ch6:hw_attraction}

% Ported from: cwaight_systems_paper_with_revisions2.tex
Following the test plan above, 80 vortex and 80 saddle trials with three robots and 79 vortex trials with four were analyzed. Due to file management errors, one four-robot vortex trial was lost.

The three-robot cluster converged to the critical point (Fig.~\ref{fig:ch6:vortextomiddle}) in 99\% of the saddle trials and 91\% of the vortex trials at the 0.05~m threshold, and in 80\% and 39\% at 0.03~m. No trial left the map, and the critical point ended inside the robots' triangle in 100\% of trials. Table~\ref{tab:ch6:estimation_results} summarizes the results. In the vortex field, systematic bias was 0.007~m with precision of 0.036~m. In the saddle field, systematic bias was 0.002~m with precision of 0.025~m. The bias is smaller than the mean final error because final positions on opposite sides of $\mathbf{p}^*$ cancel in the mean, and the estimate ended closer to $\mathbf{p}^*$ than the cluster.

With four robots in the vortex, 100\% of trials ended within 0.05~m and 65\% within 0.03~m. Precision improved to 0.014~m from 0.036~m with three, but the bias of 0.024~m matched the localization error of 0.025~m.

% Ported from: cwaight_systems_paper_with_revisions2.tex
\begin{figure}[htbp]
    \centering
    \includegraphics[width=0.9\textwidth]{ch6_fig_hw_attraction_2026.pdf}
    \caption{Convergence trajectories from eight starting positions to the critical point of the vortex (a) and saddle (b), with each run thin and the mean path from each starting position thick. (c) Distance to $\mathbf{p}^*$ over time (log scale), mean and one standard deviation across runs, for three robots in both fields and four in the vortex; the friction ring $v_c/k$ of Section~\ref{sec:ch6:stability} is dashed.}
    \label{fig:ch6:vortextomiddle}\label{fig:ch6:Tomiddlewithtime}
\end{figure}

% Ported from: cwaight_systems_paper_with_revisions2.tex
\begin{table}[htbp]
    \scriptsize
    \centering
    \caption{Hardware experimental results, three robots (3R) and four (4R), with matched simulations of each attraction trial.}
    \label{tab:ch6:estimation_results}\label{tab:ch6:orbital_radius_comparison}
    \setlength{\tabcolsep}{3pt}
    \begin{tabular*}{0.8\textwidth}{@{\extracolsep{\fill}}l|cc|cc}
\hline
 & \multicolumn{2}{c|}{\textbf{Vortex}} & \multicolumn{2}{c}{\textbf{Saddle}} \\
 & \textbf{3R} & \textbf{4R} & \textbf{3R} & \textbf{4R} \\
\hline
\multicolumn{5}{l}{\textit{(a) Attraction}} \\
\hline
Runs & 80 & 79 & 80 & -- \\
Bias (mm) & 7 & 24 & 2 & -- \\
$\sigma$ (mm) & 36 & 14 & 25 & -- \\
$\bar{e}_f \pm s_e$ (mm) & 34$\pm$14 & 26$\pm$9 & 23$\pm$9 & -- \\
Loc.\ error (mm) & 15 & 25 & 14 & -- \\
$T_{\mathrm{conv}}$ (s) & 3.9 & 2.8 & 3.4 & -- \\
$E_u$ (m$^2$/s) & 0.187 & 0.218 & 0.204 & -- \\
$\le$\,3/4/5 cm (\%) & 39/69/91 & 65/86/100 & 80/98/99 & -- \\
Sim.\ $\bar{e}_f$ anal./rec.\ (mm) & 23/19 & 22/19 & 23/26 & -- \\
\hline
\multicolumn{5}{l}{\textit{(b) Orbit radial bias $\pm$ between-trial spread (mm), by $r_d$ (m)}} \\
\hline
\quad 0.01 & +60$\pm$11 & +57$\pm$11 & +66$\pm$15 & +57$\pm$14 \\
\quad 0.10 & +28$\pm$2 & +28$\pm$7 & +36$\pm$11 & +27$\pm$6 \\
\quad 0.20 & +7$\pm$14 & +36$\pm$25 & $-$1$\pm$5 & +31$\pm$25 \\
\quad 0.30 & +43$\pm$27 & +98$\pm$20 & +33$\pm$16 & +40$\pm$19 \\
\quad 0.40 & +73$\pm$23 & +98$\pm$8 & +32$\pm$10 & $-$13$\pm$73 \\
\quad 0.50 & +20$\pm$9 & -- & +5$\pm$11 & -- \\
\hline
\multicolumn{5}{l}{\textit{(c) Orbit within-run spread (mm), by $r_d$ (m)}} \\
\hline
\quad 0.01 & 20 & 18 & 12 & 23 \\
\quad 0.10 & 14 & 17 & 12 & 22 \\
\quad 0.20 & 36 & 38 & 15 & 32 \\
\quad 0.30 & 57 & 57 & 56 & 55 \\
\quad 0.40 & 69 & 55 & 71 & 76 \\
\quad 0.50 & 76 & -- & 83 & -- \\
\hline
\end{tabular*}
    \end{table}

\subsection{Orbital Navigation}
\label{sec:ch6:hw_orbit}

% Ported from: cwaight_systems_paper_with_revisions2.tex
Orbital control was evaluated on both fields at commanded radii from 0.01~m to 0.50~m with three robots and to 0.40~m with four. For each radius, the orbital controller (\ref{eq:ch6:11}) was run for 200 to 300 cycles at 10~Hz, and the mean radial error and its standard deviation were recorded after the first 3~s, as in simulation (Table~\ref{tab:ch6:orbital_radius_comparison}).

 With three robots in the vortex field, mean radial errors ranged from 0.007~m at a commanded radius of 0.20~m to 0.073~m at 0.40~m, with within-run standard deviations ranging from 0.014 to 0.076~m. In the saddle field, mean radial errors ranged from $-0.001$~m at 0.20~m to 0.066~m at 0.01~m, with within-run standard deviations increasing from 0.012 to 0.083~m. From $r_d = 0.10$~m upward the fluctuation within one orbit exceeds the variation between repeats. By the $0.25\,r_d$ criterion the practical range is $r_d \ge 0.20$~m in both fields. Below these radii the controller holds a bounded orbit but does not track the commanded radius.

The four-robot square behaved like the triangle at small radii, with a bias of +0.057~m at $r_d = 0.01$~m in both fields and a similar within-run spread in the vortex, but nearly twice the triangle's in the saddle (Table~\ref{tab:ch6:estimation_results}(c)). Its practical range begins at 0.20~m in the saddle; in the vortex it met the $0.25\,r_d$ criterion at 0.20 and 0.40~m but not at 0.30~m, where its bias was +0.098~m. Unlike in attraction, the fourth robot did not improve the orbits: from $r_d = 0.20$~m the square's bias exceeded the triangle's at every radius in the vortex and at 0.20 and 0.30~m in the saddle. The matched simulations predicted no such loss: on the reconstructed fields from $r_d = 0.20$~m the square's bias stayed within 8~mm of the triangle's and its within-run spread was 16 to 43\% smaller; because the model reads one fixed distortion through identical, noise-free sensors, the square's shortfall on hardware points to an error the model omits. No orbit run with either formation left the map.

\subsection{Classification}
\label{sec:ch6:hw_class}

% Ported from: cwaight_systems_paper_with_revisions2.tex
Across every cycle of the orbit runs, the classifier matched the printed saddle on 100\% (three robots) and 99.6\% (four robots) of cycles. In the vortex, it identified rotational flow on 99.9\% and 99.6\% of cycles, respectively, but labeled only 45\% and 47\% strictly as a center; as in Section~\ref{sec:ch6:sim_class}, field distortion and measurement noise introduce small real parts in the eigenvalues that push the center across its boundary into a spiral.

\subsection{Offline Classification of the Attraction Trials}
\label{sec:ch6:classification_offline}

% New text. Analysis: experiments/classify_from_hardware_logs.py, results
% in experiments/outputs/classify_hw/summary.json (VF_Robot), run
% 2026-09-29. Those logs are the three-robot attraction trials of the
% revision (2025 campaign) less saddle runs 43 and 64, which the extras
% later filled. It has not been re-run with the extras.
The same fit that locates the critical point also names its type. The
eigenvalues of $\hat{\mathbf{J}}$ place it in one of the classes of
Table~\ref{tab:ch3:field_summary}, at no extra sampling cost. The
attraction trials were designed to test location, but every control cycle
logged the positions and sensed readings needed to recompute
$\hat{\mathbf{J}}$. Recomputing it offline reproduces the controller's
logged critical point estimates to within $5\times10^{-7}$~m, so the
Jacobians classified here are the ones the controller computed. The
logs hold 78 saddle trials and 80 vortex trials of 101 cycles each. These
are the three-robot attraction trials of
Section~\ref{sec:ch6:hw_attraction}, less saddle runs 43 and 64, which
were filled later by extra runs.

Each estimate was classified as a saddle when $\det\hat{\mathbf{J}} < 0$,
a node when the eigenvalues are real and of one sign, and a rotational
type when they form a complex pair $\alpha \pm i\nu$. A measured center
never has $\alpha$ exactly zero, so a complex pair was labeled a center
when $|\alpha| < 0.2\,\nu$ and a spiral otherwise. That threshold is a
choice, and Table~\ref{tab:ch6:classification_hw} also gives the results
at 0.1 and 0.3.

\begin{table}[htbp]
\centering
\caption{Offline classification of the hardware Jacobian estimates.
The terminal window is the final 2~s of each trial, and the trial vote is
the majority label over that window.}
\label{tab:ch6:classification_hw}
\begin{tabular}{lcc}
\hline
 & Saddle map & Vortex map \\
\hline
Trials & 78 & 80 \\
Class correct (saddle, rotational, node), all cycles & 99.9\% & 99.9\% \\
Class correct, terminal window & 100\% & 100\% \\
Class correct, trial vote & 78/78 & 80/80 \\
Exact type, terminal window, threshold 0.1 & 100\% & 71\% \\
Exact type, terminal window, threshold 0.2 & 100\% & 89\% \\
Exact type, terminal window, threshold 0.3 & 100\% & 99\% \\
\hline
\end{tabular}
\end{table}

The estimate named the class of both fields correctly in effectively
every cycle, including cycles more than 0.8~m from the critical point.
Every vortex error is a stable spiral with a small negative real part,
and none is a saddle or a node. A small damping in the printed map or a
small bias in the fit is enough to turn a center into a weak spiral, and
the data cannot separate the two.

Classification carries one condition that location does not. A
reflection of the sensed vectors reverses the sign of
$\det\hat{\mathbf{J}}$, which exchanges saddles with nodes but leaves the
estimated critical point where it was. The two printed maps were placed
under different frame conventions, including reflections, rotations, and
camera mirroring. The results above use the registration from the field
reconstructions of Section~\ref{sec:ch2:reconstruction}, which negates
the sensed $v$ component for the saddle map only. The controller applied
no registration, since it used the estimate only for location, and its
unregistered Jacobian read the saddle map as a stable node or spiral
throughout the terminal window. A deployment that reports feature type
therefore needs its sensor frame registered to its navigation frame.

\section{Discussion}
\label{sec:ch6:discussion}

\subsection{Error Analysis}
\label{sec:ch6:error}

% Ported from: cwaight_systems_paper_with_revisions2.tex
Estimation and control errors arise from four sources: robot dynamics, field distortion and measurement noise, cluster geometry, and field conditioning. Sections~\ref{sec:ch6:estimation} and~\ref{sec:ch6:stability} bound the effect of each, and this section traces the errors of Sections~\ref{sec:ch6:sim} and~\ref{sec:ch6:hw} to them.

Positive radial bias at small orbital radii is a result of robot dynamics. At small radii, the commanded angular rate $k_t / r_d$ is high relative to the radial correction bandwidth $gk_r$, and the outward drift from the delay and lag sets the bias (\ref{eq:ch6:orbit_cf}). The closed form matches simulation within 3~mm at every commanded radius at the settings of Section~\ref{sec:ch6:sim_setup}. Across sweeps of the gains, delay, lag, sample period, and dead zone, it matches to a median of 1.2~mm. On hardware at $r_d = 0.01$~m the vortex orbit's radial bias was 0.071~m about $\hat{\mathbf{p}}^*$, which the orbital law regulates, and 0.060~m about $\mathbf{p}^*$, so the small-radius offset arises in the closed loop.

Comparing the noise-free and reconstructed fields in simulation isolates the effect of field distortion: the fixed distortion of the printed maps appears to be a driving factor of the simulation-reality gap (Section~\ref{sec:ch6:sim_attraction}). Time-varying noise is far smaller; with the formation at rest, a robot's reading changes by about $2\times10^{-4}$~m/s between cycles, while the calibration error of Section~\ref{sec:ch2:calibration} is about 0.02~m/s and fixed at each location. In the final error of Section~\ref{sec:ch6:stability}, friction sets the spread through $v_c/k$, and the distortion sets the bias.

The simulations matched to each hardware run show how much of the hardware error the model reproduces. They reproduce the saddle's attraction error, 0.026~m against 0.023~m, while the hardware vortex ends 0.4 to 1.5~cm farther from $\mathbf{p}^*$ (Table~\ref{tab:ch6:estimation_results}(a)). In orbit at $r_d = 0.01$~m the hardware bias exceeds the matched reconstructed-field simulations by 14 to 19~mm. In these runs the live estimate ran 1.7 to 2.8~cm ahead of the cluster along its motion, a lead the model omits. From $r_d = 0.30$~m the reconstructed fields produce 27 to 63\% of the hardware within-run spread and the analytical fields at most 1~mm, so the fixed distortion of the printed maps accounts for about half of it. Some errors may also have arisen from modeling discrepancies such as modeling all robots identically with the same dead zone, speed limit, and time constant, though their fitted gains range from 0.80 to 0.89.

To avoid errors from collinear robots (the condition number of $\mathbf{A}$ becomes infinite), an equilateral triangle was chosen as the cluster shape. This shape minimizes the condition number of $\mathbf{A}$ in (\ref{eq:ch6:4}) among formations of the same size when the origin is at the centroid, making it optimal for sampling.

Further errors could arise from poor conditioning of the estimated Jacobian, which amplifies coefficient errors into errors in the critical point location (\ref{eq:ch6:9}) by the factor $1/\sigma_{\min}(\hat{\mathbf{J}})$ in (\ref{eq:ch6:bound}). Unlike the condition number of $\mathbf{A}$, this factor is an intrinsic property of the field and is not under the designer's control. A center is non-hyperbolic yet has $\det\mathbf{J} > 0$, so (\ref{eq:ch6:9}) locates it as well as any other critical point.

\subsection{The Orbit as a Limit Cycle}
\label{sec:ch6:limit_cycle}

% Ported from: cwaight_systems_paper_with_revisions2.tex
The orbit converges not to the commanded circle but to a limit cycle of the closed loop, at every commanded radius in simulation. The cycle is clearest on hardware at $r_d = 0.10$~m in the vortex, where ten three-robot runs (four started inside the cycle and six outside) met on one curve (Fig.~\ref{fig:ch6:limit_cycle}(b)) with mean radius 0.127~m and 2.9~mm between runs, against 0.111~m for the simulated cycle on the reconstructed vortex (Fig.~\ref{fig:ch6:limit_cycle}(a)). The closed form (\ref{eq:ch6:orbit_cf}) sets the cycle radius at $r^* = 0.120$~m (a circle in the analytical fields). The printed map's fixed distortion bends the cycle out of round, so its radius varies around the loop by 0.029~m in simulation and 0.043~m on hardware.

% Ported from: cwaight_systems_paper_with_revisions2.tex
\begin{figure}[htbp]
    \centering
    \includegraphics[width=0.9\textwidth]{ch6_fig_limit_cycle_2026.pdf}
    \caption{Limit cycle in the vortex at $r_d = 0.10$~m (dotted circle). (a) Simulation on the reconstructed vortex from starting positions inside (blue) and outside (red) the cycle, which meet on one closed curve (black). (b) Ten hardware runs, with the hardware cycle (black) and the simulated cycle (dashed).}
    \label{fig:ch6:limit_cycle}
\end{figure}

\subsection{Comparison with the Vector-to-Scalar Primitive}
\label{sec:ch6:comparison}

% Ported from: cwaight_systems_paper_with_revisions2.tex
The modular architecture allows navigation controllers to be interchanged while cluster and robot controller layers remain fixed. As a result, performance differences between controllers reflect algorithmic rather than implementation variations.
The vector-to-scalar primitive \cite{26} reduces the flow to its magnitude, and a plane fit to the speeds $S_i = \|\mathbf{v}(\mathbf{p}_i)\|$ at the three robots gives their gradient $\mathbf{g}$. For attraction the cluster descends $\mathbf{g}$, and for orbits the primitive was modified to move tangentially to the direction of steepest descent. To separate what the primitive lacks (an explicit critical point coordinate) from its steering direction, we simulated it with the robot model, gains, and formation of Section~\ref{sec:ch6:sim_setup} in three forms (Table~\ref{tab:ch6:baselines}): the primitive alone; its tangential steering paired with the laws (\ref{eq:ch6:10})--(\ref{eq:ch6:11}) given an oracle (the true critical point) as target and center; and the same laws fed a coordinate recovered from magnitudes alone, $\hat{\mathbf{p}}_m = \mathbf{p}_c - S_c\mathbf{g}/\|\mathbf{g}\|^2$, where $S_c$ is the fitted centroid speed.

On the six analytical fields the speed grows equally in every direction, so $\mathbf{g}$ points almost directly away from $\mathbf{p}^*$ and the primitive reached $\mathbf{p}^*$ in every trial. Its command is the fitted gradient, which near $\mathbf{p}^*$ shrinks with distance at a gain of up to $1/R$ (six times $k$). The primitive settles faster and closer to $\mathbf{p}^*$, because friction stops a proportional command on the ring $v_c/k$, but at 1.6 times the control effort. In orbit, the primitive alone has no mechanism to detect or correct radial error with respect to a specific point, so every orbit drifted outward. Given the oracle, the same tangent with the radial term of (\ref{eq:ch6:11}) matched the Jacobian method within 0.004~m RMSE at every radius. With its own speed estimate the primitive orbited inside the commanded radius, because the estimate's inward bias outweighs the lag drift of Section~\ref{sec:ch6:stability}; in attraction it left the cluster still moving after 10~s.

The gradient $\mathbf{g}$ points away from $\mathbf{p}^*$ only where the speed grows equally in every direction. The anisotropic node $\mathbf{J} = \operatorname{diag}(-1, -0.2)$ is affine, but its speed grows five times more slowly along one axis. There the primitive reached $\mathbf{p}^*$ in 6\% of trials, and most of the rest ended on the slow axis. Even with the oracle 88\% of orbits at $r_d = 0.1$~m were lost, because the speed contours are ellipses; the Jacobian estimate is exact there, as on any affine field. In the gyre the speed vanishes at the center and the saddles alike, so the primitive cannot tell them apart, while the Jacobian method labeled the type on every cycle (Section~\ref{sec:ch6:sim_class}). The primitive lost every gyre orbit, and with the oracle 92\% at $r_d = 0.3$~m and all at 0.5~m.

% Ported from: cwaight_systems_paper_with_revisions2.tex
\begin{table}[htbp]
\centering
\scriptsize
\caption{Comparison with the vector-to-scalar (V2S) primitive in simulation, 500 attraction trials of 10~s and 50 orbit runs of 60~s per setting. Accuracy uses successful trials, gyre success counts the center, and a dash marks mostly lost orbits.}
\label{tab:ch6:baselines}
\setlength{\tabcolsep}{1.8pt}
\begin{tabular}{ll|cccc|ccc}
\hline
 & & \multicolumn{4}{c|}{\textbf{Attraction}} & \multicolumn{3}{c}{\textbf{Orbit RMSE (m)}} \\
\textbf{Field} & \textbf{Navigation} & \textbf{Succ.} & \textbf{$\bar{e}_f$} & \textbf{$T_{\mathrm{conv}}$} & \textbf{$E_u$} & \multicolumn{3}{c}{\textbf{at $r_d$ (m)}} \\
 & & \textbf{(\%)} & \textbf{(m)} & \textbf{(s)} & \textbf{(m$^2$/s)} & 0.1 & 0.3 & 0.5 \\
\hline
Linear (6) & Jacobian & 100 & 0.023 & 2.3 & 0.067 & 0.018 & 0.007 & 0.004 \\
 & V2S & 100 & 0.005 & 1.3 & 0.106 & 0.252 & 0.158 & 0.111 \\
 & V2S + oracle & 100 & 0.023 & 2.3 & 0.067 & 0.018 & 0.010 & 0.008 \\
 & V2S + speed est. & 100 & 0.017 & 1.4 & 0.568 & 0.075 & 0.052 & 0.035 \\
\hline
Aniso.\ node & Jacobian & 100 & 0.023 & 2.3 & 0.067 & 0.018 & 0.007 & 0.004 \\
 & V2S & 6 & 0.057 & 5.6 & 0.082 & -- & -- & -- \\
 & V2S + oracle & 100 & 0.023 & 2.3 & 0.067 & -- & 0.125 & 0.117 \\
 & V2S + speed est. & 6 & 0.067 & 9.9 & 0.376 & -- & -- & -- \\
\hline
Gyre & Jacobian & 50 & 0.024 & 1.8 & 0.041 & 0.018 & 0.066 & -- \\
 & V2S & 49 & 0.005 & 1.0 & 0.073 & -- & -- & -- \\
 & V2S + oracle & 100 & 0.023 & 2.3 & 0.067 & 0.019 & -- & -- \\
 & V2S + speed est. & 49 & 0.017 & 1.1 & 0.579 & 0.075 & 0.273 & -- \\
\hline
\end{tabular}
\end{table}

\section{Summary}
\label{sec:ch6:summary}

% Ported from: cwaight_systems_paper_with_revisions2.tex
This chapter presented a framework for multirobot navigation in vector fields containing critical points. A Jacobian-based approach used instantaneous measurements from three or more robots to identify critical point locations without prior field knowledge, supporting attraction and orbital control laws built on these estimates. Simulation across six canonical field types confirmed convergence of these control primitives. On curved fields the sign of $\det\hat{\mathbf{J}}$ on the first cycle decided which type of critical point the cluster reached. In 459 hardware runs, the cluster converged within 0.05~m in 99\% of saddle and 91\% of vortex trials with three robots and in every vortex trial with four, with mean position errors of 0.034~m or less. Orbital control experiments maintained commanded radii within a quarter of the radius from 0.20~m upward with three robots in both fields, and below that held bounded orbits that did not track the commanded radius; at 0.10~m the vortex orbit settled onto a limit cycle sized by the robots' delay and lag. Simulations matched to each run reproduced the saddle's attraction error and about half of the orbit fluctuation.
Together these results establish cooperative Jacobian estimation as a practical kinematic foundation for multirobot critical point navigation. The printed fields do not advect the robots, so compensating for a physical flow that carries the platform is the natural next step toward field deployment.

% New text
Offline, the hardware attraction estimates also named the type of both printed
fields in every trial by majority vote, once each map's sensor frame was
registered to the motion-capture frame.

% ===================================================================
% CHAPTER 7: SECOND-ORDER PRIMITIVES (Separatrix / OW paper, Draft 14)
% ===================================================================
\chapter{Second-Order Primitives: Coherent Structures}
\label{ch:second_order}

\section{Introduction}
\label{sec:ch7:intro}

% New text
Chapter~\ref{ch:first_order} located isolated critical points from a
first-order fit. This chapter moves to the second rung of the ladder of
Section~\ref{sec:ch4:ladder}, where the features are curves.

% Ported from: Draft_14.tex
A vector field also carries well-defined curves. In a dynamic ocean flow some of these are coherent structures, boundaries that influence where floating material such as debris and larvae accumulates or separates \cite{56,50}.

% Ported from: Draft_14.tex
Coherent structures are typically identified in one of two ways. The
first finds Lagrangian coherent structures (LCS) \cite{57}, typically as
ridges of the finite-time Lyapunov exponent (FTLE) field \cite{52}. The
second uses instantaneous Eulerian criteria, such as the Okubo-Weiss
parameter, which partition the flow by the local velocity gradient alone
\cite{4,48}. This instantaneous approach has been generalized to three
dimensions \cite{58}, refined to second order \cite{59}, and applied at
scale in mesoscale eddy censuses \cite{11}. Their tradeoffs remain under
study \cite{51,49,60,61}.

% Ported from: Draft_14.tex
Mapping these curves tells a search where floating material has gathered
\cite{50}, and ocean sampling campaigns already steer vehicles toward
dynamic features such as fronts \cite{23,62,22}. Without adaptive
navigation, a large fleet must survey the basin and reconstruct the
structures offline; adaptive navigation maps them locally with one small
cluster.

% Ported from: Draft_14.tex
Michini et al. travel along a presumed stable or
unstable manifold with three robots using only local velocity
measurements \cite{27}, an approach based on the PIM triple procedure
that its authors extend to larger teams \cite{63} and to robots in a flow
tank \cite{28}. The saddle point is addressed by neither the
controller nor its analysis.
Kularatne and Hsieh confirm an attracting boundary by computing the local
FTLE field on board, at the cost of a nonzero integration horizon
\cite{6}.
Their analysis assumes that the robots start straddling the structure
\cite{27,63} or that the center robot starts on it \cite{6}.

% Ported from: Draft_14.tex
This chapter shows that a second-order polynomial fit of the local flow can
acquire and track such a structure from instantaneous measurements. The
chapter contributes as follows.
First, a cooperative estimator recovers the local quadratic model of
a two-component field from one synchronized set of measurements across six
robots, the fewest that can, extending the scalar gradient and Hessian
estimation of \cite{15} to both components of a vector field. Second, two surrogate fields computed from that fit,
the determinant field $D$ and the smaller strain eigenvalue field $s_1$,
turn the search for a separatrix into the search for a trench, the kind of
feature the scalar primitives of \cite{3} follow. Third, two AN
primitives, the $D$ tracker and the $s_1$ tracker, navigate to the trench
from a start that does not straddle it and ride along it. Finally, a
comparison shows that the $D$ tracker tolerates more measurement noise,
while the $s_1$ tracker is built from objective quantities, so that
observers rotating relative to each other agree on its path except where
its tangent swaps eigenvector. A rotation-to-strain ratio from the same
fit indicates when either trench likely marks a separatrix.

% Ported from: Draft_14.tex
Both primitives run on a simulated double gyre, three variants, and
measured Santa Barbara Channel radar currents.

\section{Six-Robot Quadratic Estimation}
\label{sec:ch7:estimation}

% Ported from: Draft_14.tex
The critical point estimator of Chapter~\ref{ch:first_order} fits a plane to each
component of the vector field and solves for the point where $\mathbf{v} = \mathbf{0}$. 

This chapter extends the linear approximation to the second-order Taylor expansion about the
centroid $\mathbf{p}_c$, with $x$ and $y$ measured from $\mathbf{p}_c$.

\begin{align}
    u(\mathbf{p}) &= a_1 + a_2 x + a_3 y
        \nonumber \\
        &\quad + a_4 xy + \tfrac{1}{2} a_5 x^2 + \tfrac{1}{2} a_6 y^2
        + \mathcal{O}(\|\mathbf{p}\|^3),
        \label{eq:ch7:quad_model_u} \\
    v(\mathbf{p}) &= b_1 + b_2 x + b_3 y
        \nonumber \\
        &\quad + b_4 xy + \tfrac{1}{2} b_5 x^2 + \tfrac{1}{2} b_6 y^2
        + \mathcal{O}(\|\mathbf{p}\|^3),
        \label{eq:ch7:quad_model_v}
\end{align}
where $\mathcal{O}(\|\mathbf{p}\|^3)$ is the truncation error.
The twelve coefficients are constants, and the one-half factors make each
fitted coefficient an estimate of a derivative of the flow at $\mathbf{p}_c$. The flow itself is $a_1$,
$a_2$ and $a_3$ are $\partial u/\partial x$ and $\partial u/\partial y$,
and $a_4$, $a_5$, and $a_6$ are $\partial^2 u/\partial x\partial y$,
$\partial^2 u/\partial x^2$, and $\partial^2 u/\partial y^2$, with
$b_1, \ldots, b_6$ the same for $v$.

To estimate these coefficients with robot positions $\mathbf{p}_i = (x_i,
y_i)$, we define the quadratic basis vector
\begin{equation}
    \boldsymbol{\phi}(\mathbf{p}) =
    \bigl[\,1, \; x, \; y, \; xy, \;
    \tfrac{1}{2}x^2, \; \tfrac{1}{2}y^2\,\bigr]^\top.
\end{equation}



Stacking the $n$ basis vectors row-wise gives us the $n \times 6$ formation matrix $\boldsymbol{\Phi}$, and the coefficient estimates $\hat{\boldsymbol{\theta}}_u$ and $\hat{\boldsymbol{\theta}}_v$ solve
\begin{equation}
\boldsymbol{\Phi}\,\hat{\boldsymbol{\theta}}_u = \mathbf{m}_u, \qquad
\boldsymbol{\Phi}\,\hat{\boldsymbol{\theta}}_v = \mathbf{m}_v,
\label{eq:ch7:solve}
\end{equation}
where $\mathbf{m}_u = [u_1, \ldots, u_n]^\top$ and
$\mathbf{m}_v = [v_1, \ldots, v_n]^\top$ collect the components $u_i$
and $v_i$ that robot $i$ measures.

\subsection{Minimality}
\label{sec:ch7:minimality}

% Ported from: Draft_14.tex
Six robots is the minimum for recovery of the unconstrained local
quadratic model from one set of instantaneous measurements, the
second-order counterpart of the three-robot minimum for a linear fit
(Proposition~\ref{prop:min}). Each robot contributes one equation per component of
(\ref{eq:ch7:solve}), and each system has six unknowns. With five robots or
fewer, a method must either return a family of solutions or impose a
prior such as $\nabla\cdot\mathbf{v} = 0$, which would bias every
coefficient of a horizontal surface current, since such currents are not
divergence-free. Six robots suffice provided $\boldsymbol{\Phi}$ is
nonsingular, and more than six need $\boldsymbol{\Phi}$ to have full
column rank for least squares. Either condition fails exactly when the
robots lie on a common conic, so a ring formation is degenerate at any
size. A fit from more than six robots also leaves a residual, which
carries both the measurement noise and the truncation error.

We use the pentagon-plus-center formation of \cite{15} throughout. It is
the smallest uniform ring plus center on which \cite{15} recovers a scalar
Hessian exactly for quadratic signals at any orientation, and the same
symmetry makes the noise gains of the fit here independent of heading.

\subsection{Estimator Sensitivity}
\label{sec:ch7:sensitivity}

% Ported from: Draft_14.tex
Three things limit the fit: the formation, the noise on each reading,
and the error in each robot's position.

\subsubsection{Formation}
The 2-norm condition number $\kappa(\boldsymbol{\Phi})$ indicates how much
the fit amplifies noise. On the radius-normalized $\boldsymbol{\Phi}$, with
each column of order $q$ divided by $\rho^{\,q}$ for formation radius
$\rho$, the pentagon-plus-center formation gives
$\kappa(\boldsymbol{\Phi}) = 8.26$ at every radius. The formation's size
changes instead how much noise reaches each coefficient. A coefficient of order $q$ is recovered
from differences in readings across the footprint divided by
$\rho^{\,q}$, so its noise scales as $\rho^{-q}$. A larger formation
therefore suppresses noise most in the second-order terms, but it also
increases the truncation error, since the ring robots sit at
$\|\mathbf{p}\| = \rho$. Moving the center robot out to $0.99\rho$
places all six robots within $0.6\%$ of a common circle and raises
$\kappa(\boldsymbol{\Phi})$ to $564$, so the fit degrades smoothly near
a conic.

\subsubsection{Measurement Noise}
Every robot reads the current with its own sensor, and each fit uses only
the readings of its own cycle. The noise in one fit is therefore
independent across robots and across cycles. It is modeled as additive on
each reading, $u_i^m = u(\mathbf{p}_i) + \eta_i^u$ and
$v_i^m = v(\mathbf{p}_i) + \eta_i^v$ with
$\eta \sim \mathcal{N}(0, \sigma_{uv}^2)$. The fit is linear in the
readings, so the noise reaches the coefficients as
$\boldsymbol{\Phi}^{-1}\boldsymbol{\eta}$, where $\boldsymbol{\eta}$
stacks the noise of the six robots. It adds no bias, and its covariance is
$\mathrm{Cov}(\hat{\boldsymbol{\theta}}) =
\sigma_{uv}^2(\boldsymbol{\Phi}^\top\boldsymbol{\Phi})^{-1}$ for each of
$\hat{\boldsymbol{\theta}}_u$ and $\hat{\boldsymbol{\theta}}_v$, which
are independent. For the pentagon-plus-center formation its diagonal does
not depend on heading. A reading noise $\sigma_{uv}$ gives the flow estimate a standard deviation
$\sigma_{uv}$, each first derivative $0.632\,\sigma_{uv}/\rho$, the
mixed second derivative $1.265\,\sigma_{uv}/\rho^2$, and each pure
second derivative $2.53\,\sigma_{uv}/\rho^2$.

\subsubsection{Position Noise}
A robot with position error
$\boldsymbol{\xi}_i \sim \mathcal{N}(\mathbf{0}, \sigma_p^2\mathbf{I})$
samples the field at $\mathbf{p}_i + \boldsymbol{\xi}_i$, but the fit is
built from its nominal position $\mathbf{p}_i$. The formation matrix is
therefore wrong by $\Delta\boldsymbol{\Phi}$, whose row $i$ is
$\boldsymbol{\phi}(\mathbf{p}_i + \boldsymbol{\xi}_i)^{\top} -
\boldsymbol{\phi}(\mathbf{p}_i)^{\top}$. Under the quadratic model the
measurements are
$\mathbf{m}_u = (\boldsymbol{\Phi} + \Delta\boldsymbol{\Phi})
\boldsymbol{\theta}_u$, and the fit returns
\begin{equation}
    \hat{\boldsymbol{\theta}}_u = \boldsymbol{\theta}_u
        + \boldsymbol{\Phi}^{-1}\Delta\boldsymbol{\Phi}\,
          \boldsymbol{\theta}_u ,
    \label{eq:ch7:pos_noise}
\end{equation}
and likewise for $\hat{\boldsymbol{\theta}}_v$. To first order, entry
$i$ of $\Delta\boldsymbol{\Phi}\,\boldsymbol{\theta}_u$ is
$\nabla u \cdot \boldsymbol{\xi}_i$, the change in $u$ across the
displacement. This error vanishes in uniform flow and grows with the
local gradient, so unlike measurement noise it depends on the field.

\section{Eulerian Surrogates from the Fit}
\label{sec:ch7:surrogates}

% Ported from: Draft_14.tex
Lagrangian methods locate coherent structures by integrating particle
trajectories over a finite time window \cite{57,52}. This chapter instead
reads them from the instantaneous geometry of the velocity field, as
Eulerian coherent structures \cite{49}. Two fields computed from the
six-robot fit can mark the separatrix. The first is the determinant of the velocity
gradient, and the second is the smaller eigenvalue of its symmetric
part.

\subsection{Determinant Field}
\label{sec:ch7:det_field}

% Ported from: Draft_14.tex
The fitted coefficients define the velocity gradient
$\mathbf{J} = \nabla\mathbf{v}$ at every point of the footprint. Differentiating the quadratic model gives a Jacobian whose
entries are affine in position,
\begin{equation}
    \hat{\mathbf{J}}(\mathbf{p}) =
    \begin{bmatrix}
        \hat{a}_2 + \hat{a}_5 x + \hat{a}_4 y &
        \hat{a}_3 + \hat{a}_4 x + \hat{a}_6 y \\
        \hat{b}_2 + \hat{b}_5 x + \hat{b}_4 y &
        \hat{b}_3 + \hat{b}_4 x + \hat{b}_6 y
    \end{bmatrix}.
    \label{eq:ch7:Jhat}
\end{equation}
The determinant $\hat{D}(\mathbf{p}) = \det\hat{\mathbf{J}}(\mathbf{p})$
is therefore exactly quadratic over the footprint. Its value, gradient,
and Hessian at the centroid are polynomials in the coefficients. The
value is $\hat{D}_0 = \hat{a}_2\hat{b}_3 - \hat{a}_3\hat{b}_2$, and the
gradient is
\begin{equation}
    \hat{\mathbf{g}}_0 =
    \begin{bmatrix}
        \hat{a}_5\hat{b}_3 + \hat{a}_2\hat{b}_4 - \hat{a}_4\hat{b}_2 - \hat{a}_3\hat{b}_5 \\
        \hat{a}_4\hat{b}_3 + \hat{a}_2\hat{b}_6 - \hat{a}_6\hat{b}_2 - \hat{a}_3\hat{b}_4
    \end{bmatrix}.
    \label{eq:ch7:grad_det}
\end{equation}
The Hessian is
\begin{equation}
    \hat{\mathbf{H}}_{D,0} =
    \begin{bmatrix}
        2(\hat{a}_5\hat{b}_4 - \hat{a}_4\hat{b}_5) &
        \hat{a}_5\hat{b}_6 - \hat{a}_6\hat{b}_5 \\
        \hat{a}_5\hat{b}_6 - \hat{a}_6\hat{b}_5 &
        2(\hat{a}_4\hat{b}_6 - \hat{a}_6\hat{b}_4)
    \end{bmatrix},
    \label{eq:ch7:hess_det}
\end{equation}
which is constant over the footprint.

% New text
Section~\ref{sec:ch3:ow} gives the physical reading of $D$, which
separates strain-dominated from rotation-dominated flow.

\subsection{Strain Eigenvalue Field}
\label{sec:ch7:s1_field}

% Ported from: Draft_14.tex
The second surrogate comes from the rate-of-strain tensor
$\hat{\mathbf{S}} = \tfrac{1}{2}(\hat{\mathbf{J}} + \hat{\mathbf{J}}^{\top})$,
the symmetric part of the same fitted Jacobian.
At the centroid its mean, deviatoric normal, and shear strain are
$\hat{\mu} = (\hat{a}_2 + \hat{b}_3)/2$,
$\hat{s}_n = (\hat{a}_2 - \hat{b}_3)/2$, and
$\hat{s}_s = (\hat{a}_3 + \hat{b}_2)/2$, and each is affine in position.
The eigenvalues of $\hat{\mathbf{S}}$ are
$\hat{s}_{1,2} = \hat{\mu} \mp \hat{r}$ with
$\hat{r} = \sqrt{\hat{s}_n^2 + \hat{s}_s^2}$. Their eigenvectors
$\mathbf{e}_1$ and $\mathbf{e}_2$ are the compression and stretching
directions, and the gradient of the smaller eigenvalue is
\begin{equation}
    \nabla\hat{s}_1 = \nabla\hat{\mu}
        - \frac{\hat{s}_n\nabla\hat{s}_n + \hat{s}_s\nabla\hat{s}_s}{\hat{r}},
    \label{eq:ch7:grad_s1}
\end{equation}
where the second-order coefficients give the constant gradients
\begin{equation}
    \begin{aligned}
        \nabla\hat{\mu} &= \tfrac{1}{2}(\hat{a}_5 + \hat{b}_4,\; \hat{a}_4 + \hat{b}_6), \\
        \nabla\hat{s}_n &= \tfrac{1}{2}(\hat{a}_5 - \hat{b}_4,\; \hat{a}_4 - \hat{b}_6), \\
        \nabla\hat{s}_s &= \tfrac{1}{2}(\hat{a}_4 + \hat{b}_5,\; \hat{a}_6 + \hat{b}_4).
    \end{aligned}
    \label{eq:ch7:grad_strain}
\end{equation}
The quotient stays bounded as $\hat{r} \to 0$ but its direction is
undefined at $\hat{r} = 0$, so the computation floors $\hat{r}$ at
$10^{-12}$. Unlike $\hat{D}$, $\hat{s}_1$ has no usable Hessian. It is an
affine function minus the norm of an affine map, so it is concave over the
footprint for every field and formation.

% New text
Section~\ref{sec:ch3:oecs} defines the attracting and repelling OECS
that the trenches of $s_1$ mark, and Section~\ref{sec:ch3:relation}
relates the two fields, including how each changes for a rotating
observer.

\subsection{Fit Check on the Double Gyre}
\label{sec:ch7:fit_check}

% New text. The closed forms of the double gyre are in Section~\ref{sec:ch3:dg_example};
% these two paragraphs of Draft 14 compare the fit against them.

% Ported from: Draft_14.tex
These closed forms also check the fit. Along the separatrix at
$\rho = 0.075$ and without noise, the fitted $\hat{\mathbf{J}}$, $\hat{D}$,
$\nabla\hat{D}$, and $\nabla\hat{s}_1$ stay within 1.4,
2.8, 8.2, and 4.0\% of the largest true value of each,
errors that scale as $\rho^2$ for the first two and $\rho$ for the
gradients. Under measurement noise, the spread of each fitted pure second
derivative matches $2.53\,\sigma_{uv}/\rho^2$ to within 1\%.

The fitted Hessian $\hat{\mathbf{H}}_{D,0}$ is the exception, for a
structural reason. As the Hessian of the determinant of the fitted
Jacobian, it keeps products of second derivatives but drops products of
first and third derivatives, which a quadratic model does not carry. Near
the saddles, where the velocity gradient is largest, the dropped terms
dominate: at $y = -0.45$ the fitted curvature across the trench is
0.43, against the true $2\pi^6 A^2 = 19.2$. Its eigenvectors
survive: for $|y| \leq 0.40$ they lie within 3.5$^\circ$ of the true
trench directions, and the across-trench eigenvalue stays positive up to
the saddles.

\section{Six-Robot Control Architecture}
\label{sec:ch7:architecture}

% New text
The cluster runs on the three-layer architecture of
Section~\ref{sec:ch4:architecture}, with the pentagon-plus-center
formation of Section~\ref{sec:ch4:pentagon}.

% Ported from: Draft_14.tex
This chapter fills two blocks of the first layer:
feature estimation runs the quadratic fit of Section~\ref{sec:ch7:estimation}, and the
navigation controller runs the $D$ or $s_1$ tracker, which emits a desired
centroid velocity $\dot{\mathbf{p}}_c$. A formation controller adds
proportional shape regulation (gain 1.0 on lengths, 0.1 on angles), and
the heading holds its initial value.

% Ported from: Draft_14.tex
In the robot controller layer, each robot follows its velocity command
through the model of Section~\ref{sec:ch2:dynamics}, identified on the Decabot indoor
omnidirectional testbed \cite{42,30}. A robot reaches a share
$g = 0.86$ of its command after a delay $L = 0.055$~s, through
a first-order lag with time constant $\tau = 0.02$~s. It ignores
commands below a dead zone $v_c = 0.025$~m/s and is capped at 0.3~m/s. A
Python simulation integrates each robot's position by explicit Euler at the
control period $\Delta t = 0.1$~s. The simulation code and the scripts
behind every figure are available from the authors on request.

\section{The $D$ Tracker}
\label{sec:ch7:sep_controller}

% Ported from: Draft_14.tex
The $D$ tracker is built for a mountain pass in $D$. Two minima are
joined by a trench that rises to a saddle point of $D$, the crest. On the
benchmark the trench is the separatrix and the minima are the flow saddles
(Fig.~\ref{fig:ch3:dg_fields}(b)). The tracker steers onto the trench, over
the crest, and down to the next minimum, with one channel across the trench,
one along it.

Across a trench $D$ curves upward, and along it $D$ can curve either way.
The tracker reads both directions from the eigenvectors of the fitted
Hessian, $\hat{\mathbf{H}}_{D,0}\mathbf{w}_j = \lambda_j \mathbf{w}_j$,
with $\|\mathbf{w}_j\| = 1$ and $\lambda_1 \leq \lambda_2$. The
eigenvector $\mathbf{w}_2$ points across the trench and $\mathbf{w}_1$
along it, signed by $\hat{\mathbf{v}}_0^\top \mathbf{w}_1 \geq 0$ to point
downstream, where $\hat{\mathbf{v}}_0 = (\hat{a}_1, \hat{b}_1)$ is the
measured flow at the centroid. This follows the height-ridge definition
of a trench \cite{64}, and on the benchmark the fitted eigenvectors lie
close to the true trench directions (Section~\ref{sec:ch7:fit_check}).

Away from the crest, both channels take Newton steps of the fitted
quadratic, $-\mathbf{w}_2^\top\hat{\mathbf{g}}_0/\lambda_2$ toward the
trench bottom and $|\mathbf{w}_1^\top\hat{\mathbf{g}}_0|/|\lambda_1|$
along it, which slows the cluster near the crest.

At the crest the along-trench gradient vanishes, and the measured flow
carries the cluster over instead. The tracker detects the crest as a
fitted saddle of $D$ whose stationary point lies inside the formation, a
crest set $\mathcal{B}$ where $\lambda_1\lambda_2 < 0$ and
\mbox{$\|\hat{\mathbf{H}}_{D,0}^{-1}\hat{\mathbf{g}}_0\| < c_{\mathrm{crest}}\,\rho$}.
The left side is the Newton step to the stationary point, so
$c_{\mathrm{crest}}$ is a fraction of the formation radius and needs no
knowledge of the field's scale.

Combining the two channels, the crest test selects the along-trench term,
\begin{equation}
    v_\parallel =
    \begin{cases}
        T\,\hat{\mathbf{v}}_0^\top \mathbf{w}_1,
            & \mathbf{p}_c \in \mathcal{B}, \\[6pt]
        \dfrac{|\mathbf{w}_1^\top \hat{\mathbf{g}}_0|}{|\lambda_1|},
            & \text{otherwise},
    \end{cases}
    \label{eq:ch7:vpar}
\end{equation}
and the centroid velocity command is
\begin{equation}
    \dot{\mathbf{p}}_c = k\Bigl[
    \underbrace{\mathrm{sat}(v_\parallel)\,\mathbf{w}_1}_{\text{along the trench}}
    + \underbrace{\mathrm{sat}\!\left(-\frac{\mathbf{w}_2^\top
      \hat{\mathbf{g}}_0}{\lambda_2}\right)\mathbf{w}_2}_{\text{across the trench}}\Bigr].
    \label{eq:ch7:d_law}
\end{equation}
Every term in the brackets is a distance, with $T$ one time unit in the
crest branch, so the navigation gain $k$ has units of inverse time. The
function $\mathrm{sat}(c) = c_{\max}\tanh(c/c_{\max})$ saturates each
channel at the distance $c_{\max}$, which caps the centroid speed at
$\sqrt{2}\,k\,c_{\max}$. It also bounds each Newton step where its
eigenvalue nears zero, and each eigenvalue in a denominator is floored at
$10^{-9}$ in magnitude, so no step divides by zero. Both branches of
$v_\parallel$ are nonnegative, so the along-trench motion never opposes
the flow, and where the fitted Hessian is definite the same law runs
unchanged.

The tracker stops at the end of the pass with the critical point test of
Chapter~\ref{ch:first_order}, which fires when the fitted flow has a saddle inside the
formation,
$\det\hat{\mathbf{J}}_0 < 0$ and
$\|\hat{\mathbf{J}}_0^{-1}\hat{\mathbf{v}}_0\| < \rho$, where $\hat{\mathbf{J}}_0$ is the fitted Jacobian at the centroid and
$\hat{\mathbf{J}}_0^{-1}\hat{\mathbf{v}}_0$ is the Newton step to the
stagnation point. The test uses only the zeroth- and first-order
coefficients, the least noisy the fit returns, and certifies a saddle of
the fitted linear flow, which on the benchmark is the flow saddle at the
end of the trench. Once it passes, the tracker
holds for the rest of the run with the attraction law of Chapter~\ref{ch:first_order},
$\dot{\mathbf{p}}_c = k\,\mathrm{sat}(-\hat{\mathbf{J}}_0^{-1}\hat{\mathbf{v}}_0)$
with $\mathrm{sat}$ applied to each component. Both read the measured
flow, so neither is objective.

\section{The $s_1$ Tracker}
\label{sec:ch7:s1_controller}

% Ported from: Draft_14.tex
The fitted $\hat{s}_1$ is concave for any field and formation, so it
never shows the upward curvature a
trench has across it, and the $D$ tracker's Newton step has no counterpart
here. Both channels instead use $\nabla\hat{s}_1$ and the strain
eigenframe.

The trench tangent is a strain eigenvector (Section~\ref{sec:ch3:oecs}),
and the tracker must choose between $\mathbf{e}_1$ and $\mathbf{e}_2$ at
every step. On the trench the transverse derivative of $s_1$ vanishes, so
$\nabla\hat{s}_1$ lies nearly along the trench away from a stationary
point. The tangent is the eigenvector most nearly parallel to it, signed
against a reference,
\begin{equation}
    \mathbf{t}_{\mathrm{raw}} =
    \arg\max_{\mathbf{e} \in \{\mathbf{e}_1, \mathbf{e}_2\}}
    \bigl\lvert \nabla\hat{s}_1^\top \mathbf{e} \bigr\rvert,
    \qquad
    \mathbf{t} = \mathrm{sgn}\bigl(\mathbf{t}_{\mathrm{ref}}^\top
    \mathbf{t}_{\mathrm{raw}}\bigr)\,\mathbf{t}_{\mathrm{raw}},
    \label{eq:ch7:tangent_select}
\end{equation}
where $\mathbf{t}_{\mathrm{ref}}$ is the previous step's tangent
(Fig.~\ref{fig:ch7:s1_channels}(a)). The first ride step has no previous
tangent, so the measured flow $\hat{\mathbf{v}}_0$ serves as
$\mathbf{t}_{\mathrm{ref}}$, the one non-objective input. After that,
objectivity holds at three levels. The quantities the tracker reads,
$\nabla\hat{s}_1$ and the strain eigenframe, are objective. The command
is rotationally equivariant only where the sign test is decisive: an
observer rotating at rate $\Omega$ turns the stored tangent by
$\Omega\,\Delta t$ between steps, which flips the sign only where
\mbox{$|\mathbf{t}_{\mathrm{ref}}^\top\mathbf{t}_{\mathrm{raw}}| <
\sin\Omega\,\Delta t$}, as at an eigenframe swap. The path, which a
rotating observer also sees carried by the frame's transport velocity,
matches across frames only empirically.

From the first step, the command pairs a constant-speed ride along
$\mathbf{t}$ with gradient descent across the trench,
\begin{equation}
    \dot{\mathbf{p}}_c = k\Bigl[
    \underbrace{\mathrm{sat}(c_{\max})\,\mathbf{t}}_{\text{along the trench}}
    + \underbrace{\mathrm{sat}\bigl(-g_\perp\,\mathbf{t}_\perp^\top
      \nabla\hat{s}_1\bigr)\,\mathbf{t}_\perp}_{\text{across the trench}}\Bigr],
    \label{eq:ch7:s1_law}
\end{equation}
with transverse gain $g_\perp$ and $\mathbf{t}_\perp$ the unit normal to
$\mathbf{t}$. As in (\ref{eq:ch7:d_law}) each bracketed term is a distance, so
$g_\perp$ carries units of length squared times time. The cluster rides at
speed $k\,c_{\max}\tanh 1$ throughout, and descends only across the
trench (Fig.~\ref{fig:ch7:s1_channels}(b)).

The ride speed does not depend on the slope, so the cluster crosses the
crest without the crest test the $D$ tracker needs. In incompressible flow
$s_1 = -r \leq 0$ reaches zero only at an isotropic point, where the
strain vanishes. In general such a point is isolated, so a crest is a
smooth saddle where $\nabla\hat{s}_1$ vanishes and the tangent rule is
least reliable. On the benchmark the shear strain vanishes identically, the
isotropic points form curves, and the crest sits on one at a kink of
$s_1$. There the tangent swaps eigenvector, and the rule follows the swap
with no dedicated test, since $\nabla\hat{s}_1$ stays large on either
side. Above the origin $u_x < 0$ and the trench follows $\mathbf{e}_2$,
an attracting OECS; below, $u_x > 0$ gives $\mathbf{e}_1$, a repelling
OECS.

The tracker stops at a minimum by a gradient test. The concave fit cannot
certify a minimum, but the fitted gradient tracks the true gradient even
where the fitted curvature has the wrong sign. Each end of the trench is
a true two-dimensional minimum of $s_1$, so a gradient that is small
against the depth suffices. Capture requires $\hat{s}_1 < 0$ and
$q = \rho\,\|\nabla\hat{s}_1\|/|\hat{s}_1| < q_{\mathrm{cap}}$, where $q$ is the fraction by which $\hat{s}_1$ changes across the
footprint. The ratio is dimensionless, so $q_{\mathrm{cap}}$ needs no
knowledge of the field's strain rate. It is large along the ride and wherever $\hat{s}_1$ nears zero, as at an isotropic point or a
gyre center, but small at any stationary point with $\hat{s}_1 < 0$, so
the test can also pass at a smooth crest. On the benchmark the crest sits
on an isotropic point, so only the minima pass. Once the test passes, the
tracker
drops the ride term and descends $\hat{s}_1$ for the rest of the run with
$\dot{\mathbf{p}}_c = -k\,\mathrm{sat}_{\mathrm{v}}(g_\perp\nabla\hat{s}_1)$,
where $\mathrm{sat}_{\mathrm{v}}(\mathbf{a}) =
\mathrm{sat}(\|\mathbf{a}\|)\,\mathbf{a}/\|\mathbf{a}\|$, with
$\mathrm{sat}_{\mathrm{v}}(\mathbf{0}) = \mathbf{0}$, saturates the
vector's magnitude and preserves its direction, which keeps the descent
independent of the coordinate axes.

\begin{figure}[htbp]
    \centering
    \includegraphics[width=0.95\textwidth]{ch7_s1_channels.png}
    \caption{The two channels of the $s_1$ tracker on the benchmark trench
    $x = 0$. (a) Tangent selection: the strain eigenframe, selected
    eigenvector heavy, with $\nabla\hat{s}_1$ in orange. The tangent is
    $\mathbf{e}_2$ on the attracting segment and $\mathbf{e}_1$ on the
    repelling one, swapping at the origin. (b) The command: an open-loop
    ride along $\mathbf{t}$ and a descent across it.}
    \label{fig:ch7:s1_channels}
\end{figure}

\section{Closed-Loop Behavior}
\label{sec:ch7:closed_loop}

% Ported from: Draft_14.tex
Both trackers pull the cluster onto the trench through one saturated loop
across it, of the form analyzed in Section~\ref{sec:ch6:stability} on the same robots. With the
delay and gain of Section~\ref{sec:ch7:architecture}, the sampled loop is
exponentially stable for loop gains below $1/(gL) = 21.1$ per unit
time, or 16.1 with the lag included, and as a stable linear loop it
keeps the across-trench error bounded under a bounded estimate error
(Section~\ref{sec:ch6:stability}). The result is local and linearized, not a guarantee for the
full switched, saturated controller.

Neither channel is exactly linear. The $D$ tracker divides by the fitted
curvature, so its gain is $k$ times the ratio of true to fitted curvature
across the trench, $1/\sin^2\pi y_f$ on the benchmark. The gain is
$k = 3$ at the crest and exceeds the bound only within 0.14 of the
saddles (Section~\ref{sec:ch7:fit_check}). The $s_1$ tracker's gain is
$k\,g_\perp\,\partial^2 s_1/\partial n^2$, or $3\pi^4 A\,|\cos\pi y_f|$
on the benchmark, which exceeds the bound for $|y| > 0.19$ and
reaches 29 at the saddles. Each channel saturates a few thousandths
from the trench, which keeps any resulting oscillation in a narrow band.
At the end of the ride, the $D$ tracker's hold is the
attraction law of Chapter~\ref{ch:first_order}, so the cluster settles inside the friction
ring $v_c/k = 0.008$ around the fitted stagnation point.
Table~\ref{tab:ch7:params} collects the operating points.

% Ported from: Draft_14.tex
\begin{table}[htbp]
\caption{Operating points. The double-gyre column is non-dimensional,
with one time unit equal to one second of the Decabot's identified
dynamics. The ocean column uses the scaling of the Santa Barbara Channel
trial.}
\label{tab:ch7:params}
\centering
\begin{tabular}{llll}
\hline
Parameter & Double gyre & Ocean & Units \\
\hline
$k$                     & 3.0   & 1.94   & 1/time \\
$c_{\max}$              & 0.04  & 0.033  & length \\
$\rho$                  & 0.075 & 0.098  & length \\
$c_{\mathrm{crest}}$    & 0.1   & 0.1    & none \\
$g_\perp$               & 1.0   & 1.0    & length$^2\cdot$time \\
$q_{\mathrm{cap}}$      & 0.03  & 0.03   & none \\
robot $g$, $L$, $\tau$   & 0.86, 0.055, 0.02 & 1, 0, 0 & --, time, time \\
dead zone $v_c$, cap     & 0.025, 0.3 & 0.002, 0.3 & length/time \\
\hline
\end{tabular}
\end{table}

\section{Double-Gyre Benchmark}
\label{sec:ch7:dg_bench}

\subsection{Setup}

% Ported from: Draft_14.tex
The double-gyre experiments use the steady field of
Section~\ref{sec:ch3:dg_example} and the nominal formation, with ring radius
$\rho = 0.075$. Trials vary only the initial heading, set to zero in the
noise-free runs and drawn uniformly from $[0, 2\pi)$ in the Monte Carlo
sweeps, for up to 500 steps. The task is to ride the separatrix downstream
and capture the saddle at its end, $\mathbf{p}^*_1 = (0, -0.5)$; a run
that ends at $\mathbf{p}^*_2$ has ridden against the flow and fails. A
run acquires the separatrix at the first of ten consecutive
steps with its centroid within 0.05 of $x = 0$. Noise-free trials stop once
the centroid has stayed within 0.06 of a flow saddle for 20 consecutive
steps, while noisy trials succeed if the centroid comes within 0.06 of
$\mathbf{p}^*_1$. The formation straddles the separatrix when robots lie
on both sides of $x = 0$. Because losing the straddle is the failure
\cite{27} reports, we measure noise tolerance by retention, the fraction
of trials that succeed and never lose the straddle once they have it.

\subsection{Clean Runs}
\label{sec:ch7:disc_clean}

% Ported from: Draft_14.tex
\begin{figure}[htbp]
    \centering
    \includegraphics[width=0.9\textwidth]{ch7_traverse_vs_logic_c_seven.png}
    \caption{$s_1$ tracker versus $D$ tracker from seven matched starts
    ($\bullet$), noise-free, at S1 $(-0.10, 0.30)$, S2 $(0.15, 0.25)$, S3 $(0, 0.35)$,
    S4 $(0, 0)$, S5 $(0, -0.25)$, S6 $(0.15, -0.15)$, and S7 $(0.10, -0.20)$,
    each ending at the downstream saddle ($\times$, squares). Shading marks
    the 0.1 cells from which a start reaches that saddle in at least 90\%
    (darker) or 50\% (lighter) of random headings, over \prelim{8000}
    noise-free starts per tracker.}
    \label{fig:ch7:traverse_vs_logic_c}
\end{figure}

% Ported from: Draft_14.tex
\begin{figure}[htbp]
    \centering
    \includegraphics[width=0.9\textwidth]{ch7_separatrix_formation_seven.png}
    \caption{The six robots from start S1 (star) under both trackers,
    zoomed to the separatrix $x = 0$ (dashed): robot paths (thin), the
    centroid (heavy), and the formation at four instants.}
    \label{fig:ch7:formation}
\end{figure}

% Ported from: Draft_14.tex
Without noise, both trackers acquire the separatrix, traverse it in the
same direction, and capture the downstream saddle from all seven matched
starts (Fig.~\ref{fig:ch7:traverse_vs_logic_c}). Three starts lie on the
separatrix, at heights $0.35$, $0$, and $-0.25$, and four lie off it by
$1.3\rho$ or $2\rho$, none straddling it at the first step. Both acquire
the separatrix in 0 to 10 steps. Runs from above the origin ride the upper
segment, an attracting OECS, through the isotropic point at the origin and
onto the lower, repelling one, with the formation straddling the
separatrix throughout (Fig.~\ref{fig:ch7:formation}). The $s_1$ tracker
crosses the origin on ordinary tangent selection.

The seven starts sample a wider basin (Fig.~\ref{fig:ch7:traverse_vs_logic_c},
shading). Near mid-height both trackers reach $\mathbf{p}^*_1$ in at least
90\% of headings from any start within \prelim{0.40} of the separatrix, more
than five formation radii, and the basin narrows to \prelim{0.25} at
$|y| = 0.2$. Farther out, the nearest trench need not be the separatrix.

Both hold $\mathbf{p}^*_1 = (0, -0.5)$ in all seven runs. The $D$ tracker
settles 0.005 to 0.006 from it, inside the friction ring of
Section~\ref{sec:ch7:closed_loop}, and the $s_1$ tracker within 0.009.
Where its gain exceeds the bound, the $s_1$ path wiggles about the
separatrix (Fig.~\ref{fig:ch7:formation}), a median 0.004 from it below
$y = -0.19$ from S6 and S7, against under 0.001 for the $D$
tracker.

\subsection{Behavior Under Noise}
\label{sec:ch7:disc_noise}

% Ported from: Draft_14.tex
\begin{figure}[htbp]
    \centering
    \includegraphics[width=0.65\textwidth]{ch7_noise_seven_starts.png}
    \caption{Downstream-saddle success of both trackers from the seven starts of
    Fig.~\ref{fig:ch7:traverse_vs_logic_c}, one per row, \prelim{1000} trials per point,
    with each tracker's retention dashed. Ticks and labels mark the noise
    level at which retention falls to 50\%, and each inset marks the start
    (dot), the separatrix, and $\mathbf{p}^*_1$ (star). (a) Against
    $\sigma_{uv}$, $\sigma_p = 0$. (b) Against $\sigma_p$, $\sigma_{uv} = 0$.}
    \label{fig:ch7:flip_resolution}
\end{figure}

% Ported from: Draft_14.tex
Both primitives sweep measurement and position noise separately from the
seven starts, at \prelim{1000} trials per noise level
(Fig.~\ref{fig:ch7:flip_resolution}), and both degrade smoothly as noise
grows.

The closed form of Section~\ref{sec:ch7:sensitivity} predicts where both must
fail. Each tracker steers on the second-order coefficients, and each pure
second derivative carries a standard deviation of
$2.53\,\sigma_{uv}/\rho^2$, which at $\rho = 0.075$ matches the largest
second derivative in the field, $\pi^3 A \approx 3.1$, near
$\sigma_{uv} = 0.007$. By $\sigma_{uv} = 0.01$ neither tracker keeps the
straddle in more than \prelim{60}\% of trials from any start. The same form
predicts how formation size moves the limit. From S3 and S4, the $D$
tracker's 50\% retention level rises by a factor of \prelim{5.0} and
\prelim{4.7} from $\rho = 0.05$ to $0.075$ and is never reached at
$\rho = 0.1125$, and the $s_1$ tracker's rises by \prelim{2.6} from
$\rho = 0.075$ to $0.1125$. Both beat the factor of 2.25 the coefficient
noise alone predicts, since a wider ring also straddles a larger offset,
so over this range the truncation error has not yet set the limit.

The $D$ tracker tolerates more noise from every start, under both noise
types. Against measurement noise its 50\%
retention level is \prelim{4.4} to \prelim{5.6} times the $s_1$ tracker's above
the origin, \prelim{3.5} times at it, and \prelim{1.3} to \prelim{1.5} times from S6
and S7, and from S5 its retention never falls to 50\%. Against position
noise the ratios are \prelim{2.8} to \prelim{4.9} and \prelim{1.3} to \prelim{1.4}.
Only from S6 and S7 do the retention curves cross, near
$\sigma_{uv} = \prelim{0.005}$. Halving $g_\perp$, which brings the $s_1$
tracker's gain under the bound of Section~\ref{sec:ch7:closed_loop}
everywhere, changes its retention by at most \prelim{2.3} points.

The $s_1$ tracker's main failure is a reversal: at
$\sigma_{uv} = 0.002$, \prelim{96} to \prelim{100}\% of its trials that miss
$\mathbf{p}^*_1$ end at $\mathbf{p}^*_2$. The $D$ tracker re-signs its
direction every cycle against the measured flow, while the $s_1$ tracker,
to stay objective, carries its sign forward. Near the saddles the
along-trench slope $\pi^3 A|\sin\pi y_f|$ that tangent selection relies on
goes to zero, so one noisy step can select the transverse eigenvector and
leave the next sign arbitrary. The chance of a first reversal per riding
step is \prelim{0.13} in the top tenth of the separatrix and
\prelim{0.002} to \prelim{0.004} midway.

The $D$ tracker's smaller margin from S6 and S7 comes from its fitted
Hessian, whose across-trench curvature falls short near the saddles
(Section~\ref{sec:ch7:fit_check}), so noise can make it definite before the
cluster reaches the separatrix. At $\sigma_{uv} = 0.005$ this happens in \prelim{95}\% and
\prelim{86}\% of trials from S6 and S7, which meet the separatrix near
$y = -0.30$, and only \prelim{32}\% and \prelim{38}\% of those trials reach
$\mathbf{p}^*_1$, against \prelim{85}\% and \prelim{81}\% of the rest. From S1
and S2, which meet it inside $|y| < 0.25$, the same happens in \prelim{75}\%
and \prelim{84}\% of trials without changing success.

\subsection{Behavior Under a Rotating Observer}
\label{sec:ch7:rotating}

% Ported from: Draft_14.tex
\begin{figure}[htbp]
    \centering
    \includegraphics[width=0.9\textwidth]{ch7_objectivity_rotating.png}
    \caption{Rotating-observer trial from S1 (star), $10^4$ trials per
    frame, pulled back to inertial coordinates. Lines are mean paths with
    one standard deviation shaded, and squares are median final
    positions. Pentagons show heading-zero runs at $\Omega = \pm 0.18$.
    Only the $D$ tracker's path depends on the frame.}
    \label{fig:ch7:oecs_objectivity}
\end{figure}

% Ported from: Draft_14.tex
A cluster can sense the water without error and still see the flow from a
rotating frame, when each robot adds its own velocity, differenced from
position fixes in drifting axes, to the water velocity it senses. The flow the cluster measures is then
$\mathbf{v}'(\mathbf{p}', t) = \mathbf{Q}\mathbf{v}(\mathbf{Q}^\top \mathbf{p}')
+ \Omega\,(-y', x')^\top$, the physical flow seen in rotated axes plus the transport velocity of a
frame rotating at rate $\Omega$, clockwise for $\Omega > 0$.

Both trackers start from S1 on the noise-free double gyre in the inertial
frame and in frames rotating at $\Omega = \pm 0.10$ and $\pm 0.18$,
with $10^4$ trials of random heading per frame, each stopping once its
centroid stalls or after 500 steps. The frame carries
$\mathbf{p}^*_1$ at $0.5|\Omega|$, and the $s_1$ capture delivers at most
$g\,k\,c_{\max} = 0.103$, so it can hold $\mathbf{p}^*_1$ only while
$|\Omega| < 0.21$. The faster rate stays 13\% inside that
bound.

The $s_1$ tracker follows the same path in every frame
(Fig.~\ref{fig:ch7:oecs_objectivity}, right), with a median largest offset
from the separatrix of 0.010 in the inertial frame and at most
0.015 in the rotating ones. No trial loses the straddle, and every
trial holds $\mathbf{p}^*_1$, ending 0.012 to 0.015 from it.
Its channels and capture test use quantities a rotation leaves invariant,
its tangent seed agrees in every frame from S1, and its across-trench
feedback absorbs the frame's transport velocity, so the match is empirical
rather than exact.

The $D$ tracker's path depends on the frame
(Fig.~\ref{fig:ch7:oecs_objectivity}, left). Its median largest offset is
0.004 in the inertial frame, 0.035 at $|\Omega| = 0.10$, and
0.063 at $|\Omega| = 0.18$, where 28 to 32\% of trials
lose the straddle, the failure \cite{27} reports. The frame's added
vorticity $2\Omega$ shifts $D'$ by (\ref{eq:ch3:D_not_objective}). On the
separatrix the added term $\Omega\omega$ vanishes, but its transverse
gradient $2\pi^3 A\sin\pi y_f$ does not, and against the trench's
curvature $2\pi^6 A^2$ it moves the trench by
$\Omega\sin(\pi y_f)/(\pi^3 A)$, 0.032 and 0.058 at the origin,
close to the measured offsets. The transport velocity also moves the
fitted stagnation point $0.5|\Omega|/(\pi^2 A)$ from $\mathbf{p}^*_1$,
which the hold trails by $0.5|\Omega|/(g\,k)$. At $|\Omega| = 0.10$ the $D$ tracker ends a median
0.072 from $\mathbf{p}^*_1$, against 0.012 for the $s_1$
tracker, and at $0.18$, where the shift reaches 0.091, beyond $\rho$,
it ends 0.138 away.

\subsection{Beyond the Steady Gyre}
\label{sec:ch7:beyond}

% Ported from: Draft_14.tex
\begin{figure}[htbp]
    \centering
    \includegraphics[width=\textwidth]{ch7_beyond_steady.png}
    \caption{Each tracker rides a trench of its own field. Separatrices
    (black), the true trenches of $D$ and $s_1$ (thin), and one noise-free
    ride of each tracker (heavy) from S1 on the gyres and from B2, on the
    separatrix, on the jet; filled end markers reached the downstream
    saddle. (a) On the steady gyre the curves coincide. (b, c) Bending the
    separatrix bows both trenches past it. (d) On the Bickley jet the $s_1$
    trench runs nearly flat below the separatrix.}
    \label{fig:ch7:beyond}
\end{figure}

% Ported from: Draft_14.tex
\begin{table}[htbp]
\caption{Noise-free outcomes beyond the steady gyre, $D$/$s_1$, in
percent of all trials from the seven starts: success, retention against
the separatrix (Ret.) and against the tracker's own true trench (Own),
reversal, and the median fitted $|\hat{\omega}|/2\hat{r}$ along the
ride.}
\label{tab:ch7:beyond}
\centering
\footnotesize
\setlength{\tabcolsep}{2.5pt}
\begin{tabular}{lccccc}
\hline
Field & Succ. & Ret. & Own & Rev. & $|\hat{\omega}|/2\hat{r}$ \\
\hline
Steady gyre                        & 100/100 & 100/100 & 100/100 & 0/0 & \prelim{0.01}/\prelim{0.02} \\
Time-periodic, $\epsilon = 0.1$    & \prelim{97}/\prelim{96} & \prelim{97}/\prelim{61} & & \prelim{3}/\prelim{4} & \prelim{0.11}/\prelim{0.07} \\
Time-periodic, $\epsilon = 0.25$   & \prelim{41}/\prelim{62} & \prelim{13}/\prelim{15} & & \prelim{15}/\prelim{26} & \prelim{0.10}/\prelim{0.20} \\
Bent, $\delta = 0.1$               & \prelim{71}/\prelim{100} & \prelim{67}/\prelim{29} & \prelim{71}/\prelim{70} & 0/$<$1 & \prelim{0.25}/\prelim{0.33} \\
Bent, $\delta = 0.2$               & \prelim{51}/\prelim{62} & \prelim{8}/\prelim{9} & \prelim{50}/\prelim{26} & 0/\prelim{9} & \prelim{0.52}/\prelim{0.53} \\
Bickley jet                        & \prelim{9}/0 & 0/0 & & \prelim{20}/0 & \prelim{0.61}/\prelim{0.60} \\
\hline
\end{tabular}
\end{table}

% Ported from: Draft_14.tex
The steady double gyre is kind to both trackers. Its shear strain
vanishes, its separatrix is straight and still, and the rotation on the
separatrix vanishes, so both trenches lie exactly on it
(Fig.~\ref{fig:ch7:beyond}(a)). Three further fields take these properties
away a few at a time. Each runs without noise from the seven starts, with
\prelim{500} trials of random heading per start (\prelim{300} on the bent
gyre). The time-periodic member of
the family of \cite{52}, with period 10, swings the separatrix sideways at
up to 0.063 for $\epsilon = 0.1$ and 0.157 for $\epsilon = 0.25$ and puts
shear and rotation on it. A bent gyre, constructed here, replaces $x_f$ by
$x_f + \delta\sin\pi y_f$ in the stream function $A\sin\pi x_f\sin\pi y_f$
of (\ref{eq:ch3:dg_u}), which stays steady and incompressible but bows the
separatrix by $\delta$. A Bickley jet \cite{65}
with stream function
$-UL\tanh(y/L) + cy + UL\epsilon\,\mathrm{sech}^2(y/L)\cos k(x + 0.5)$,
$U = 0.3$, $L = 0.2$, $c = 0.3U$, $\epsilon = 0.3$, and $k = 2\pi$, seen in
the frame of its wave, has cat's-eye eddies bounded by U-shaped
separatrices, along most of which rotation exceeds strain.

Each tracker rides a trench of its own field, which marks the separatrix
only where strain outweighs rotation on it (Fig.~\ref{fig:ch7:beyond}). On the bent gyre the trench of $D$ bows up to
0.041 and 0.078 past the separatrix and the trench of $s_1$ up to 0.085
and 0.155. Scored against its own trench, the $s_1$ tracker keeps the
straddle in \prelim{70}\% of trials at $\delta = 0.1$ instead of
\prelim{29}\% (Table~\ref{tab:ch7:beyond}), so it has not lost its curve but
rides a different one. Retention against the
separatrix falls as the median fitted $|\hat{\omega}|/2\hat{r}$ along the
ride rises, from \prelim{97}\% near 0.1 to \prelim{29} to \prelim{67}\% near 0.3
and at most \prelim{9}\% above 0.5. On the Bickley jet the $s_1$ rides end a
median \prelim{0.007} from their own nearly flat trench, and the $D$
tracker, with no trench near much of it, succeeds in \prelim{9}\% of
trials.

Two failures have other causes. On the time-periodic gyre at
$\epsilon = 0.1$ the $s_1$ tracker keeps the straddle in only \prelim{61}\%
of trials, and every successful trial that loses it does so within 0.1 of
the crest, where the strain on the separatrix falls to
$\pi A\epsilon = 0.031$. There $s_1$ is nearly flat across the
moving trench, and its slope-driven channel cannot keep up, while the $D$
tracker's Newton step, a distance, does not weaken. At $\epsilon = 0.25$
the separatrix outruns the $g\,k\,c_{\max} = 0.103$ either channel
can deliver. On the bent gyre the $D$ tracker also fails from S6 and S7,
which the bend leaves 3.2$\rho$ and 2.4$\rho$ from the
separatrix, by acquiring the nearer bottom-wall trench instead.

Because $|\hat{\omega}|/2\hat{r}$ uses only first-order coefficients, a
cluster can compute it as it rides and flag an unreliable trench, provided
its heading is referenced to an inertial frame.

\section{Santa Barbara Channel Trial}
\label{sec:ch7:ocean}

\subsection{Setup}
\label{sec:ch7:ocean_sim}

% Ported from: Draft_14.tex
The ocean experiment uses the public HFRNet U.S.~West Coast surface-current
product for the Santa Barbara Channel on its 2~km grid \cite{66,67}, with
the record of \cite{27}: 29 hourly frames from May 16, 2012 08:00 GMT to
May 17, 2012 12:00 GMT, in the same channel and month used by \cite{6}.
Radar coverage is partial, with 74 to 87\% of the water cells valid in
each frame. Missing cells, including land inside the coverage, are filled
by linear interpolation, the field is zero beyond valid data, and
interpolation is bilinear in space, linear in time.

The simulation maps $[-0.65, 0.65]^2$ onto a square centered on
$(34.2^{\circ}\mathrm{N}, 120.4^{\circ}\mathrm{W})$, with one length unit
of 51.4~km and one time unit of 6000~s, a uniform dilation once longitude
is scaled by $\cos(34.2^{\circ})$. Each control step advances
the field by 600~s, so the 168-step trial covers the record and any vessel
delay or lag is negligible (Table~\ref{tab:ch7:params}). The speed cap
$\sqrt{2}\,k\,c_{\max}$ is then 0.78~m/s, within reach of a small
autonomous surface vessel. Both trackers start north of the channel
entrance at $(34.38^{\circ}\mathrm{N}, 120.41^{\circ}\mathrm{W})$ with
heading $335^\circ$. A run ends at the end of the record or at landfall,
when any robot lies over land on the Natural Earth 1:10 million coastline,
so filled land cells reach only fits near the coast.

An offline FTLE field, available to neither controller, serves as the
Lagrangian reference \cite{52}. It advects a seed grid four times finer
than the data by fourth-order Runge-Kutta at 600~s steps, with land held
at zero velocity. The water points at or above the 95th percentile of the
24-h forward field anchored at the record start form a high-FTLE band that
contains the dominant ridge (Fig.~\ref{fig:ch7:ocean_ftle}), and a distance to
the ridge here means a distance to this band. Recomputing a 12-h field at
four anchor times shows that this ridge persists while the finer structure
downstream evolves. The pentagon is enlarged to $\rho = 0.098$, 5.0~km,
so that it spans the several grid cells a quadratic fit needs. The radius $\rho$ and gain $k$ are tuned
against the FTLE reference, as they set when the cluster reaches the
evolving ridge's bifurcation and so which branch it meets.

\subsection{Results}
\label{sec:ch7:disc_ocean}

% Ported from: Draft_14.tex
\begin{figure}[htbp]
    \centering
    \includegraphics[width=\textwidth]{ch7_ocean_progression_2km.png}
    \caption{$D$ and $s_1$ tracker paths from a single shared start at
    (a) 7, (b) 14, (c) 21, and (d) 28~h, each drawn up to that time, over
    the 24-h forward FTLE field anchored at the record start. Both follow
    the dominant ridge from the channel entrance to the bifurcation and take
    the same branch past it.}
    \label{fig:ch7:ocean_ftle}
\end{figure}

% Ported from: Draft_14.tex
Both trackers follow the dominant ridge of the forward FTLE field from the
channel entrance to the bifurcation and take the same branch past it
(Fig.~\ref{fig:ch7:ocean_ftle}). This agreement is not built
in, as the field is divergent and not vorticity-free. Along both paths
the fit is strain-dominated, with $\hat{D} < 0$ on about 80\% of
steps and a median $|\hat{\omega}|/2\hat{r}$ of 0.15 to 0.20 from 7 to 21~h.
$D$ is then close to $s_1 s_2 = \mu^2 - r^2$, where divergence enters only
as $\mu^2$, a tenth to a fifth of $r^2$. Divergence enters
$s_1 = \mu - r$ linearly, however, so the 1.3~km mean separation of the
two paths is an observation on this record.

The $s_1$ tracker also types what it rides. Its tangent followed
$\mathbf{e}_2$, an attracting OECS, on 34 of 42 steps from 7 to 14~h, and
$\mathbf{e}_1$ on 37 of 42 steps from 14 to 21~h. A 12-h backward FTLE
field anchored at 14~h marks attracting structures and checks both
segments. Over 7 to 14~h the $s_1$ path ran a mean 1.7~km from its
top 5\%, against 5.6~km for the $D$ path. Over 14 to 21~h it ran
6.5~km from that band and 2.0~km from the forward one. In this
divergent field a trench along $\mathbf{e}_1$ is a repelling OECS only
where it is also a ridge of $s_2$, which held on 49\% of the later steps.

The two trackers differ in speed early on. From 2 to 6.5~h the $s_1$
tracker crossed a patch near the entrance where rotation and divergence
rival the strain ($\hat{D} > 0$ on 70\% of steps), switching tangents 12
times in 27 steps, and in the first 7~h it made 2.1~km of net progress
against 8.4~km for the $D$ tracker. Past the patch both kept pace, a mean
1.6 and 2.1~km from the ridge. The $D$ tracker reached land at 23~h. The
$s_1$ tracker drifted 3.8~km off the weakening ridge in the southern
channel, as the robots of \cite{27} also did, and remained at sea at 28~h.

Small changes to the shared start show which of these behaviors are
reliable. Over 100 runs that move the start by up to 1~km, turn the
heading by up to $15^\circ$, and scale $\rho$, $k$, and $c_{\max}$ by up
to 5\%, both trackers stay on the same ridge, a median 1.5 and
0.9~km from it, but ride it in different directions. The $D$
tracker signs its direction against the measured flow every cycle, so it
always rides downstream: 99 of its runs reach the southern channel
and 66 land at the island. The $s_1$ tracker signs each tangent
against the one before it, and in the patch near the entrance that chain
breaks. There the strain eigenframe swaps, and a swap compares the new
tangent with a perpendicular old one, so its sign is close to a coin toss.
In 98 runs the $s_1$ tracker turned against the flow at a median
2.5~h, where the fitted $|\hat{\omega}|/2\hat{r}$ had risen to
0.46 against 0.27 along the $D$ rides, then rode the ridge
north to the mainland. The shared start of Fig.~\ref{fig:ch7:ocean_ftle} is
the one setting of the 101 tested in which the tosses came out
downstream. This is the reversal of the double-gyre runs, set off by
rotation instead of noise, and the ratio $|\omega|/2r$ flags where it
happens.

\section{Summary}
\label{sec:ch7:summary}

% Ported from: Draft_14.tex
This chapter has shown that six robots sampling at once can fit a local
quadratic model of a planar flow, which gives the velocity gradient as a
field over the formation. Two surrogate fields read from that fit, the
Jacobian determinant $D$ and the smaller strain eigenvalue $s_1$, have
trenches that, in the tested environments, lie closest to coherent
structures where strain dominates rotation, and on each we built an adaptive navigation primitive that acquires
and rides those trenches from instantaneous measurements, with no
integration horizon and no pre-straddling. Without noise, both traverse
the double-gyre separatrix from every start and hold the downstream
saddle. The $D$ tracker withstands more measurement and position noise
from every start tested, while the $s_1$ tracker, built from objective
quantities, alone holds its path under a rotating observer. On the Santa
Barbara Channel record both acquire the dominant transport corridor, but
under small changes to the start only the $D$ tracker keeps its
direction. Choosing between them is therefore a trade-off: rotational
robustness favors the $s_1$ tracker, while noise tolerance and directional
reliability favor the $D$ tracker. A cluster with an absolute heading
reference, such as a compass or GNSS, gives up little by carrying the $D$
tracker and can trust its own rotation-to-strain ratio.

% ===================================================================
% CHAPTER 8: CROSS-CUTTING FINDINGS
% ===================================================================
% New text throughout, written 2026-09-30 to Waight_style.md.
\chapter{Cross-Cutting Findings}
\label{ch:crosscutting}

Each rung of the ladder of Section~\ref{sec:ch4:ladder} reads more of the
field than the one below it, and it pays for that in robots, in formation
geometry, and in noise. This chapter compares the rungs on one flow, on
the landscapes each can read, and on those costs. Short numerical
experiments supply the evidence where no chapter already does.

\section{One Flow at Three Orders}
\label{sec:ch8:one_flow}

% New text. Checked 2026-11-19 in experiments/thesis_ch8_one_flow.py: the
% shifted field equals the left gyre of Section~\ref{sec:ch3:dg_example} to
% machine precision, and det J vanishes on the diamond.
A cell of the double gyre is a bounded eddy: water circulates about a
center and is held between saddles at the corners of the cell. Two of the
benchmark fields of this dissertation are this one flow.
Chapter~\ref{ch:first_order} tests the first-order estimate on a single
cell centered on the origin,
\begin{equation}
    \mathbf{v} = 0.1\pi\,\bigl(\cos\pi X\,\sin\pi Y,\;
    -\sin\pi X\,\cos\pi Y\bigr),
    \label{eq:ch8:cell}
\end{equation}
and Chapter~\ref{ch:second_order} tracks the separatrix between the two
cells of the steady double gyre of Section~\ref{sec:ch3:dg_example}.
Shifting the left cell by $X = x + 0.5$, with $Y = y$, turns the
double-gyre field into (\ref{eq:ch8:cell}) term for term. The right wall
of the cell, $X = 0.5$, is then the separatrix $x = 0$, and the saddles at
its ends are the flow saddles of Chapter~\ref{ch:second_order}. On the
cell the determinant of the velocity gradient is
$D = \tfrac{1}{2}\pi^4 A^2(\cos 2\pi X + \cos 2\pi Y)$ with $A = 0.1$,
which vanishes on the diamond $|X| + |Y| = 1/2$.

Each rung of the ladder was run on this flow, and
Table~\ref{tab:ch8:one_flow} records what happened.

\begin{table}[htbp]
\centering
\caption{The gyre cell at each rung.}
\label{tab:ch8:one_flow}
\small
\begin{tabular}{>{\raggedright\arraybackslash}p{1.6cm} >{\raggedright\arraybackslash}p{4.0cm} >{\raggedright\arraybackslash}p{7.2cm}}
\hline
Rung & What the cluster reads & What it did on this flow \\
\hline
Zeroth & the speed & It could not tell the center from the saddles, since the speed vanishes at all of them, and it lost every orbit (Section~\ref{sec:ch6:comparison}). \\
First & $\hat{\mathbf{J}}$ at one point, from three robots & The sign of $\det\hat{\mathbf{J}}$ at the first step decided which critical point the cluster reached, and orbits were lost once the formation reached the diamond (Section~\ref{sec:ch6:sim_attraction}). \\
Second & the quadratic fit over the footprint, from six robots & Both trackers rode the separatrix, the right wall of the cell, to the saddle at its end (Section~\ref{sec:ch7:disc_clean}). \\
\hline
\end{tabular}
\end{table}

The determinant is the function behind the last two rows, read at two
orders. At first order the cluster reads one value of $D$ at a point, and
its sign sorts the starts into basins: the attraction law is a Newton step
toward a zero of the fitted flow, which cannot cross the diamond
$\det\hat{\mathbf{J}} = 0$. At second order the cluster reads the
landscape of $D$ over its footprint, and the trenches of that landscape
are separatrices (Section~\ref{sec:ch7:surrogates}). The two readings meet
at one point. On the separatrix $D$ is zero only at the crest, where the
diamond touches the wall (Fig.~\ref{fig:ch8:one_flow}). The crest is also
where the along-trench slope of $D$ vanishes, so it is where the Newton
step of the $D$ tracker stalls and the measured flow carries the cluster
over (Section~\ref{sec:ch7:sep_controller}).

\begin{figure}[htbp]
    \centering
    \includegraphics[width=0.7\textwidth]{ch8_one_flow.png}
    \caption{One gyre cell at three orders, in the cell coordinates of
    (\ref{eq:ch8:cell}). The shading is $D$, and the dotted diamond is
    $\det\mathbf{J} = 0$. Dots are the 1000 starts of the first-order gyre
    trials, blue where the cluster reached the center and orange where it
    reached a saddle. Lines are one noise-free ride of each second-order
    tracker from start S1 of Chapter~\ref{ch:second_order}, shifted into
    cell coordinates; both ride the separatrix, the right wall of the cell,
    to the saddle at its lower corner. Crosses mark the saddles of the cell,
    and the diamond marker is the crest.}
    \label{fig:ch8:one_flow}
\end{figure}

\section{Vector Fields as Scalar Landscapes}
\label{sec:ch8:scalar}

Earlier work in this laboratory navigated scalar fields with a library of
primitives for extrema finding, contour following, ridge and trench
following, and saddle point station keeping \cite{3}. A vector field
is not a scalar field, but it supplies several. Each robot reads the
magnitude $\lVert\mathbf{v}\rVert$ directly. The divergence, the
vorticity, the determinant $D$, and the strain eigenvalue $s_1$ are all
built from the velocity gradient. Each of these is a landscape the
cluster can climb, descend, or follow.

Table~\ref{tab:ch8:scalar} shows what each fit order makes available. A
fit of order $k$ recovers derivatives of the field up to order $k$. A
quantity built from derivatives of order $m$ therefore has a value at
order $m$, a gradient at order $m + 1$, and a curvature at order $m + 2$,
wherever it is differentiable. The magnitude is not differentiable where
the velocity vanishes, which matters at the critical points themselves,
and $s_1 = \mu - r$ is not differentiable at an isotropic point, where
the strain vanishes. A primitive needs the gradient to steer on a landscape and the curvature
to know what shape it is on. The magnitude sits one order below the
quantities built from $\mathbf{J}$. This is why a three-robot fit can
steer on the magnitude but only read the value of $D$ and $s_1$, and why
steering on $D$ or $s_1$ needs six robots. The true curvature of $D$,
which the second-order trackers approximate with the Hessian of the
fitted $\hat{D}$, needs ten.

\begin{table}[htbp]
\centering
\caption{Scalar fields supplied by a planar vector field, and the fit
order needed for each field's value, gradient, and curvature. A field is
objective when an observer rotating relative to the flow reads the same
value (Section~\ref{sec:ch3:relation}).}
\label{tab:ch8:scalar}
\begin{tabular}{llcccc}
\hline
Field & Built from & Value & Gradient & Curvature & Objective \\
\hline
Magnitude $\lVert\mathbf{v}\rVert$ & $\mathbf{v}$ & 0 & 1 & 2 & no \\
Divergence $\operatorname{tr}\mathbf{J}$ & $\mathbf{J}$ & 1 & 2 & 3 & yes \\
Vorticity $\omega$ & $\mathbf{J}$ & 1 & 2 & 3 & no \\
Determinant $D$ & $\mathbf{J}$ & 1 & 2 & 3 & no \\
Strain eigenvalue $s_1$ & $\mathbf{J}$ & 1 & 2 & 3 & yes \\
\hline
\end{tabular}
\end{table}

Critical points enter this table through the magnitude. A critical point
is where the velocity vanishes, so it is a zero of
$\lVert\mathbf{v}\rVert$. Near a nondegenerate critical point
$\mathbf{p}^*$,
\begin{equation}
    \lVert\mathbf{v}\rVert^2 \approx
    (\mathbf{p} - \mathbf{p}^*)^{\top}\mathbf{J}^{\top}\mathbf{J}\,
    (\mathbf{p} - \mathbf{p}^*),
    \label{eq:ch8:mag_quad}
\end{equation}
and $\mathbf{J}^{\top}\mathbf{J}$ is positive definite whenever
$\det\mathbf{J} \neq 0$, whatever the type. Every nondegenerate critical
point is therefore a minimum of the magnitude, saddles included.
Descending the magnitude reaches a critical point but cannot say which
type it reached, since a saddle and a center can present the same
magnitude landscape. The first-order estimate of
Chapter~\ref{ch:first_order} attacks the same zero directly. Its step
$\hat{\mathbf{p}}^* - \mathbf{p}_c = -\hat{\mathbf{J}}^{-1}\hat{\mathbf{v}}_c$
is one Newton step toward $\mathbf{v} = \mathbf{0}$, and the same
$\hat{\mathbf{J}}$ names the type.

A closed-loop test confirms this with the three-robot simulator of
Chapter~\ref{ch:first_order}, on the robot model of
Section~\ref{sec:ch2:dynamics}. The vector-to-scalar primitive of
Section~\ref{sec:ch5:vector_to_scalar} descended the magnitude from 1000
random starts in a 1~m box around a saddle. On the canonical saddle of
Appendix~\ref{app:fields} it finished within 0.005~m of the saddle in
every trial, the same as on the vortex, whose magnitude landscape is
identical. On a saddle with rates 1 and $-1/3$ the magnitude is an
elliptic cone with a sharp point at the saddle, and a plane fit across
that point is biased. The cluster then settled a median of 0.048~m from
the saddle, with every trial within 0.05~m. The attraction law of
Chapter~\ref{ch:first_order} finished within 0.024~m in every trial on
all three fields, at the edge of the friction ring $v_c/k = 0.025$~m of
Section~\ref{sec:ch6:friction} and in agreement with the 22.6~mm mean
final error of Table~\ref{tab:ch6:simulation_results}
(Fig.~\ref{fig:ch8:magnitude}).

\begin{figure}[htbp]
    \centering
    \includegraphics[width=0.55\textwidth]{ch8_magnitude_saddle.png}
    \caption{Magnitude descent (blue) and the attraction law (red) from ten
    matched starts on a saddle with rates 1 and $-1/3$, simulated with the
    robot model of Chapter~\ref{ch:first_order}. The cross marks the
    saddle. Magnitude descent follows the steepest slope of the elliptic
    magnitude cone and settles a few centimeters from the saddle, while the
    attraction law runs nearly straight to it.}
    \label{fig:ch8:magnitude}
\end{figure}

\section{Sampling Geometry}
\label{sec:ch8:geometry}

A fit of order $k$ has $n_k = \tfrac{1}{2}(k+1)(k+2)$ coefficients per
component (Section~\ref{sec:ch4:minimality}), so it needs 3, 6, 10, and 15
robots for $k = 1$ to 4. With exactly $n_k$ robots the fit fails when the
robots lie on a common algebraic curve of degree at most $k$, a line for the
first order, a conic for the second, and a cubic for the third.
Interpolation theory calls a point set that avoids this unisolvent.

The failure has a physical reading. When the robots share such a curve,
the null space of $\boldsymbol{\Phi}_k$ holds the curve's own
coefficients. Any multiple of that polynomial can be added to the fitted
field without changing a single reading, so the robots cannot tell the
true field from the true field plus the curve. A regular hexagon lies on
its circumscribed circle, and its quadratic fit is singular at any size.
Adding robots to the same circle does not help, since each of them reads
the same ambiguity. Nine robots on a ring with a tenth at the center fail
a cubic fit in two ways at once. Their null space is two-dimensional,
the circle times any line through the center robot.

A formation that avoids these curves can still come close to them, and
its condition number indicates how much the fit amplifies noise.
Table~\ref{tab:ch8:geometry} lists the formations of this dissertation and
candidates for a cubic fit, with $\kappa$ computed on the
radius-normalized formation matrix of Section~\ref{sec:ch7:sensitivity}.
The equilateral triangle gives $\kappa(\boldsymbol{\Phi}_1) = 1.41$ and
the pentagon plus center gives $\kappa(\boldsymbol{\Phi}_2) = 8.26$. Ten
robots are harder to place. A triangular lattice of 1, 2, 3, and 4 robots
supports a cubic fit but gives $\kappa = 216$. A search over two-ring
layouts found $\kappa = 44.3$ for an inner pentagon at 0.65 of the outer
radius, rotated $36^\circ$ against the outer pentagon. The best cubic
layout found therefore has a condition number about five times that of
the pentagon plus center, before formation size enters at all.

\begin{table}[htbp]
\centering
\caption{Condition numbers of the radius-normalized formation matrix.
Singular formations lie on a common curve of the fit's degree.}
\label{tab:ch8:geometry}
\small
\begin{tabular}{lccl}
\hline
Formation & Order & Robots & $\kappa(\boldsymbol{\Phi}_k)$ \\
\hline
Equilateral triangle & 1 & 3 & 1.41 \\
Pentagon plus center & 2 & 6 & 8.26 \\
Regular hexagon & 2 & 6 & singular (circle) \\
Hexagon plus center & 2 & 7 & 8.84 \\
Nine-robot ring plus center & 3 & 10 & singular (circle $\times$ line) \\
Triangular lattice 1+2+3+4 & 3 & 10 & 216 \\
Inner triangle, outer hexagon, center & 3 & 10 & 211 \\
Two rotated pentagons, radii 0.65 and 1 & 3 & 10 & 44.3 \\
\hline
\end{tabular}
\end{table}

\section{Noise Growth with Fit Order}
\label{sec:ch8:noise_order}

A higher-order fit reads finer structure from the same readings, and it
pays for that structure in noise. A coefficient of order $q$ is a $q$-th
derivative of the field, estimated from differences between readings a
formation radius apart. Each additional derivative divides by another
factor of the radius.

Let a formation of fixed shape $\boldsymbol{\xi}_i$ be scaled to radius
$\rho$, so robot $i$ sits at $\rho\,\boldsymbol{\xi}_i$. Each entry of the
basis of Section~\ref{sec:ch4:basis} is a monomial of degree $|\alpha|$,
so the formation matrix factors as
\begin{equation}
    \boldsymbol{\Phi}_\rho = \boldsymbol{\Phi}_1\,\mathbf{S}_\rho, \qquad
    \mathbf{S}_\rho = \operatorname{diag}\bigl(\rho^{|\alpha|}\bigr).
    \label{eq:ch8:phi_scaling}
\end{equation}
With independent measurement noise of standard deviation $\sigma$ on each
reading, the fitted coefficient of degree $q$ has standard deviation
\begin{equation}
    \sigma_{\alpha} = \rho^{-q}\,\sigma\,
        \bigl\lVert \mathbf{e}_\alpha^{\top}\boldsymbol{\Phi}_1^{-1}\bigr\rVert,
    \qquad q = |\alpha|,
    \label{eq:ch8:coef_noise}
\end{equation}
where the row norm depends on the formation's shape alone. This scaling is
exact for a formation of fixed shape, independent noise of equal variance
on every reading, and a full-rank $\boldsymbol{\Phi}_1$.

What matters at the terminal step is the error in whatever quantity the
controller drives to zero. An error $\delta f$ in that quantity moves its
estimated zero by $\delta n \approx \delta f/|\partial f/\partial n|$,
where $n$ is the direction in which $f$ changes. If $f$ is built from
derivatives up to order $m$, then by (\ref{eq:ch8:coef_noise}) its error
$\delta f$ scales as $\rho^{-m}\sigma$, and
\begin{equation}
    \delta n \sim \frac{\rho^{-m}\,\sigma}{|\partial f/\partial n|}.
    \label{eq:ch8:terminal_noise}
\end{equation}
Unlike (\ref{eq:ch8:coef_noise}), this relation is an approximation to
first order in the noise, since $f$ is nonlinear in the coefficients. The
first-order law drives the velocity itself to zero ($m = 0$). With
the centroid on the critical point, the fit's value there is the mean of
the three readings, so its error has standard deviation $\sigma/\sqrt{3}$
per component. This recovers the estimate $\sigma_p \approx
\lVert\mathbf{J}^{-1}\rVert\,\sigma_v/\sqrt{3}$ of
Section~\ref{sec:ch6:error}, which does not depend on $\rho$. Both
second-order trackers drive a transverse gradient to zero,
$\mathbf{w}_2^{\top}\nabla\hat{D}$ or $\mathbf{P}\nabla\hat{s}_1$, which
is built from second derivatives ($m = 2$). Their terminal error should
therefore grow as $\rho^{-2}$.

The denominator of (\ref{eq:ch8:terminal_noise}) is the field's share of
the error, the inverse slope of the quantity being zeroed. For the
critical point it is $\lVert\mathbf{J}^{-1}\rVert$, which diverges as
$\det\mathbf{J} \to 0$, the $\kappa(\mathbf{J})$ effect of
Section~\ref{sec:ch6:error}. For the $D$ tracker it is $1/|\lambda_2|$,
the inverse transverse curvature that already divides its across-trench
command, and it diverges as the trench flattens. The formation's share,
$\kappa(\mathbf{A})$ or $\kappa(\boldsymbol{\Phi})$, is set by the robot
positions alone.

A Monte Carlo test on the double gyre checks both predictions
(Fig.~\ref{fig:ch8:radius}). A triangle was centered on a gyre center and
on a saddle, and a pentagon plus center on the separatrix at
$(0, 0.25)$. The radius ran from a quarter to twice the nominal 0.075,
with a random heading and measurement noise $\sigma = 0.002$ on every
reading, over 1000 trials per radius. Each trial was fit twice, with and
without noise, so the difference isolates the noise. At first order the
scatter of the estimated critical point was flat in radius, with a
log-log slope of 0.01 at the gyre center and 0.02 at the saddle, and it
matched $\lVert\mathbf{J}^{-1}\rVert\sigma/\sqrt{3}$ to within 6\%. The
transverse errors of the $D$ and $s_1$ trackers grew with slopes of
$-2.04$ and $-2.01$.

Truncation bias runs the other way. The noise-free fits carried an RMS
error that grew as $\rho^{3.0}$ at first order and as $\rho^{0.95}$ and
$\rho^{0.98}$ for the two trackers. At first order the bias overtook the
noise at about 1.5 times the nominal radius, and for the $D$ tracker the
two were nearly equal at twice the nominal radius.

At every radius the transverse error of the $s_1$ tracker was about 3.6
times that of the $D$ tracker. Section~\ref{sec:ch7:disc_noise} traces
the $s_1$ tracker's failures under noise to the sign of its tangent,
which is selected from the same gradient. The ordering matches the
retention levels of Section~\ref{sec:ch7:disc_noise}, where the $D$
tracker withstood \prelim{4.4} to \prelim{5.6} times the measurement noise
of the $s_1$ tracker from the starts above the origin. The two ratios
measure different things. The first compares the fit errors at one point
of the separatrix at a fixed noise level. The second compares the noise
levels at which each tracker's retention falls to 50\%, and acquisition
and reversals enter only the second.

\begin{figure}[htbp]
    \centering
    \includegraphics[width=0.95\textwidth]{ch8_radius_sweep.png}
    \caption{Terminal error against formation radius on the double gyre,
    1000 trials per point. Left, the noise alone at $\sigma = 0.002$. The
    first-order critical point estimate is flat in radius, while the
    transverse errors of both second-order trackers fall as $\rho^{-2}$.
    Right, the RMS truncation error of the noise-free fits.}
    \label{fig:ch8:radius}
\end{figure}

Each rung has a signal-to-noise ratio of its own. The signal is the next
derivative of the field times the footprint raised to that order, and
Table~\ref{tab:ch8:snr} sets it against the noise on it and gives the
threshold that each chapter found. The thresholds are not comparable one
to one: the first-order condition is a validity condition for a
covariance prediction and not a failure point, and the zeroth rung has no
measured threshold, so none is stated. The zeroth rung has a sharp
consequence all the same. A direction read from a vector becomes unreliable
where the vector vanishes, and the vector vanishes at the critical points
that the first rung was built to find.

\begin{table}[htbp]
\centering
\caption{The signal each rung reads, the noise on it, and the threshold
found for that rung.}
\label{tab:ch8:snr}
\small
\begin{tabular}{>{\raggedright\arraybackslash}p{1.5cm} >{\raggedright\arraybackslash}p{3.0cm} >{\raggedright\arraybackslash}p{2.0cm} >{\raggedright\arraybackslash}p{5.7cm}}
\hline
Rung & Signal & Noise on it & Threshold found \\
\hline
Zeroth & $\lVert\mathbf{v}\rVert$ & sensing error & None measured. Small vectors gave large direction shifts on the RGB testbed (Section~\ref{sec:ch2:hsv}). \\
First & $\sigma_{\min}(\mathbf{J})R$ & $\sigma_v$ & The covariance prediction of $\bar{e}_f$ holds within 5\% when $\sigma_{\min}(\mathbf{J})R \geq 3\sigma_v$ (Section~\ref{sec:ch6:discussion}). \\
Second & $\pi^3 A\rho^2$, the largest second derivative of the gyre times $\rho^2$ & $2.53\,\sigma_{uv}$ & The two are equal near $\sigma_{uv} = 0.007$. By $\sigma_{uv} = 0.01$ neither tracker keeps the straddle in more than \prelim{60}\% of trials from any start (Section~\ref{sec:ch7:disc_noise}). \\
\hline
\end{tabular}
\end{table}

\section{Steady Offsets: Speed over Loop Gain}
\label{sec:ch8:offsets}

% New text. Every row of the table is a result of Chapters 5 to 7 or a
% statement of Chapter 9; nothing here was run for this chapter, except the
% fit of the vector-sum drift in Section~\ref{sec:ch6:limit_cycle}.
Every primitive of this dissertation closes a loop around a target, and
every one settles away from it for the same reason. When some speed
reaches the loop uncancelled, the cluster settles offset from its target by
that speed divided by the loop gain. If the speed exceeds the largest
command the loop can issue, the cluster loses the target.
Table~\ref{tab:ch8:offsets} lists the speeds that reached the loop in the
experiments.

\begin{table}[htbp]
\centering
\caption{Speeds that the loop does not cancel, the loop gain that opposes
them, and the offset or loss that follows. $T$ is the response delay of
Section~\ref{sec:ch6:orbits}.}
\label{tab:ch8:offsets}
\footnotesize
\begin{tabular}{>{\raggedright\arraybackslash}p{3.4cm} c >{\raggedright\arraybackslash}p{5.4cm} >{\raggedright\arraybackslash}p{2.4cm}}
\hline
Speed the loop does not cancel & Loop gain & Result & Where \\
\hline
Dead zone $v_c$ & $k$ & The cluster settles inside the friction ring $v_c/k$: 0.025~m in Chapter~\ref{ch:first_order} and 0.008 in the hold of Chapter~\ref{ch:second_order}. & Sec.~\ref{sec:ch6:friction}, \ref{sec:ch7:closed_loop} \\
Outward drift from late turns, $TV^2/r^*$ & $gk_r$ & Limit cycle, $r^* - r_d = (TV^2/r^*)/(gk_r)$. & Sec.~\ref{sec:ch6:orbits} \\
The same drift with no radial loop & $0$ & $r^2$ grows at $2TV^2$ without bound, as the vector-sum primitive drifted. & Sec.~\ref{sec:ch5:trials}, \ref{sec:ch6:limit_cycle} \\
A critical point moving at speed $V$ & $gk$ & The cluster trails it by $V/(gk)$ (stated, not tested). & Sec.~\ref{sec:ch9:limitations} \\
Frame transport, $0.5\lvert\Omega\rvert$ & $gk$ & The hold trails by $0.5\lvert\Omega\rvert/(gk)$. The $s_1$ hold is lost once $0.5\lvert\Omega\rvert > gkc_{\max}$. & Sec.~\ref{sec:ch7:rotating} \\
A separatrix moving faster than $gkc_{\max} = 0.103$ & saturated & Both channels lose it. & Sec.~\ref{sec:ch7:beyond} \\
\hline
\end{tabular}
\end{table}

The friction ring is the simplest case. The dead zone ignores any command
below $v_c$, so the cluster stops where $k\lVert\mathbf{p}_c -
\hat{\mathbf{p}}^*\rVert = v_c$. The orbit's limit cycle has the same
form. The late turn drifts the cluster outward at $TV^2/r^*$, and the
radial loop of gain $gk_r$ balances it at
$r^* - r_d = (TV^2/r^*)/(gk_r)$, which is (\ref{eq:ch6:orbit_cf})
rearranged. A primitive with no radial loop leaves the drift uncorrected.

All of these settle a loop whose gain is bounded above. The sampled loop
is stable only below $1/(gL) = 21.1$ per unit time, or 16.1 with the lag
included (Section~\ref{sec:ch7:closed_loop}). At first order the designer
sets the gain, which is $k$. At second order the field sets part of it: $k$
times the ratio of true to fitted curvature for the $D$ tracker, and
$k g_\perp\,\partial^2 s_1/\partial n^2$ for the $s_1$ tracker. This is why
the $s_1$ path wiggles below $y = -0.19$, where its gain passes the bound.

\section{Scalability}
\label{sec:ch8:scalability}

More robots than $n_k$ turn the fit into least squares, and the extra
readings average down the noise. How much they help depends on where the
robots go (Fig.~\ref{fig:ch8:gain_vs_N}). For a ring with a center robot,
adding ring robots leaves the noise on the centroid value unchanged and
barely reduces the noise on the second-order coefficients, whose gain
falls only as $N^{-0.22}$. Every ring robot sits on one circle, so the
ring cannot separate the constant term from the curvature term
$x^2 + y^2$, and the center robot carries both alone. With the robots
split between two rings and a center, the gains fall as $N^{-0.46}$ to
$N^{-0.50}$ for $N \geq 10$, close to the $N^{-1/2}$ of independent
averaging, and the condition number holds near 7.4.

\begin{figure}[htbp]
    \centering
    \includegraphics[width=0.6\textwidth]{ch8_gain_vs_N.png}
    \caption{Noise gain of the quadratic fit's coefficients against the
    number of robots, at radius 1. Dashed, a ring with a center robot.
    Solid, a center robot and two rings, shown from ten robots, where each
    ring holds at least four. The dotted lines fall as $N^{-1/2}$.}
    \label{fig:ch8:gain_vs_N}
\end{figure}

The same geometry decides whether a cluster survives losing a robot. A
formation with exactly $n_k$ robots cannot lose any of them. A ring with a
center robot keeps its quadratic fit when any ring robot fails, but
losing the center leaves the rest on a circle, which is fatal. Every
two-ring formation from seven to fifteen robots kept its quadratic fit
after the loss of any single robot.

The cost of more robots falls on the control architecture. Cluster space
control is centralized, requires global state, and becomes intractable
for suitably large clusters \cite{53}. Each added robot adds two
cluster variables and two rows to the inverse kinematic Jacobian of
Section~\ref{sec:ch4:cluster_space}. The claims of
Chapter~\ref{ch:second_order} are for six robots with ample bandwidth
(Section~\ref{sec:ch9:second_order}), and a ten-robot cubic cluster would
carry the same limitation with more state to share.

\section{Fitted Features and True Features}
\label{sec:ch8:fitted_true}

% New text.
A cluster steers on a feature of its fit, and the fit stands in for the
flow. Each rung therefore has two features: the one the cluster steers on,
which belongs to the fitted polynomial, and the one it wants, which belongs
to the flow. Table~\ref{tab:ch8:fitted_true} gives, for each rung, when the
two coincide and where they came apart.

\begin{table}[htbp]
\centering
\caption{The feature each rung steers on, the condition under which it is
the true feature, and where it was seen to fail.}
\label{tab:ch8:fitted_true}
\small
\begin{tabular}{>{\raggedright\arraybackslash}p{1.4cm} >{\raggedright\arraybackslash}p{3.0cm} >{\raggedright\arraybackslash}p{3.4cm} >{\raggedright\arraybackslash}p{4.6cm}}
\hline
Rung & Steers to & It is the true feature when & It fails visibly on \\
\hline
Zeroth & the minimum of the fitted magnitude plane & the speed grows equally in every direction & the anisotropic node, where the primitive reached the critical point in 6\% of trials, and the gyre, where the speed vanishes at the center and the saddles alike (Section~\ref{sec:ch6:comparison}). \\
First & the zero of the fitted affine field & the field is affine from the cluster out to the critical point & the gyre and the rotational saddle, where the estimate stayed 2.6 to 4.9~mm from the true critical point and the sign of $\det\hat{\mathbf{J}}$ set the basin (Section~\ref{sec:ch6:sim_attraction}). \\
Second & the trench of the fitted $D$ or $s_1$ & the rotation on the separatrix is small against the strain, which is necessary but not sufficient & the bent gyre and the Bickley jet, where each tracker rides the trench of its own field (Section~\ref{sec:ch7:beyond}). \\
\hline
\end{tabular}
\end{table}

The condition is a statement about the flow, which the cluster cannot see.
What it can do is compute quantities from its own fit that signal when the
condition is likely to fail. Table~\ref{tab:ch8:indicators} lists them.

\begin{table}[htbp]
\centering
\caption{Quantities a cluster can compute about its own fit.}
\label{tab:ch8:indicators}
\small
\begin{tabular}{>{\raggedright\arraybackslash}p{3.3cm} >{\raggedright\arraybackslash}p{1.7cm} >{\raggedright\arraybackslash}p{3.4cm} >{\raggedright\arraybackslash}p{4.0cm}}
\hline
Indicator & Order & Signals & Tested \\
\hline
$\sigma_{\min}(\hat{\mathbf{J}})$ & first & a near-singular critical point estimate, since model error is amplified by $1/\sigma_{\min}$ & Stated in Section~\ref{sec:ch6:discussion}. The attraction law does not gate on it. \\
$\lvert\hat{\alpha}\rvert/\hat{\nu}$ against the center threshold & first & a type label near a boundary & The labels were scored at three thresholds (Table~\ref{tab:ch6:classification_hw}). \\
$\lvert\hat{\omega}\rvert/2\hat{r}$ & first & a trench unlikely to mark a separatrix & Yes: retention falls as the ratio rises (Section~\ref{sec:ch7:beyond}). It needs a heading referenced to an inertial frame. \\
The least-squares residual & any order with more than $n_k$ robots & measurement noise, where the truncation error is small & Not tested (Section~\ref{sec:ch9:higher_order}). \\
\hline
\end{tabular}
\end{table}

\section{Lower Rungs Inside Higher Ones}
\label{sec:ch8:nesting}

% New text. The decision rule below is a synthesis of the two trackers of
% Chapter 7, not a tested result.
The rungs also nest in use. The $D$ tracker is a composition of all three.
It makes its stop decision with the zeroth- and first-order coefficients,
the least noisy that the fit returns: once the fitted flow has a saddle
inside the formation, it latches the attraction law of
Chapter~\ref{ch:first_order}. It signs its direction against the measured
flow, a zeroth-order quantity. Only its directions and steps along and
across the trench, the eigenvectors and curvature of the fitted Hessian of
$D$, are of second order. The $s_1$ tracker nests the other way. It chooses
between the strain eigenvectors $\mathbf{e}_1$ and $\mathbf{e}_2$, a
first-order quantity, by the gradient of $s_1$, a second-order quantity,
and it takes the sign of that choice from memory. Its reversals therefore
concentrate where the along-trench slope goes to zero
(Section~\ref{sec:ch7:disc_noise}), and the future-work rule of
Section~\ref{sec:ch9:higher_order}, taking the eigenvector nearest the
previous tangent, would move that choice down to first order.

These two cases suggest a design rule: make each decision with the
lowest-order coefficients that can make it, since the noise on a
coefficient of order $q$ grows as $\rho^{-q}$ (Section~\ref{sec:ch8:noise_order}).
The rule accounts for the sign and stop decisions of the $D$ tracker. It
does not account for the whole noise gap between the trackers, because the
along-trench direction of the $D$ tracker is itself of second order.

The rungs meet in a second way. Critical points are the nodes of a flow,
and separatrices connect them. On the steady double gyre the separatrix is
a heteroclinic connection, since the flow runs down it from
$\mathbf{p}^*_2$ to $\mathbf{p}^*_1$. The $D$ tracker travels that
connection and captures the node at its end with the test and the law of
Chapter~\ref{ch:first_order}. Section~\ref{sec:ch9:features} lists
heteroclinic orbits as a feature that needs a second-order fit. On this one
steady flow Chapter~\ref{ch:second_order} has done it, and
Section~\ref{sec:ch7:beyond} shows that the time-periodic and jet cases
do not hold up as cleanly.

\section{Classification and Frames Across the Rungs}
\label{sec:ch8:class_frames}

% New text. The saddle-read-as-center statement follows from
% D' = D + Omega*omega + Omega^2 (Appendix D); with omega = 0 and a pure
% strain s, D' = Omega^2 - s^2, and J' has eigenvalues +/- i sqrt(Omega^2 - s^2)
% for Omega > s. Checked numerically 2026-11-19.
Every rung is asked what it is looking at. The zeroth rung cannot say: the
speed looks the same at a saddle and a center
(Section~\ref{sec:ch8:scalar}). The first names the type from the
eigenvalues of $\hat{\mathbf{J}}$ (Chapter~\ref{ch:first_order}). The
second types what it rides. At sea the tangent of the $s_1$ tracker
followed $\mathbf{e}_2$, an attracting OECS, on 34 of 42 steps from 7 to
14~h, and a backward FTLE field checked it (Section~\ref{sec:ch7:disc_ocean}).
The hard cases sit on the classifier's boundaries at both rungs. A center
has $\operatorname{tr}\mathbf{J} = 0$, and on hardware only 45\% and 47\%
of its cycles were labeled strictly as a center with three and four robots
(Section~\ref{sec:ch6:hw_class}). At an isotropic point, where the strain
vanishes, the tangent of the $s_1$ tracker swaps eigenvector and its
reversals begin.

Table~\ref{tab:ch8:scalar} marks which fields are objective, and the
primitives need the same column (Table~\ref{tab:ch8:objectivity}). A
first-order fit is not objective either. With no vorticity and a pure
strain $s$, an observer rotating at $\Omega$ reads $D' = \Omega^2 - s^2$
(Appendix~\ref{app:frame_equivariance}), so one that rotates faster than
the strain rate reads an irrotational saddle as a center.

\begin{table}[htbp]
\centering
\caption{Objectivity of the primitives. A primitive is objective when
observers rotating relative to each other agree on what it steers on.}
\label{tab:ch8:objectivity}
\small
\begin{tabular}{>{\raggedright\arraybackslash}p{3.4cm} >{\raggedright\arraybackslash}p{3.8cm} >{\raggedright\arraybackslash}p{5.2cm}}
\hline
Primitive & Reads & Objective \\
\hline
Vector-sum, vector-to-scalar & the sensed vector and its magnitude & no \\
Attraction, orbital & the fitted affine field & no \\
$D$ tracker & $\hat{D}$ and its Hessian & no \\
$s_1$ tracker & $\hat{s}_1$, its gradient, the strain eigenframe & the features, yes; the command, where its sign test is decisive; the path, only empirically (Appendix~\ref{app:frame_equivariance}) \\
\hline
\end{tabular}
\end{table}

``Instantaneous'' also needs one scope sentence. Estimation is memoryless,
but control is not: the $D$ tracker latches its stop test, and the $s_1$
tracker carries one sign bit. That bit buys its objectivity and is the
source of its reversals.

\section{Design Guidelines}
\label{sec:ch8:guidelines}

% New text. The ratios 0.14 and 0.11 are arithmetic on stated numbers:
% 0.025 m over the center-to-robot distance of the 0.30 m triangle,
% 0.30/sqrt(3) = 0.173 m, and 0.008 m over the pentagon radius 0.075 m.
The thresholds of the second-order trackers are stated in formation radii,
which is why the same trackers ran on the nondimensional double gyre and,
after a uniform dilation, on the Santa Barbara Channel, where one length
unit is 51.4~km (Section~\ref{sec:ch7:ocean_sim}).
Table~\ref{tab:ch8:groups} collects the dimensionless groups behind the
results of this dissertation.

\begin{table}[htbp]
\centering
\caption{Dimensionless groups behind the results of this dissertation.}
\label{tab:ch8:groups}
\small
\begin{tabular}{>{\raggedright\arraybackslash}p{4.4cm} >{\raggedright\arraybackslash}p{4.0cm} >{\raggedright\arraybackslash}p{4.0cm}}
\hline
Group & Governs & Where \\
\hline
$\rho$ over the length on which the field curves & truncation bias & Sections~\ref{sec:ch4:size}, \ref{sec:ch8:noise_order} \\
Signal over noise, by rung & noise failure & Table~\ref{tab:ch8:snr} \\
$gkL$ & stability of the sampled loop & Section~\ref{sec:ch7:closed_loop} \\
$v_c/(k\rho)$ & the friction ring in formation radii, 0.14 of the center-to-robot distance of the triangle in Chapter~\ref{ch:first_order} and 0.11 of the pentagon radius in Chapter~\ref{ch:second_order} & Table~\ref{tab:ch8:offsets} \\
$\lvert\omega\rvert/2r$ & whether a trench marks a separatrix & Section~\ref{sec:ch7:beyond} \\
Feature or frame speed over $gkc_{\max}$ & tracking a moving feature or holding in a rotating frame & Table~\ref{tab:ch8:offsets} \\
\hline
\end{tabular}
\end{table}

Four guidelines follow, each from a result above. First, size the
formation to the field. A larger formation suppresses the noise on a
coefficient of order $q$ as $\rho^{-q}$ and raises the truncation error with
$\rho$, so the radius belongs where the two curves cross
(Section~\ref{sec:ch8:noise_order}). Second, keep the robots off the
degenerate curve and watch the condition number: the pentagon plus center
gives 8.26 for a quadratic fit, and the best ten-robot layout found for a
cubic fit gives 44.3 (Section~\ref{sec:ch8:geometry}). Third, keep the loop
gain below the sampled-loop limit and the friction ring small against the
formation radius, remembering that at second order the field sets part of the
gain (Section~\ref{sec:ch8:offsets}). Fourth, choose the tracker by the
constraint: rotation of the observer favors the $s_1$ tracker, while noise
and directional reliability favor the $D$ tracker, and a cluster with an
absolute heading reference gives up little by carrying the $D$ tracker
(Chapter~\ref{ch:second_order}).

% ===================================================================
% CHAPTER 9: LIMITATIONS AND FUTURE WORK
% ===================================================================
\chapter{Limitations and Future Work}
\label{ch:limitations}

% New text
Each chapter states the limits of its own claims. This chapter collects
them. The limits of the testbed itself are given in
Section~\ref{sec:ch2:limitations}.

\section{Limitations}
\label{sec:ch9:limitations}

\subsection{First-Order Estimation and Control}
\label{sec:ch9:first_order}

% Ported from: cwaight_systems_paper_with_revisions2.tex
Several limitations constrain the current implementation and experimental validation. First, the approach is restricted to planar (2D) fields, though many applications involve three-dimensional flows; extension to 3D fields would require at minimum four non-coplanar robots for Jacobian estimation.

Second, the testbed validates performance only on static fields where critical points remain fixed during operation. The static printed field approach enables precise ground truth validation but does not advect the robots: a physical flow whose outward speed grows with distance from the critical point faster than $gk$ would destabilize the attraction law, and a critical point drifting at constant speed would be trailed by that speed divided by $gk$.

Third, robot team sizes were limited to four. The framework supports arbitrary team sizes through the overdetermined least-squares fit, which averages measurement errors (with four robots, attraction precision improved from 0.036 to 0.014~m, Section~\ref{sec:ch6:hw_attraction}) and tolerates the loss of one robot. However, spreading more robots over a wider footprint requires the affine model to hold over a larger region of the curved field, where the curvature term of (\ref{eq:ch6:bound}) grows with the square of the formation size, and larger formations present practical challenges including communication bandwidth, computational complexity, and formation control stability. Teams beyond four, and the loss of a robot during a run, were not tested.

Fourth, the estimator assumes the field is well approximated by an affine model from the cluster out to the critical point (Section~\ref{sec:ch6:estimation}). On fields that are highly nonlinear at the scale of the formation, this local linear fit degrades, biasing the recovered critical point location. Section~\ref{sec:ch6:sim_attraction} probes this regime on a gyre cell and a curved saddle, but fields with stronger curvature, or with critical points closer together than one formation side, were not tested.

Finally, the computation is centralized (Section~\ref{sec:ch2:testbed}). The estimate needs only the positions and readings of one cycle in a common frame, so a robot receiving them from its neighbors could compute it. A dropout below three readings leaves the fit underdetermined, and latency mixes readings from different cycles into one fit. Onboard relative localization, inter-robot communication, and robustness to lower update rates and dropouts were not tested.

% New text. What the fitted robot model still misses (Chapter 6).
The fitted robot model of Section~\ref{sec:ch2:dynamics} closes part of the
gap between simulation and hardware that Chapter~\ref{ch:zeroth} found.
Chapter~\ref{ch:first_order} reports three differences that it leaves. The
live estimate led the cluster along its motion by 1.7 to 2.8~cm in small
orbits, a lead that the model omits. The hardware vortex ended 0.4 to
1.5~cm farther from the critical point than the matched simulations. The
orbits of the four-robot square fell short of the triangle's, which the
matched simulations did not predict. The shape of the low-speed response
is also unmeasured, since no run commanded less than about 0.06~m/s.

\subsection{Second-Order Estimation and Tracking}
\label{sec:ch9:second_order}

% Ported from: Draft_14.tex
Cluster space control is centralized and needs global state and reliable
high-bandwidth communication \cite{53}. Unlike distributed or
event-triggered consensus protocols \cite{68}, the estimator gathers all six
readings on one processor, so the claims here are for small clusters with
ample bandwidth.

Both trackers follow trenches of their own fields, which leave a
separatrix once the rotation on it rivals the strain
(Section~\ref{sec:ch7:beyond}). The $s_1$ tracker can reverse wherever its
eigenframe swaps, and its capture test can pass at a smooth crest. The
$D$ tracker's fitted Hessian omits the third-derivative terms that
dominate near the saddles, and the tracker is not objective.

The evaluation is kinematic simulation without flow advection, and the
ocean trial uses one record without added noise. In general neither
instantaneous Eulerian field is guaranteed to mark a Lagrangian material
transport structure \cite{51}. No head-to-head against \cite{27,63,28} or
\cite{6} was run. Their robots ride the flow and steer relative to it,
while ours command velocity over ground, so a fair comparison needs one
advection model for both, and their analyses assume a straddling start
that four of our seven starts do not provide. The no-advection premise is
exact on the planned Decabot hardware, where the field is only encoded
\cite{30}, but untested in a physical flow.

\section{Future Work}
\label{sec:ch9:future}

\subsection{Time-Varying Fields}
\label{sec:ch9:time_varying}

% Ported from: cwaight_systems_paper_with_revisions2.tex
The estimation framework uses only instantaneous measurements and carries no state between control cycles. This memoryless property means the estimator does not assume field stationarity, suggesting applicability to environments where the field itself varies with time, that is, where $\mathbf{v}(\mathbf{p}, t)$ changes shape or strength during operation. Validating the estimator on a genuinely time-varying field remains future work.

% New text
The Santa Barbara Channel trial of Section~\ref{sec:ch7:ocean} ran the
second-order trackers on a field that changes in time, but on one record
and with no analytic guarantee. On a time-periodic double gyre,
Section~\ref{sec:ch7:beyond} found that both trackers follow the separatrix
at the smaller amplitude, and that at the larger the separatrix outruns the
speed that either channel can deliver.

\subsection{Advection-Aware Control}
\label{sec:ch9:advection}

% Ported from: cwaight_systems_paper.tex
As noted in Section~\ref{sec:ch2:dynamics}, the robot dynamics (\ref{eq:ch2:continuous})--(\ref{eq:ch2:momentum}) respond only to commanded velocity. In a physical flow such as an ocean current, the field advects the platform, adding a disturbance term that the present controller does not compensate for. Extending the control architecture to account for this direct effect on robot motion is a priority for future work.

\subsection{Higher-Order Fits and Formation Reshaping}
\label{sec:ch9:higher_order}

% Ported from: Draft_14.tex
Future work targets three extensions. First, a fit from more than six
robots becomes least squares, which averages the measurement noise and
leaves a residual that, where the truncation error is small, estimates
that noise online. A cluster could then detect the high-noise regime in
which neither tracker is reliable. Second, the
$s_1$ tracker could take as its tangent the strain eigenvector nearest the
previous tangent. That rule does not rest on the along-trench slope and
could carry the direction through the eigenframe swaps behind both kinds
of reversal. Finally, a cubic
fit from ten robots off any common cubic curve would recover the third
derivatives on which the true $\mathbf{H}_D$ and the transverse curvature
of an $s_1$ trench depend.

\subsection{Choosing the Fit Order}
\label{sec:ch9:switching}

% New text. Built from the indicators of Table~\ref{tab:ch8:indicators}; none of
% this was built or tested.
The indicators of Table~\ref{tab:ch8:indicators} give a cluster the inputs
for choosing its own fit order. A cluster could hold a first-order fit while
$\lvert\hat{\omega}\rvert/2\hat{r}$ stays small and its type labels sit away
from their thresholds, and move to a second-order fit when a curve is worth
following. A fit from more robots than $n_k$ could signal, through its
residual, when the measurement noise is too large for either tracker. None
of this was built or tested.

\subsection{Three-Dimensional Fields}
\label{sec:ch9:3d}

% New text
The estimator of Section~\ref{sec:ch4:estimation} extends to three
dimensions by counting monomials in $x$, $y$, and $z$. A fit of order $k$
then has $\tfrac{1}{6}(k+1)(k+2)(k+3)$ coefficients per component. A
first-order fit needs four robots and fails when they are coplanar. A
second-order fit needs ten robots and fails when they lie on a common
quadric surface. The sensing and formation control of a
three-dimensional cluster are outside this dissertation.

\subsection{Additional Topological Features}
\label{sec:ch9:features}

% Ported from: cwaight_systems_paper_with_revisions2.tex
Beyond critical point detection and orbital control, vector fields contain other topological features of practical interest. Heteroclinic orbits connecting different critical points through their manifolds offer transit routes through complex fields. The three-robot fit returns a spatially constant Jacobian, so detecting these features calls for a second-order fit, which six robots that do not lie on a common conic can provide.

% New text. Chapter 7 followed one heteroclinic connection on the steady flow.
On the steady double gyre the separatrix between the two saddles is a
heteroclinic connection, and the $D$ tracker travels it
(Section~\ref{sec:ch8:nesting}). The bent gyre and the Bickley jet of
Section~\ref{sec:ch7:beyond} show that a trench of a fitted surrogate
follows such a connection only where strain outweighs rotation, so
following heteroclinic connections in general fields remains future work.

\subsection{Testbed Extensions}
\label{sec:ch9:testbed}

% Ported from: IDETC 2025 (DETC2025-167604), Sec. 6, transcribed from the published PDF
The supervised learning calibration technique can be extended to
greyscale sensors for encoding and decoding scalar field monochrome color
maps, creating a unified platform for testing both vector and scalar
field navigation paradigms simultaneously. The modular architecture can
also be enhanced to facilitate controller substitution, enabling
comparative studies between cluster space controllers and alternative
approaches such as swarm-based methods \cite{44}. This modularity
would allow researchers to isolate the effects of different control
strategies while maintaining consistent environmental conditions.
Outdoor testing with autonomous surface vessels \cite{2} and UAVs for
environmental monitoring \cite{69} would examine transferability to
larger scales and real-world deployment challenges.

% New text. The IDETC digital-twin and controller-substitution plans, now done in part.
Part of this has been done. Section~\ref{sec:ch2:dynamics} fits the dynamics
of the Decabots for simulation, and Section~\ref{sec:ch6:comparison}
substituted the vector-to-scalar primitive for the Jacobian method with the
cluster and robot layers held fixed. One further test would settle the dead
zone of Section~\ref{sec:ch2:dynamics}: a single robot stepped through
commands from 0.01 to 0.08~m/s, from rest and from motion.

% -------------------------------------------------------------------
% Parked research note, kept by the author's decision (commit ad76bbf,
% 2026-07-02). Not rendered: it describes the earlier D = 0 contour
% tracker, which this dissertation does not present.
% -------------------------------------------------------------------
%% \subsection{Okubo-Weiss Corner Branch Selection (Deferred 2026-07-02)}
%% \label{sec:ch9:ow_corner_branch}
%%
%% % Research note carried over unchanged from phd_thesis_template.tex's
%% % prior Chapter 5 (git history, this file). Relocated here because it
%% % is a specific, evaluable future-work item, not a general limitation.
%% % Sources: experiments/ow_clean_runs.py and
%% % experiments/outputs/ow_clean/ under
%% % trunk/Python_Simulations/Vector_Fields/VF_Robot/, git commits 7dc01a8
%% % and c19732f, and plan.md (repo root) Phase 3 and Phase 7 entries.
%%
%% \textbf{Context.} The Okubo-Weiss boundary tracker of the second paper
%% drives $\hat{D} = \det(\hat{\mathbf{J}})$ to zero and drifts along the
%% level-set tangent $\mathrm{perp}(\nabla\hat{D})$, sign-aligned with the
%% ambient flow. On the steady double gyre the $D = 0$ set is not a pair of
%% closed diamonds but a network of straight lines,
%% $x_f \pm y_f \in \tfrac{1}{2} + \mathbb{Z}$, which cross at the diamond
%% corners where $\nabla D = \mathbf{0}$. Noise-free experiments (four
%% starts) showed lock-on within 13 to 16 control steps and mean $|D|$
%% between 0.003 and 0.007 on the boundary, 0.3 to 0.8 percent of the well
%% depth $\pi^4 A^2 = 0.974$, but the tangent follows its line straight
%% through every corner crossing: the gradient-floor fallback carries the
%% team across the degenerate point and the Newton action resumes on the
%% continuing line. Circulation around a single vortex core therefore never
%% occurs with the published law. The paper presents this straight-through
%% behavior honestly; the wording of Problem 2 was changed from ``circulate
%% along'' to ``travel along'' the boundary.
%%
%% \textbf{The deferred idea.} At an X-crossing the field is locally
%% $D(\mathbf{p}^* + \mathbf{t}) \approx \tfrac{1}{2}\mathbf{t}^\top
%% \mathbf{H}_D \mathbf{t}$ with $\mathbf{H}_D$ indefinite, so the zero set
%% consists of the two asymptote directions $\mathbf{t}$ satisfying
%% $\mathbf{t}^\top \mathbf{H}_D \mathbf{t} = 0$. In the eigenbasis of
%% $\mathbf{H}_D$ with $\lambda_1 < 0 < \lambda_2$ these lie at angles
%% $\pm\arctan\sqrt{-\lambda_1/\lambda_2}$ from the $\lambda_2$
%% eigenvector. A corner rule can therefore be computed onboard from
%% $\hat{\mathbf{H}}_D$ alone: inside the gradient-floor ball, evaluate
%% both estimated asymptotes, discard the one parallel to the incoming
%% tangent, and exit along the turning branch, sign-chosen by the flow.
%% Selecting the turning branch at every corner closes the circuit around
%% one diamond; always selecting the continuing branch reproduces the
%% published straight-through behavior, so the rule is a second
%% one-parameter behavior switch, mirroring the park-versus-traverse knob
%% of the separatrix tracker.
%%
%% \textbf{Known risk.} The $\hat{\mathbf{H}}_D$ estimate carries a
%% formation-radius-independent structural bias (third-derivative terms
%% absent from the quadratic model; see the estimator accuracy results of
%% the second paper). Eigen-directions survive that bias to within a few
%% tenths of a degree on the separatrix, but the accuracy of the estimated
%% asymptote directions at the corners specifically is untested, and near
%% the corners the eigenvalues themselves are strongly biased. This is the
%% first thing to measure.
%%
%% \textbf{Validation plan when revisited.} Implement as a new primitive
%% (working name \textit{logic\_g\_fable\_step}) alongside the untouched
%% \textit{logic\_g\_newton\_contour\_pentagon}, following the pattern of
%% \textit{separatrix\_logic\_fable\_step} (bit-exact equivalence to the
%% parent law when the knob is off, then the new branch enabled).
%% Head-to-head at all eight diamond corners plus the domain-center
%% crossing: success is loop closure around one diamond; metrics are
%% circumnavigation time and mean $|D|$ per circuit, against the
%% straight-through baseline of \textit{ow\_clean\_runs.py}.

% ===================================================================
% CHAPTER 10: CONCLUSION
% ===================================================================
\chapter{Conclusion}
\label{ch:conclusion}

\section{Summary of Contributions}
\label{sec:ch10:summary}

% Ported from: IDETC 2025 (DETC2025-167604), Sec. 7, transcribed from the published PDF
The testbed of Chapter~\ref{ch:tools} overcame previous testbed
limitations through HSV color space representation and neural network
calibration, achieving $R^2$ accuracies of 0.96 for direction and 0.91 for
magnitude while maintaining precise formation control of the robot
cluster. On it, the vector-sum and vector-to-scalar primitives ran in
continuous motion, confirming behaviors previously demonstrated only in
simulation.

% Ported from: cwaight_systems_paper_with_revisions2.tex
Chapter~\ref{ch:first_order} presented a framework for multirobot navigation in vector fields containing critical points. A Jacobian-based approach uses instantaneous measurements from three or more robots to identify critical point locations without prior field knowledge, supporting attraction and orbital control laws built on these estimates. Simulation across six canonical field types confirmed convergence of these control primitives. On curved fields the sign of $\det\hat{\mathbf{J}}$ on the first cycle decided which type of critical point the cluster reached. In 459 hardware runs, the cluster converged within 0.05~m in 99\% of saddle and 91\% of vortex trials with three robots and in every vortex trial with four, with mean position errors of 0.034~m or less. Orbital control experiments maintained commanded radii within a quarter of the radius from 0.20~m upward with three robots in both fields, and below that held bounded orbits that did not track the commanded radius; at 0.10~m the vortex orbit settled onto a limit cycle sized by the robots' delay and lag. Simulations matched to each run reproduced the saddle's attraction error and about half of the orbit fluctuation.

% Ported from: Draft_13.tex
Chapter~\ref{ch:second_order} showed that six robots sampling at once can fit a local
quadratic model of a planar flow, which gives the velocity gradient as a
field over the formation. Two surrogate fields read from that fit, the
Jacobian determinant $D$ and the smaller strain eigenvalue $s_1$, mark
coherent structures of the flow as trenches. On each field we built an
adaptive navigation primitive that acquires and rides those trenches from
instantaneous measurements, with no integration horizon and no
pre-straddling. Without noise, both acquire and traverse the
double-gyre separatrix from every start and hold the downstream saddle. Measured
by the noise level at which retention falls to 50\%, the $D$ tracker
withstands more measurement and position noise than the $s_1$ tracker from
every start tested. The $s_1$ tracker is objective past its
one seed step, and it alone holds the same path under a rotating observer.
A cluster with an absolute heading reference, such as a compass or GNSS,
should therefore carry the $D$ tracker, and one whose flow estimate
carries its own rotation should carry the $s_1$ tracker.

% New text. Primitive library, moved from Chapter 8 on 2026-09-30.
Table~\ref{tab:ch10:library} collects every primitive in this
dissertation with its fit order, the smallest cluster it was run with,
the quantity it steers on, its stability result, and how it was verified.

\begin{table}[htbp]
\centering
\footnotesize
\caption{Primitive library. Stability entries are as stated in the
chapter that defines each primitive.}
\label{tab:ch10:library}
\setlength{\tabcolsep}{3pt}
\begin{tabular}{>{\raggedright\arraybackslash}p{2.0cm} c c >{\raggedright\arraybackslash}p{2.7cm} >{\raggedright\arraybackslash}p{3.4cm} >{\raggedright\arraybackslash}p{3.0cm}}
\hline
Primitive & Order & Robots & Steers on & Stability & Verification \\
\hline
Vector-sum & 0 & 3 & mean sensed vector & none stated & hardware, fixed and sinking vortex, 10 runs each \\
Vector-to-scalar & 0 & 3 & slope of the sensed magnitude & none stated & hardware, fixed vortex, 10 runs plus randomized starts \\
Attraction & 1 & 3 & $\hat{\mathbf{p}}^* = -\hat{\mathbf{J}}^{-1}\hat{\mathbf{h}}$ & exponential convergence to a stationary estimate, $\tau = 1/k$ & simulation, 11 fields; hardware, 239 trials \\
Orbital & 1 & 3 & $\hat{\mathbf{p}}^*$ with radial and tangential terms & radial error $\dot{e}_r = -k_r e_r$ & simulation, 11 fields; hardware, 220 runs \\
$D$ tracker & 2 & 6 & $\hat{D}$ and $\hat{\mathbf{H}}_{D,0}$ & transverse practical stability, $\limsup|n| \leq \delta/a_\perp$ & simulation, double gyre and Santa Barbara Channel \\
$s_1$ tracker & 2 & 6 & $\hat{s}_1$ and $\nabla\hat{s}_1$ & same transverse bound, frame equivariant past its seed step & simulation, double gyre and Santa Barbara Channel \\
\hline
\end{tabular}
\end{table}

\section{Final Thoughts}
\label{sec:ch10:final}

% New text
The primitives in this dissertation are admittedly simple, reactive, and
minimal. Each uses the smallest cluster its fit order allows and only the
measurements of the current instant. Even in this form, however, they
show that the order of a cooperative fit sets which features of a flow a
cluster can reach, from a heading, to a critical point, to a separatrix.
Ongoing and future work targets a) hardware trials of the second-order
trackers on the Decabot testbed, b) control that accounts for a flow that
advects the robots, and c) a third-order fit from ten robots.


% ===================================================================
% REFERENCES
% ===================================================================
\clearpage
\singlespacing
\addcontentsline{toc}{chapter}{References}

\begin{thebibliography}{99}

% Generated by sources/assemble/build_bib.py. Order of first citation.

\bibitem{1} J. Cortés and M. Egerstedt, ``Coordinated control of multi-robot systems: A survey,'' \textit{SICE Journal of Control, Measurement, and System Integration}, vol. 10, no. 6, pp. 495--503, 2017.

\bibitem{2} T. Adamek, C. A. Kitts, and I. Mas, ``Gradient-based cluster space navigation for autonomous surface vessels,'' \textit{IEEE/ASME Transactions on Mechatronics}, vol. 20, no. 2, pp. 506--518, 2014.

\bibitem{3} C. A. Kitts, R. T. McDonald, and M. A. Neumann, ``Adaptive navigation control primitives for multirobot clusters: Extrema finding, contour following, ridge/trench following, and saddle point station keeping,'' \textit{IEEE Access}, vol. 6, pp. 17625--17642, 2018.

\bibitem{4} R. T. McDonald, C. A. Kitts, and M. A. Neumann, ``Navigation of scalar fronts with multirobot clusters in simulation and experiment,'' \textit{IEEE Systems Journal}, vol. 14, no. 3, pp. 3755--3766, 2020.

\bibitem{5} L. Briñón-Arranz, A. Renzaglia, and L. Schenato, ``Multirobot symmetric formations for gradient and hessian estimation with application to source seeking,'' \textit{IEEE Transactions on Robotics}, vol. 35, no. 3, pp. 782--789, 2019.

\bibitem{6} A. Okubo, ``Horizontal dispersion of floatable particles in the vicinity of velocity singularities such as convergences,'' \textit{Deep-Sea Res.}, vol. 17, no. 3, pp. 445--454, 1970.

\bibitem{7} P. Keramea, K. Spanoudaki, G. Zodiatis, G. Gikas, and G. Sylaios, ``Oil spill modeling: A critical review on current trends, perspectives, and challenges,'' \textit{Journal of Marine Science and Engineering}, vol. 9, no. 2, p. 181, 2021.

\bibitem{8} D. Kularatne and M. A. Hsieh, ``Tracking attracting manifolds in flows,'' \textit{Autonomous Robots}, vol. 41, no. 8, pp. 1575--1588, 2017.

\bibitem{9} X. Tricoche, G. Scheuermann, and H. Hagen, ``Continuous topology simplification of planar vector fields,'' in \textit{Proc. IEEE Visualization}, pp. 159--166, 2001.

\bibitem{10} W. G. Pichel et al., ``Marine debris collects within the North Pacific subtropical convergence zone,'' \textit{Marine Pollution Bulletin}, vol. 54, no. 8, pp. 1207--1211, 2007.

\bibitem{11} M. Mateus, R. Canelas, L. Pinto, and N. Vaz, ``When tragedy strikes: potential contributions from ocean observation to search and rescue operations after drowning accidents,'' \textit{Frontiers in Marine Science}, vol. 7, p. 55, 2020.

\bibitem{12} J. Wang, Y. Lin, R. Liu, and J. Fu, ``Odor source localization of multi-robots with swarm intelligence algorithms: A review,'' \textit{Frontiers in Neurorobotics}, vol. 16, p. 949888, 2022.

\bibitem{13} C. Dong et al., ``Oceanic mesoscale eddies,'' \textit{Ocean-Land-Atmosphere Research}, vol. 4, p. 0081, 2025.

\bibitem{14} Z. Li, X. Wang, J. Hu, F. P. Andutta, and Z. Liu, ``A study on an anticyclonic-cyclonic eddy pair off Fraser Island, Australia,'' \textit{Frontiers in Marine Science}, vol. 7, p. 594358, 2020.

\bibitem{15} R. Deogharia, H. Gupta, S. Sil, A. Gangopadhyay, and A. Shee, ``On the evidence of helico-spiralling recirculation within coherent cores of eddies using Lagrangian approach,'' \textit{Scientific Reports}, vol. 14, no. 1, p. 11014, 2024.

\bibitem{16} P. Ögren, E. Fiorelli, and N. E. Leonard, ``Cooperative control of mobile sensor networks: Adaptive gradient climbing in a distributed environment,'' \textit{IEEE Transactions on Automatic Control}, vol. 49, no. 8, pp. 1292--1302, 2004.

\bibitem{17} R. T. McDonald, C. A. Kitts, and M. A. Neumann, ``Experimental implementation and verification of scalar field ridge, trench, and saddle point maneuvers using multirobot adaptive navigation,'' \textit{IEEE Access}, vol. 7, pp. 62950--62961, 2019.

\bibitem{18} F. Pagano, N. De Carli, E. Restrepo, A. Marino, and P. Robuffo Giordano, ``Distributed multi-robot active-sensing of a diffusive source,'' \textit{IEEE Robotics and Automation Letters}, vol. 10, no. 10, pp. 10807--10814, 2025.

\bibitem{19} S. Li, R. Kong, and Y. Guo, ``Cooperative distributed source seeking by multiple robots: Algorithms and experiments,'' \textit{IEEE/ASME Transactions on Mechatronics}, vol. 19, no. 6, pp. 1810--1820, 2014.

\bibitem{20} Z. Deng and J. Luo, ``Fully distributed algorithms for constrained nonsmooth optimization problems of general linear multiagent systems and their application,'' \textit{IEEE Transactions on Automatic Control}, vol. 69, no. 2, pp. 1377--1384, 2023.

\bibitem{21} W. Li, J. A. Farrell, S. Pang, and R. M. Arrieta, ``Moth-inspired chemical plume tracing on an autonomous underwater vehicle,'' \textit{IEEE Transactions on Robotics}, vol. 22, no. 2, pp. 292--307, 2006.

\bibitem{22} S. R. Ramp et al., ``Preparing to predict: the second autonomous ocean sampling network (AOSN-II) experiment in the Monterey Bay,'' \textit{Deep Sea Research Part II: Topical Studies in Oceanography}, vol. 56, no. 3-5, pp. 68--86, 2009.

\bibitem{23} J. Das et al., ``Coordinated sampling of dynamic oceanographic features with underwater vehicles and drifters,'' \textit{The International Journal of Robotics Research}, vol. 31, no. 5, pp. 626--646, 2012.

\bibitem{24} J. Wang et al., ``Dynamic plume tracking by cooperative robots,'' \textit{IEEE/ASME Transactions on Mechatronics}, vol. 24, no. 2, pp. 609--620, 2019.

\bibitem{25} D. Panagou, ``Motion planning and collision avoidance using navigation vector fields,'' in \textit{Proc. IEEE International Conference on Robotics and Automation}, 2014, pp. 2513--2518.

\bibitem{26} R. Cooper et al., ``Initial Study of Multirobot Adaptive Navigation for Exploring Environmental Vector Fields,'' in \textit{Proc. IEEE Aerospace Conference}, 2019.

\bibitem{27} K. Garg and D. Panagou, ``Finite-time estimation and control for multi-aircraft systems under wind and dynamic obstacles,'' \textit{Journal of Guidance, Control, and Dynamics}, vol. 42, no. 7, pp. 1489--1505, 2019.

\bibitem{28} C. Waight and C. Kitts, ``A functional indoor testbed for multirobot adaptive navigation in vector field environments,'' in \textit{Proc. ASME Int. Design Engineering Technical Conf. Computers and Information in Engineering Conf.}, vol. 89251, 2025.

\bibitem{29} M. Michini, M. A. Hsieh, E. Forgoston, and I. B. Schwartz, ``Robotic tracking of coherent structures in flows,'' \textit{IEEE Transactions on Robotics}, vol. 30, no. 3, pp. 593--603, 2014.

\bibitem{30} D. R. Nelson, D. B. Barber, T. W. McLain, and R. W. Beard, ``Vector field path following for miniature air vehicles,'' \textit{IEEE Transactions on Robotics}, vol. 23, no. 3, pp. 519--529, 2007.

\bibitem{31} W. Yao, H. G. de Marina, B. Lin, and M. Cao, ``Singularity-free guiding vector field for robot navigation,'' \textit{IEEE Transactions on Robotics}, vol. 37, no. 4, pp. 1206--1221, 2021.

\bibitem{32} X. Wang et al., ``Adaptive vector field guidance without a priori knowledge of course dynamics and wind,'' \textit{IEEE/ASME Transactions on Mechatronics}, vol. 27, no. 6, pp. 4597--4607, 2022.

\bibitem{33} Y. Zhang, M. A. Godin, J. G. Bellingham, and J. P. Ryan, ``Using an autonomous underwater vehicle to track a coastal upwelling front,'' \textit{IEEE Journal of Oceanic Engineering}, vol. 37, no. 3, pp. 338--347, 2012.

\bibitem{34} T. Inanc, S. C. Shadden, and J. E. Marsden, ``Optimal trajectory generation in ocean flows,'' in \textit{Proc. American Control Conference}, pp. 674--679, 2005.

\bibitem{35} R. N. Smith, M. Schwager, S. L. Smith, B. H. Jones, D. Rus, and G. S. Sukhatme, ``Persistent ocean monitoring with underwater gliders: Adapting sampling resolution,'' \textit{Journal of Field Robotics}, vol. 28, no. 5, pp. 714--741, 2011.

\bibitem{36} G. Knizhnik, P. Li, X. Yu, and M. A. Hsieh, ``Flow-based control of marine robots in gyre-like environments,'' in \textit{Proc. IEEE International Conference on Robotics and Automation (ICRA)}, pp. 3047--3053, 2022.

\bibitem{37} A. Matveev et al., ``Non-gradient based tracking of environmental field isolines in a changing medium by Dubins-vehicle type robots,'' \textit{IEEE Transactions on Automatic Control}, 2023.

\bibitem{38} A. Mokhtarian et al., ``A survey on small-scale testbeds for connected and automated vehicles and robot swarms: A guide for creating a new testbed,'' \textit{IEEE Robotics \& Automation Magazine}, 2024.

\bibitem{39} Z. Tian et al., ``ICAT: An indoor connected and autonomous testbed for vehicle computing,'' in \textit{Proc. IEEE International Conference on Mobility, Operations, Services and Technologies (MOST)}, 2024.

\bibitem{40} S. Schäfer and B. Alrifaee, ``From small-scale to full-scale: Assessing the potential for transferability of experimental results in small-scale CAV testbeds,'' in \textit{Proc. IEEE International Conference on Vehicular Electronics and Safety (ICVES)}, 2024.

\bibitem{41} S. Tomer et al., ``A low-cost indoor testbed for multirobot adaptive navigation research,'' in \textit{Proc. IEEE Aerospace Conference}, pp. 1--12, 2018.

\bibitem{42} F.-J. Mañas-Álvarez et al., ``Robotic park: Multi-agent platform for teaching control and robotics,'' \textit{IEEE Access}, vol. 11, pp. 34899--34911, 2023.

\bibitem{43} S. T. Hart and C. A. Kitts, ``Unifying control architecture for reactive particle swarms,'' \textit{IEEE/ASME Transactions on Mechatronics}, vol. 28, no. 2, pp. 873--883, 2022.

\bibitem{44} D. Rajabhandharaks, R. T. McDonald, M. A. Neumann, and C. A. Kitts, ``Indoor testbed for vector field multirobot adaptive navigation: Feasibility study,'' in \textit{Proc. ASME Int. Design Engineering Technical Conf. Computers and Information in Engineering Conf.}, vol. 59292, p. V009T12A043, 2019.

\bibitem{45} M. Woolsey, C. A. Kitts, and M. A. Neumann, ``Gradient-based adaptive navigation of multiple diving autonomous surface vehicles,'' in \textit{Proc. OCEANS 2022, Hampton Roads}, 2022.

\bibitem{46} R. K. Lee et al., ``Multiple UAV adaptive navigation for three-dimensional scalar fields,'' \textit{IEEE Access}, vol. 9, pp. 122626--122654, 2021.

\bibitem{47} Z. Zhou, J. Liu, and J. Yu, ``A survey of underwater multi-robot systems,'' \textit{IEEE/CAA Journal of Automatica Sinica}, vol. 9, no. 1, pp. 1--18, 2021.

\bibitem{48} J. Weiss, ``The dynamics of enstrophy transfer in two-dimensional hydrodynamics,'' \textit{Physica D}, vol. 48, no. 2-3, pp. 273--294, 1991.

\bibitem{49} M. Serra and G. Haller, ``Objective Eulerian coherent structures,'' \textit{Chaos}, vol. 26, no. 5, p. 053110, 2016.

\bibitem{50} M. Serra et al., ``Search and rescue at sea aided by hidden flow structures,'' \textit{Nat. Commun.}, vol. 11, no. 1, p. 2525, 2020.

\bibitem{51} S. C. Shadden, F. Lekien, and J. E. Marsden, ``Definition and properties of Lagrangian coherent structures from finite-time Lyapunov exponents in two-dimensional aperiodic flows,'' \textit{Physica D}, vol. 212, no. 3-4, pp. 271--304, 2005.

\bibitem{52} C. A. Kitts and I. Mas, ``Cluster space specification and control of mobile multirobot systems,'' \textit{IEEE/ASME Transactions on Mechatronics}, vol. 14, no. 2, pp. 207--218, 2009.

\bibitem{53} I. Mas and C. Kitts, ``Obstacle avoidance policies for cluster space control of nonholonomic multirobot systems,'' \textit{IEEE/ASME Transactions on Mechatronics}, vol. 17, no. 6, pp. 1068--1079, 2012.

\bibitem{54} Y. Jiao, H. Hang, J. Merel, and E. Kanso, ``Sensing flow gradients is necessary for learning autonomous underwater navigation,'' \textit{Nature Communications}, vol. 16, no. 1, p. 3044, 2025.

\bibitem{55} T. Peacock and G. Haller, ``Lagrangian coherent structures: The hidden skeleton of fluid flows,'' \textit{Phys. Today}, vol. 66, no. 2, pp. 41--47, 2013.

\bibitem{56} G. Haller and G. Yuan, ``Lagrangian coherent structures and mixing in two-dimensional turbulence,'' \textit{Physica D}, vol. 147, no. 3-4, pp. 352--370, 2000.

\bibitem{57} J. C. R. Hunt, A. A. Wray, and P. Moin, ``Eddies, streams, and convergence zones in turbulent flows,'' in \textit{Proc. Summer Program, Center Turbul. Res.}, Rep. CTR-S88, pp. 193--208, 1988.

\bibitem{58} B.-L. Hua and P. Klein, ``An exact criterion for the stirring properties of nearly two-dimensional turbulence,'' \textit{Physica D}, vol. 113, no. 1, pp. 98--110, 1998.

\bibitem{59} G. Haller, ``Lagrangian coherent structures,'' \textit{Annu. Rev. Fluid Mech.}, vol. 47, pp. 137--162, 2015.

\bibitem{60} P. J. Nolan, M. Serra, and S. D. Ross, ``Finite-time Lyapunov exponents in the instantaneous limit and material transport,'' \textit{Nonlinear Dyn.}, vol. 100, no. 4, pp. 3825--3852, 2020.

\bibitem{61} T. D. Harms, S. L. Brunton, and B. J. McKeon, ``Lagrangian gradient regression for the detection of coherent structures from sparse trajectory data,'' \textit{R. Soc. Open Sci.}, vol. 11, no. 9, p. 240586, 2024.

\bibitem{62} S. McCammon et al., ``Ocean front detection and tracking using a team of heterogeneous marine vehicles,'' \textit{J. Field Robot.}, vol. 38, no. 6, pp. 854--881, 2021.

\bibitem{63} M. Michini, H. Rastgoftar, M. A. Hsieh, and S. Jayasuriya, ``Distributed formation control for collaborative tracking of manifolds in flows,'' in \textit{Proc. Amer. Control Conf.}, pp. 3874--3880, 2014.

\bibitem{64} M. Michini, M. A. Hsieh, E. Forgoston, and I. B. Schwartz, ``Experimental validation of robotic manifold tracking in gyre-like flows,'' in \textit{Proc. IEEE/RSJ Int. Conf. Intell. Robots Syst.}, pp. 2306--2311, 2014.

\bibitem{65} C. Waight and C. A. Kitts, ``Adaptive navigation of multirobot systems to critical points in 2D vector fields,'' submitted for publication.

\bibitem{66} D. Eberly et al., ``Ridges for image analysis,'' \textit{J. Math. Imaging Vis.}, vol. 4, no. 4, pp. 353--373, 1994.

\bibitem{67} Coastal Observing Research and Development Center, Scripps Inst. Oceanogr., ``Near-real time surface ocean velocity, U.S. West Coast, 2 km resolution,'' HFRNet, NOAA NCEI Accession gov.noaa.nodc:IOOS-HFRadarRTVector, 2012. [Online]. Available: \url{https://www.ncei.noaa.gov/thredds-ocean/fileServer/ioos/hfradar/rtv/2012/201205/USWC/}. Accessed: Jul. 1, 2026.

\bibitem{68} E. Terrill et al., ``Data management and real-time distribution in the HF-Radar National Network,'' in \textit{Proc. MTS/IEEE OCEANS}, Boston, MA, USA, Sep. 2006.

\bibitem{69} S. Zhang, J. Liu, W. Wang, and Z. Zhang, ``Bipartite leader-following consensus of nonlinear switched multiagent systems under model-depended dynamic event-triggered control,'' \textit{IEEE Syst. J.}, vol. 18, no. 2, pp. 1044--1055, 2024.


\end{thebibliography}

% ===================================================================
% APPENDICES
% ===================================================================
\appendix

\chapter{Canonical Field Definitions}
\label{app:fields}

\section{Six Canonical Critical Point Fields}
\label{sec:appA:canonical}

% Ported from: cwaight_systems_paper.tex
Six vector field environments centered at the origin were used to evaluate the first-order primitives of Chapter~\ref{ch:first_order}. Let $r = \sqrt{x^2 + y^2}$ and $\theta = \text{atan2}(y, x)$. The base fields are the vortex (pure rotational flow), the sink (radial inward flow) whose negation $\mathbf{v}_{\text{source}} = -\mathbf{v}_{\text{sink}}$ is the source (radial outward flow), and the saddle, an axis-aligned saddle rotated by $45^\circ$ so that its stable and unstable directions lie along the diagonals $y = x$ and $y = -x$:
\begin{equation}\label{eq:appA:14}
\mathbf{v}_{\text{vortex}} = \begin{bmatrix} -r \sin(\theta) \\ r \cos(\theta) \end{bmatrix}, \;\;
\mathbf{v}_{\text{sink}} = \begin{bmatrix} -x \\ -y \end{bmatrix}, \;\;
\mathbf{v}_{\text{saddle}} = \begin{bmatrix} -y \\ -x \end{bmatrix}
\end{equation}

% Ported from: cwaight_systems_paper.tex
The sinking vortex (spiraling inward) and spewing vortex (spiraling outward) are linear combinations of the base fields:
\begin{equation}\label{eq:appA:15}
\begin{aligned}
\mathbf{v}_{\text{sink-vortex}} &= 0.4\,\mathbf{v}_{\text{vortex}} + 0.15\,\mathbf{v}_{\text{sink}} \\
\mathbf{v}_{\text{spew-vortex}} &= -0.4\,\mathbf{v}_{\text{vortex}} - 0.15\,\mathbf{v}_{\text{sink}}
\end{aligned}
\end{equation}

\section{Testbed Vortex Fields}
\label{sec:appA:testbed}

% New text
The fixed and sinking vortex maps of Chapter~\ref{ch:zeroth} are given by
(\ref{eq:ch5:rtheta})--(\ref{eq:ch5:sinking}). The fixed vortex there is
the linear clockwise rotation $(y - c_y,\, -(x - c_x))$, the vortex of
(\ref{eq:appA:14}) with the opposite sense. Its sinking component falls
off as $1/r$, so it differs from the linear sinking vortex of
(\ref{eq:appA:15}) and is singular at the center apart from the
$10^{-6}$ regularization.

\section{Steady Double Gyre}
\label{sec:appA:dg}

% New text
The steady double gyre of Chapter~\ref{ch:second_order} is defined in
(\ref{eq:ch3:dg_u}), with its determinant, Hessian, and smaller strain
eigenvalue in closed form in Section~\ref{sec:ch3:dg_example}.

\begin{table}[htbp]
\centering
\caption{Fields used in this dissertation.}
\label{tab:appA:summary}
\begin{tabular}{>{\raggedright\arraybackslash}p{5.2cm} >{\raggedright\arraybackslash}p{5.2cm} c}
\hline
Field & Critical points & Chapter \\
\hline
Vortex, sink, source, saddle & center, stable node, unstable node, saddle & \ref{ch:first_order} \\
Sinking and spewing vortex (linear) & stable and unstable spiral & \ref{ch:first_order} \\
Fixed vortex (testbed) & center & \ref{ch:zeroth} \\
Sinking vortex (testbed, $1/r$ sink) & singular sink at the center & \ref{ch:zeroth} \\
Steady double gyre & two saddles, two centers & \ref{ch:second_order} \\
Santa Barbara Channel record & measured, time-varying & \ref{ch:second_order} \\
\hline
\end{tabular}
\end{table}

\chapter{Three-Robot Cluster Kinematics}
\label{app:kin3}

% Ported from: cwaight_systems_paper.tex
\section{Forward Kinematics}

The forward kinematics transforms robot positions $(x_1, y_1)$, $(x_2, y_2)$, $(x_3, y_3)$ to the SAS parameters and centroid coordinates. The centroid position is:
\begin{equation}\label{eq:appB:20}
x_c = \frac{x_1 + x_2 + x_3}{3}, \quad y_c = \frac{y_1 + y_2 + y_3}{3}
\end{equation}
The side lengths, with $\mathbf{p}_i = (x_i, y_i)$ the position of robot $i$, are:
\begin{equation}\label{eq:appB:21}
p = \lVert \mathbf{p}_2 - \mathbf{p}_1 \rVert, \quad q = \lVert \mathbf{p}_3 - \mathbf{p}_2 \rVert
\end{equation}
The interior angle at robot 2 is:
\begin{equation}\label{eq:appB:23}
\beta = \arccos\left(\frac{(\mathbf{p}_1 - \mathbf{p}_2) \cdot (\mathbf{p}_3 - \mathbf{p}_2)}{pq}\right)
\end{equation}

The orientation angle is:
\begin{equation}\label{eq:appB:24}
\theta_c = \text{atan2}(y_1 - y_2, x_1 - x_2)
\end{equation}

\section{Inverse Kinematics}

Given the SAS parameters $(p, \beta, q)$ and centroid parameters $(x_c, y_c, \theta_c)$, the robot positions are recovered by constructing the triangle in a local frame and then applying rotation and translation.

In the local frame with robot 2 at the origin, the triangle is constructed as:
\begin{equation}\label{eq:appB:25}
\begin{aligned}
(x_1, y_1)^{\text{local}} &= (p,\, 0), \qquad (x_2, y_2)^{\text{local}} = (0,\, 0), \\
(x_3, y_3)^{\text{local}} &= (q\cos\beta,\, q\sin\beta)
\end{aligned}
\end{equation}

This construction assumes robot 3 lies counterclockwise from the robot 1--robot 2 edge; since $\beta$ from (\ref{eq:appB:23}) lies in $(0, \pi)$, the mirror configuration (robot 3 clockwise of the edge) maps to the same SAS parameters.

The local centroid is:
\begin{equation}\label{eq:appB:28}
x_c^{\text{local}} = \frac{x_1^{\text{local}} + x_2^{\text{local}} + x_3^{\text{local}}}{3}, \quad y_c^{\text{local}} = \frac{y_1^{\text{local}} + y_2^{\text{local}} + y_3^{\text{local}}}{3}
\end{equation}

After centering the triangle at the local origin and applying rotation $\theta_c$ and translation to $(x_c, y_c)$, the global positions become:
\begin{equation}\label{eq:appB:29}
\begin{bmatrix} x_i \\ y_i \end{bmatrix} = \begin{bmatrix} \cos\theta_c & -\sin\theta_c \\ \sin\theta_c & \cos\theta_c \end{bmatrix} \begin{bmatrix} x_i^{\text{local}} - x_c^{\text{local}} \\ y_i^{\text{local}} - y_c^{\text{local}} \end{bmatrix} + \begin{bmatrix} x_c \\ y_c \end{bmatrix}
\end{equation}
for $i \in \{1, 2, 3\}$.

\section{Inverse Jacobian}

The inverse Jacobian maps shape parameter velocities to robot velocities:
\begin{equation}\label{eq:appB:30}
\begin{bmatrix}
\delta x_1 \\ \delta y_1 \\ \delta x_2 \\ \delta y_2 \\ \delta x_3 \\ \delta y_3
\end{bmatrix} =
\mathbf{J}_c^{-1}
\begin{bmatrix}
\delta x_c \\ \delta y_c \\ \delta \theta_c \\ \delta p \\ \delta \beta \\ \delta q
\end{bmatrix}
\end{equation}

where $\mathbf{J}_c^{-1}$ is a $6 \times 6$ matrix containing partial derivatives $\frac{\partial x_i}{\partial \xi_j}$ and $\frac{\partial y_i}{\partial \xi_j}$ for robot positions and shape parameters $\xi_j \in \{x_c, y_c, \theta_c, p, \beta, q\}$. This matrix is computed analytically by taking partial derivatives of the inverse kinematics equations (\ref{eq:appB:29}) with respect to each shape parameter. The explicit entries are not listed here. The simulation software computes them by forward finite differences of the inverse kinematics (\texttt{compute\_inverse\_jacobian} in \texttt{src/control/kinematics.py}), and the \texttt{cluster\_builder} tool generates them in closed symbolic form (\texttt{clusterbuilder.py} with the \texttt{--symbolic} option).

\chapter{Six-Robot Cluster Kinematics}
\label{app:kin6}

% New text. Checked against src/control/pentagon_kinematics.py
% (VF_Robot), which implements this parameterization.
The six-robot cluster of Chapter~\ref{ch:second_order} is a cluster of
clusters. The robots form three pairs, and the pair midpoints form a
triangle described with the same SAS variables as the three-robot cluster
of Appendix~\ref{app:kin3}. The estimator of
Chapter~\ref{ch:second_order} reads only the six measured positions, so
this parameterization matters to the formation controller alone.

\section{Forward Kinematics}
\label{sec:appC:forward}

Pair A holds robots 1 and 2, pair B robots 3 and 4, and pair C robots 5
and 6. The pair midpoints are
\begin{equation}
    \mathbf{m}_A = \tfrac{1}{2}(\mathbf{p}_1 + \mathbf{p}_2), \quad
    \mathbf{m}_B = \tfrac{1}{2}(\mathbf{p}_3 + \mathbf{p}_4), \quad
    \mathbf{m}_C = \tfrac{1}{2}(\mathbf{p}_5 + \mathbf{p}_6),
    \label{eq:appC:midpoints}
\end{equation}
and the cluster centroid is their mean,
$\mathbf{p}_c = \tfrac{1}{3}(\mathbf{m}_A + \mathbf{m}_B + \mathbf{m}_C)$.
The midpoint triangle has sides
$p_1 = \lVert\mathbf{m}_B - \mathbf{m}_A\rVert$ and
$q_1 = \lVert\mathbf{m}_C - \mathbf{m}_B\rVert$, with $\beta_1$ the
interior angle at $\mathbf{m}_B$, and the heading $\theta_c$ is the
direction from $\mathbf{m}_A$ to $\mathbf{m}_B$. Each pair has a length
and an orientation,
\begin{equation}
    L_2 = \lVert\mathbf{p}_2 - \mathbf{p}_1\rVert, \qquad
    \theta_2 = \operatorname{atan2}(y_2 - y_1,\, x_2 - x_1),
    \label{eq:appC:pair}
\end{equation}
and likewise $(L_3, \theta_3)$ for pair B and $(L_4, \theta_4)$ for pair C.
The twelve cluster variables are
\begin{equation}
    (x_c,\, y_c,\, \theta_c,\, p_1,\, \beta_1,\, q_1,\,
     L_2,\, \theta_2,\, L_3,\, \theta_3,\, L_4,\, \theta_4).
    \label{eq:appC:state}
\end{equation}

\section{Inverse Kinematics}
\label{sec:appC:inverse}

The inverse kinematics builds the midpoint triangle from
$(p_1, \beta_1, q_1)$ in a local frame, as in
Appendix~\ref{app:kin3}, then rotates it by the heading and translates it
to $(x_c, y_c)$. Each pair's two robots are then placed at its midpoint
plus and minus $\tfrac{1}{2}L_k(\cos\theta_k, \sin\theta_k)$.

\section{Inverse Jacobian}
\label{sec:appC:jacobian}

The $12 \times 12$ inverse Jacobian $\mathbf{J}_c^{-1}$ maps the rates of
the twelve cluster variables to the velocities of the six robots. Its
entries are the partial derivatives of the inverse kinematics and are not
listed here. The simulation software computes them analytically
(\texttt{compute\_inverse\_jacobian} in
\texttt{src/control/pentagon\_kinematics.py}), and the
\texttt{cluster\_builder} tool generates them in closed symbolic form
(\texttt{clusterbuilder.py} with the \texttt{--symbolic} option). The
cluster's heading is not regulated, since the formation controller
commands zero spin throughout.

\chapter{Objectivity and Frame Equivariance of the $s_1$ Tracker}
\label{app:frame_equivariance}

% New text. Draft 11 states the frame shift without the algebra; the
% derivation below was checked 2026-09-29.
Section~\ref{sec:ch3:relation} states that an observer rotating at rate
$\Omega$ sees the determinant field shifted by $\Omega\omega + \Omega^2$
and the strain eigenvalues unchanged. This appendix gives the algebra and
then separates three claims about the $s_1$ tracker that differ in
strength: its features are objective, its command is covariant, and its
trajectory is equivariant only under conditions.

Split the velocity gradient into strain and spin,
\begin{equation}
    \mathbf{S} = \mu\mathbf{I} +
    \begin{bmatrix} s_n & s_s \\ s_s & -s_n \end{bmatrix},
    \qquad
    \mathbf{W} = \frac{\omega}{2}\mathbf{E}, \qquad
    \mathbf{E} = \begin{bmatrix} 0 & -1 \\ 1 & 0 \end{bmatrix},
    \label{eq:appD:split}
\end{equation}
with $r = \sqrt{s_n^2 + s_s^2}$, so the strain eigenvalues are
$s_{1,2} = \mu \mp r$. Expanding the determinant,
\begin{equation}
    D = \det(\mathbf{S} + \mathbf{W})
      = \mu^2 - r^2 + \frac{\omega^2}{4}
      = s_1 s_2 + \frac{\omega^2}{4}.
    \label{eq:appD:D_split}
\end{equation}

The rotating observer sees the field
\begin{equation}
    \mathbf{v}'(\mathbf{p}') = \mathbf{Q}\,\mathbf{v}(\mathbf{Q}^{\top}\mathbf{p}')
        + \Omega\,(-y',\, x')^{\top}.
    \label{eq:appD:rotating_field}
\end{equation}
The gradient of the first term is $\mathbf{Q}\mathbf{J}\mathbf{Q}^{\top}$.
A planar rotation commutes with $\mathbf{E}$, so the strain co-rotates to
$\mathbf{Q}\mathbf{S}\mathbf{Q}^{\top}$, which keeps $\mu$ and $r$, and
the spin is unchanged. The gradient of the second term is
$\Omega\mathbf{E}$, which is pure spin. The observer's gradient is
therefore
$\mathbf{J}' = \mathbf{Q}\mathbf{S}\mathbf{Q}^{\top} +
(\tfrac{\omega}{2} + \Omega)\mathbf{E}$, and by (\ref{eq:appD:D_split})
\begin{equation}
    D' = \mu^2 - r^2 + \Bigl(\frac{\omega}{2} + \Omega\Bigr)^2
       = D + \Omega\,\omega + \Omega^2 ,
    \label{eq:appD:shift}
\end{equation}
while $s_1' = s_1$ and $s_2' = s_2$.

The $s_1$ tracker's channels are built from $\hat{s}_1$, its gradient,
and the strain eigenframe, which all co-rotate with the observer. Three
statements follow, in increasing strength, and they are the three levels
of Section~\ref{sec:ch7:s1_controller}.

First, the features are objective. The strain tensor co-rotates as
$\mathbf{S}' = \mathbf{Q}\mathbf{S}\mathbf{Q}^{\top}$ and its eigenvalues
are unchanged, so $\hat{s}_1$ is the same scalar field in every frame,
while its gradient and the eigenvectors turn with the axes. The capture
test is a ratio of quantities that a rotation leaves invariant, so it gives
the same answer. The determinant is not objective, by
(\ref{eq:appD:shift}).

Second, the command is rotationally covariant where its sign test is
decisive. Given co-rotated inputs and a co-rotated tangent reference, the
channels return the co-rotated command $\mathbf{u}' = \mathbf{Q}\mathbf{u}$.
The tangent reference is the previous output, which a rotating observer
sees turned by $\Omega\,\Delta t$ per step, so the sign rule returns the
same tangent in both frames except where
$|\mathbf{t}_{\mathrm{ref}}^{\top}\mathbf{t}_{\mathrm{raw}}| <
\sin\Omega\,\Delta t$, as at an eigenframe swap. The tangent seed on the
first ride step is the one input that is not objective, since it is taken
from the measured flow, which carries the transport velocity of the frame.

Third, the trajectory is not equivariant in general. The co-rotated
command $\mathbf{u}' = \mathbf{Q}\mathbf{u}$ is not the transformation law
of coordinate velocity in a rotating frame,
$\dot{\mathbf{p}}' = \mathbf{Q}\dot{\mathbf{p}} + \dot{\mathbf{Q}}\mathbf{p}$,
whose second term is the transport velocity of the frame. A cluster that
realizes the co-rotated command still follows the inertial path only if
its across-trench feedback absorbs that term, and the feedback can deliver
no more than the saturated command allows. Section~\ref{sec:ch7:rotating}
tests this in simulation and finds matching paths while the transport
velocity stays inside the delivered speed. This appendix does not prove it.

\chapter{Transverse Stability of the Trackers}
\label{app:stability}

% Ported from: Draft_11_6ec7e51.tex
Both arguments assume the cluster realizes $\dot{\mathbf{p}}_c$
\cite{53,3} and neglect robot lag. Near a trench, with $n$ the distance
across it and $\ell$ the arclength along it, the surrogate is
$f = h(\ell) + \tfrac{1}{2}\kappa_\perp(\ell)\,n^2$ to second order, with
$\kappa_\perp > 0$.

On a quadratic model with Hessian eigenpair $(\lambda_j, \mathbf{w}_j)$,
the gradient component $g_j = \mathbf{w}_j^\top\nabla f$ obeys
$\dot{g}_j = \lambda_j\mathbf{w}_j^\top\dot{\mathbf{p}}_c$, so the
saturated Newton channel
$\mathbf{w}_j^\top\dot{\mathbf{p}}_c = -k\,\mathrm{sat}(x_j)$,
$x_j = g_j/\lambda_j$, gives
\begin{equation}
    \frac{d}{dt}\,\tfrac{1}{2}g_j^2
    = -k\,\lambda_j^2\,x_j\,\mathrm{sat}(x_j) \leq 0
    \label{eq:appE:newton_lyap}
\end{equation}
for either sign of $\lambda_j$, with equality only at $g_j = 0$. A full
Newton step would therefore halt at the crest, so the $D$ tracker keeps
only the across-trench channel.

\begin{proposition}
Let $\dot{n} = -k\,\mathrm{sat}(a_\perp n + e)$ with $a_\perp > 0$ and
$|e| \leq \delta$. Then $\limsup_{t\to\infty}|n| \leq \delta/a_\perp$.
If $e = 0$, $|n|$ falls below $c_{\max}/a_\perp$ in finite time, then
decays at rate at least $k\,a_\perp\tanh 1$.
\end{proposition}

\begin{proof}
With $V = \tfrac{1}{2}n^2$,
$\dot{V} = -k\,n\,\mathrm{sat}(a_\perp n + e) < 0$ for
$|n| > \delta/a_\perp$, where the argument of $\mathrm{sat}$ has the
sign of $n$. With $e = 0$, $|\dot{n}| \geq k\,c_{\max}\tanh 1$
while $a_\perp|n| \geq c_{\max}$, and $|\tanh x| \geq |x|\tanh 1$ for
$|x| \leq 1$ gives the rate.
\end{proof}

Let the fitted transverse gradient be $\kappa_\perp n + \epsilon$,
$|\epsilon| \leq \delta_g$. In (\ref{eq:ch7:d_law}),
$a_\perp = \kappa_\perp/\hat{\kappa}_\perp$ and
$e = \epsilon/\hat{\kappa}_\perp$. In (\ref{eq:ch7:s1_law}),
$a_\perp = g_\perp\kappa_\perp$ and $e = g_\perp\epsilon$. Both bounds
equal $\delta_g/\kappa_\perp$, independent of gain and fitted
curvature, but only the $D$ tracker needs $\hat{\kappa}_\perp > 0$.
Where $\kappa_\perp$ vanishes over a length $\Delta\ell$, as at the
kink of $s_1$, the $s_1$ ride crosses it in time $\Delta\ell/(k\,c_{\max}\tanh 1)$
and $|n|$ grows by at most $\Delta\ell/\tanh 1$. A frame rotating at
$\Omega$ perturbs its $\dot{n}$ by at most $\Omega\lVert\mathbf{p}\rVert$,
leaving $n$ bounded while $\Omega\lVert\mathbf{p}\rVert < k\,c_{\max}$.

Both branches of (\ref{eq:ch7:vpar}) are nonnegative, so the $D$ tracker
never reverses along the trench. If $\mathrm{sat}(v_\parallel) \geq v_{\min} > 0$ outside
the flow-saddle neighborhoods where (\ref{eq:ch7:d_capture}) applies, it crosses a
segment of length $L_\Gamma$ within $L_\Gamma/(k\,v_{\min})$.

\chapter{Statistical Methods}
\label{app:statistics}

% New text
\section{Bias and Precision}
Chapter~\ref{ch:first_order} summarizes each set of convergence trials by
its final centroid positions. The systematic bias is the distance from
the mean final position to the true critical point, and it measures what
averaging cannot remove. The precision is the standard deviation of the
final positions, and it measures repeatability. Orbital trials are
summarized by the mean and standard deviation of the radial error.

\section{Circular Statistics for Hue}
Hue is an angle, so a reading near 0 and a reading near 1 are the same
direction. The hue network of Section~\ref{sec:ch2:calibration} predicts
the sine and cosine of hue for this reason, and its accuracy is reported
with circular statistics.

\section{Success Rates}
The noise sweeps of Chapter~\ref{ch:second_order} report success rates
over $N = 10^4$ trials per noise level. Each trial succeeds or fails
independently, so the standard error of a rate $\hat{p}$ is
$\sqrt{\hat{p}(1-\hat{p})/N} \leq 1/(2\sqrt{N}) = 0.005$.


\end{document}
